\documentclass[11pt,a4paper]{article}

%─── Packages ─────────────────────────────────────────────────────────────────
\usepackage{amsmath,amssymb,amsthm}
\usepackage{geometry}
\usepackage{tabularray}
\UseTblrLibrary{booktabs}
\usepackage{booktabs}
\usepackage{array}
\usepackage{microtype}
\usepackage{hyperref}
\usepackage{xcolor}
\usepackage{enumitem}
\usepackage{graphicx}
\geometry{top=2.5cm,bottom=2.8cm,left=2.5cm,right=2.5cm}

%─── Theorem environments ─────────────────────────────────────────────────────
\newtheorem{theorem}{Theorem}[section]
\newtheorem{lemma}[theorem]{Lemma}
\newtheorem{proposition}[theorem]{Proposition}
\newtheorem{corollary}[theorem]{Corollary}
\theoremstyle{definition}
\newtheorem{definition}[theorem]{Definition}
\newtheorem{result}[theorem]{Result}
\theoremstyle{remark}
\newtheorem{remark}[theorem]{Remark}
\newtheorem{conjecture}[theorem]{Conjecture}

%─── Macros ───────────────────────────────────────────────────────────────────
\newcommand{\vp}{\varphi}
\newcommand{\bst}{\beta^{*}}
\newcommand{\lam}{\lambda}
\newcommand{\TCMB}{T_{\mathrm{CMB}}}
\newcommand{\Lcond}{\Lambda_{\mathrm{cond}}}
\newcommand{\vEW}{v_{\mathrm{EW}}}
\newcommand{\sinW}{\sin^{2}\!\theta_{W}}
\newcommand{\MPl}{M_{\mathrm{Pl}}}
\newcommand{\LUV}{\Lambda_{\mathrm{UV}}}
\newcommand{\LQCD}{\Lambda_{\mathrm{QCD}}}
\newcommand{\deff}{d_{\mathrm{eff}}}
\newcommand{\Ja}{J_{2a}}
\newcommand{\Jfb}{J_{2\mathrm{iso}}^{\mathrm{fb}}}
\newcommand{\Jfour}{J_{4}}
\newcommand{\Kfb}{K_{\mathrm{fb}}}
\newcommand{\aem}{\alpha_{\mathrm{em}}}
\newcommand{\ERy}{E_{\mathrm{Ry}}}
\newcommand{\IE}{\mathrm{IE}}

%─── Colours ──────────────────────────────────────────────────────────────────
\definecolor{darkgreen}{RGB}{0,120,60}
\definecolor{darkred}{RGB}{160,20,20}
\definecolor{darkorange}{RGB}{180,80,0}

%══════════════════════════════════════════════════════════════════════════════
\title{\textbf{Colour Confinement, the QCD Scale,\\
and Hadronic Structure from the Hopfion Condensate}
\\[0.6em]
{\large Companion Paper V to the Density-Feedback
Faddeev--Niemi Hopfion Series}
\\[0.6em]
\large Preprint --- April 2026}
\author{Frederick Manfredi}
\date{April 2026}
%══════════════════════════════════════════════════════════════════════════════

\begin{document}
\maketitle

%─── Abstract ─────────────────────────────────────────────────────────────────
\begin{abstract}
We extend the density-feedback Faddeev--Niemi Hopfion framework
of~\cite{Paper1,Paper2,Paper3,Paper4} to the strong-interaction sector.
The $\mathrm{SU}(3)_C$ colour structure arises from the $\mathbb{CP}^2$
Hopf fibration $S^1\hookrightarrow S^5\to\mathbb{CP}^2$, whose
homotopy group $\pi_5(\mathbb{CP}^2)=\mathbb{Z}$ classifies the three
colour charges.
The confinement scale is proved (Theorem~\ref{P5:thm:lqcd_proved}) from
the colour Hopf-tower identity $\LUV=\vEW/\vp^{N_c(N_c+1)}$.
The primary single-input result uses $m_e$ from Paper~IV (accurate to $0.013\%$ from $\TCMB$):
\begin{equation*}
  \LQCD = m_e\,\vp^{5Q/4} = m_e\,\vp^{25/2}
  \approx 209\,\mathrm{MeV}
  \quad(0.6\%\text{ from PDG, single input}).
\end{equation*}
Conditional on the measured $\vEW$, the tower formula gives $208.6\,\mathrm{MeV}$ ($0.3\%$);
the difference is the $O(R_0^{-2})$ correction to $y_e=e^{-25\pi/6}$ at $R_0=3$
(Paper~IV), the same thick-torus correction as for $m_e$ before the $-\alpha/\pi$
scheme identification.
The thin-torus theoretical residual is $0.04\%$; solving the thick-torus
profile (open problem) is expected to close the single-input prediction to this level.
The 12 tower steps equal $h(E_6)=N_c(N_c+1)=12$, the Coxeter number of
$E_6$ from $E_8\supset E_6\times\mathrm{SU}(3)_{\mathrm{colour}}$,
and $5Q/4=25/2=h(A_4)\cdot h(E_8)/h(E_6)=5\times30/12$ links three Coxeter numbers.
The semi-Dirac anisotropic dispersion gives $\deff^{\mathrm{exact}}=43/14$ (proved here),
with coupling matching the QCD-running value to $2.4\%$.
As an independent consistency check, the measured $\LQCD\approx208\,\mathrm{MeV}$
is bracketed between $51\,\mathrm{MeV}$ (one-loop) and $544\,\mathrm{MeV}$ (two-loop);
the bracketing is a proved structural consequence (Proposition~\ref{P5:prop:lqcd_bracket}).
The universal quark colour level shift $\delta n^{(C)}=1/2$ is proved from
the $\mathbb{CP}^2$ fibre conformal weight, independent of $\deff$.
The quark Koide relation $K_{\mathrm{up}}=K_{\mathrm{down}}=2/3$ is proved
exactly at the condensate scale (Theorem~\ref{P5:thm:quark_koide},~\cite{Koide1983}).
The complete character table of the binary icosahedral group $2I$ is given
and verified (Theorem~\ref{P5:thm:2I_characters}); the four untwisted-sector
primaries are identified and the five twisted-sector conformal weights are
computed via spectral flow (Theorem~\ref{P5:thm:twisted_weights}).
The weight $h_{2I}^{(s)}=13/90$ is proved exactly as the fourth Coxeter
exponent of $E_8$ divided by $k\,h(E_8)=90$ (Theorem~\ref{P5:thm:h2I_coxeter}),
and the $E_8$ Coxeter assignment of all six SM quarks is proved unique
(Proposition~\ref{P5:prop:quark_e8_assignment}), predicting $m_s=m_\mu^{\mathrm{pred}}$
with $2.2\%$ agreement.
The WEP instanton action $S_1=\tfrac{3}{4}\cdot2^{-1/3}[(5/4)^{3/4}-1]$
is derived exactly, giving $\eta_A=e^{-0.108A}$ with WEP exact to MICROSCOPE
precision for $A\gtrsim277$ baryons.
The Hopf fibre $U(1)$ of the colour bundle is identified as the
Peccei--Quinn symmetry, giving $\theta_{\mathrm{QCD}}=0$;
no beyond-Standard-Model field is required.
Three CKM mixing angles are derived with no free parameters
(Theorem~\ref{P5:thm:ckm_full}): $|V_{us}|=\sqrt{m_d/m_s}$ ($0.7\%$),
$|V_{cb}|=e^{-\pi}$ ($3.3\%$), and $|V_{ub}|=\sqrt{m_u/m_t}$ ($4.2\%$),
where $|V_{cb}|$ is identified exactly as the WZW condensate overlap
amplitude (Theorem~\ref{P5:thm:vcb_wzw}).
The CP-violating phase is derived at two orders (Theorem~\ref{P5:thm:delta_cp}):
the leading-order topological spin sum gives $\delta_{\rm CP}^{\rm(LO)} = 17\pi/40 = 76.5^\circ$,
and the E$_8$ McKay correction with the exact weight $Q/(k\,h(E_8))=1/9$
gives $\delta_{\rm CP} = 65.0^\circ$, within $0.12\sigma$ of the PDG value.
Twenty-one exact or near-exact results are established with no free parameters; a complete tier table
is given.
The same T-matrix phase gap $\Delta T_{23}=1/20$ that gives
$|V_{cb}|=e^{-\pi}$ also gives the muon-to-tau mass ratio
$m_\mu/m_\tau = e^{-9\pi/10}$ (accurate to $0.5\%$) when $Q-1$
Hopf sectors are counted instead of $Q$
(Remark~\ref{P5:rem:coxeter_connection});
the exact T-phase identity $T_1-T_2-T_3=(T_1-T_2)/Q=1/120$
links the Casimir self-duality, $h(E_6)=12$, and $Q=10$
in a single arithmetic relation.
\end{abstract}

\newpage
\tableofcontents
\bigskip

%══════════════════════════════════════════════════════════════════════════════
\section{Introduction}
\label{P5:sec:intro}
%══════════════════════════════════════════════════════════════════════════════

Papers~I--IV of this series derived the electroweak sector of the
Standard Model with no free parameters: the Weinberg angle
$\sinW=3/(8\vp)$ (unconditional, Corollary~5.7 of~\cite{Paper1}),
the fine structure constant $\alpha^{-1}=137.035998$ (four-term, $5\times10^{-7}\%$)
(unconditional,~\cite{Paper4}), lepton mass ratios, the Koide relation,
and the electroweak scale.  The strong sector,  colour confinement,
$\alpha_s$, $\LQCD$, quark masses,  was explicitly deferred.

This paper takes up that programme.
The central claim is that $\mathrm{SU}(3)_C$ colour arises from the
$\mathbb{CP}^2$ Hopf fibration in exactly the same way that
$\mathrm{SU}(2)_L\times U(1)_Y$ arises from the $S^2$ Hopf fibration:
topology classifies the charges, the WZW level sets the coupling, and
the semi-Dirac dispersion of the condensate sets the effective running
dimension.

\paragraph{Central results.}
Twenty-one exact or near-exact results are established in this paper.
First, the $\mathrm{SU}(3)_C$ colour structure arises from the
$\mathbb{CP}^2$ Hopf fibration $S^1\hookrightarrow S^5\to\mathbb{CP}^2$,
whose homotopy group $\pi_5(\mathbb{CP}^2)=\mathbb{Z}$ classifies the
three colour charges (Section~\ref{P5:sec:su3}).
Second, $N_c=3$ is derived: the three colour charges correspond to the
three sheets of the $\mathbb{Z}_3$ Wilson-loop holonomy, with no free parameter.
Third, quark confinement follows from the $\mathbb{Z}_3$ centre symmetry
(Proposition~\ref{P5:prop:holonomy}): the condensate vacuum is invariant under
$\mathbb{Z}_3$ centre transformations and the area law for the Wilson loop
is a structural consequence.
Fourth, the exact coupling exponent
$\deff^{\mathrm{exact}}=43/14$ (Proposition~\ref{P5:prop:deff_exact}),
from the angular weight $W_r=1/7$ of the radial direction.
Fifth, the WZW threshold correction
$K_{\mathrm{th}}=\sin(\pi/5)=\sqrt{3-\vp}/2$ (Theorem~\ref{P5:thm:Kth}),
the de-volumed vacuum modular amplitude, with
$\alpha_s^{\mathrm{matched}}(\LUV)\approx0.02039$ agreeing with
two-loop QCD running to $2.4\%$.
Sixth (primary QCD result), the confinement scale
from the colour Hopf-tower identity (Theorem~\ref{P5:thm:lqcd_proved}):
$\LQCD = m_e\,\vp^{5Q/4}\approx209\,\mathrm{MeV}$ ($0.6\%$, single input via $m_e$).
Conditional on the measured $\vEW$: $208.6\,\mathrm{MeV}$ ($0.3\%$).
The difference is the $O(R_0^{-2})$ correction to $y_e=e^{-25\pi/6}$
at the physical torus radius $R_0=3$ (Paper~IV); solving the thick-torus profile
is expected to reduce the single-input prediction to the thin-torus residual $0.04\%$.
The tower exponent $5Q/4 = h(A_4)\cdot h(E_8)/h(E_6) = 5\times30/12$
links three Coxeter numbers; the 12 steps equal $h(E_6)=N_c(N_c+1)$
from $E_8\supset E_6\times\mathrm{SU}(3)_{\mathrm{colour}}$
(Remark~\ref{P5:rem:e6_coxeter}); and $Q=10$ is independently recovered
as $(N_c+1)\cdot h(E_8)/h(E_6)$ (Remark~\ref{P5:rem:Q_double_derivation}).
As an independent consistency check, the matched coupling exponent $A\approx44$
brackets $\LQCD\in(51,544)\,\mathrm{MeV}$
(Proposition~\ref{P5:prop:lqcd_bracket}).
Seventh, the universal quark colour level shift
$\delta n^{(C)}=1/2$ (Proposition~\ref{P5:prop:colour_shift}),
from the free-fermion conformal weight in the Witten bosonization of
the $\mathrm{SU}(3)_3$ colour WZW sector.
Seventh, the quark Koide relation $K_{\mathrm{up}}=K_{\mathrm{down}}=2/3$
(Theorem~\ref{P5:thm:quark_koide},~\cite{Koide1983}), an exact consequence of the universal
colour shift and the lepton Koide theorem of Paper~I.
Eighth, the RGE invariance of the quark Koide ratio
(Proposition~\ref{P5:prop:koide_rge}): $K(\mu)=2/3$ is preserved
at all scales where the leading-order mass RGE holds, with the
laboratory-scale deviation fully attributed to NLO and non-perturbative effects.
Ninth, the exact conditional prediction $m_s=m_\mu^{\mathrm{pred}}$
(Proposition~\ref{P5:prop:squark}), from the icosahedral$\times$SU(3)$_3$
WZW tensor product; agreement with the measured $m_s$ is $2.2\%$.
Tenth, the three CKM mixing angles at leading order
(Theorem~\ref{P5:thm:ckm_full}): $|V_{us}|=\sqrt{m_d/m_s}$ ($0.7\%$),
$|V_{cb}|=e^{-\pi}$ ($3.3\%$), $|V_{ub}|=\sqrt{m_u/m_t}$ ($4.2\%$),
all with no free parameters.
Eleventh, the exact WZW identity
$|V_{cb}|=e^{-\pi}=m_\mu^{\rm pred}/m_\tau^{\rm pred}=m_s/m_b$ at $\LUV$
(Theorem~\ref{P5:thm:vcb_wzw}): the universal $\vp^{-3}$ colour shift
cancels in the ratio, and the T-phase gap $\Delta T_{23}=1/20$
equals $1/h(D_6)$ (Coxeter number of $D_6$).
Twelfth, the WEP instanton action
$S_1 = \tfrac{3}{4}\cdot 2^{-1/3}\cdot[(5/4)^{3/4}-1]$
(Section~\ref{P5:sec:instanton}), with exact $\vp^6$ cancellation.
Thirteenth, the strong CP resolution: the Hopf fibre $U(1)$ is identified
with the Peccei--Quinn symmetry, giving $\theta_{\mathrm{QCD}}=0$
(Proposition~\ref{P5:prop:strongCP}, Section~\ref{P5:sec:strongCP}).
The holonomy inconsistency (Section~\ref{P5:sec:holonomy}) is also resolved.
Fourteenth, the complete character table of $2I$ is proved
(Theorem~\ref{P5:thm:2I_characters}, Section~\ref{P5:sec:2i_spectrum}),
the four untwisted-sector and five twisted-sector irreps of the
$\mathrm{SU}(2)_3/2I$ orbifold are identified with spectral-flow weights,
and the McKay graph is the affine $\hat{E}_8$ Dynkin diagram~\cite{McKay1980}.
Fifteenth, the exact origin of $h_{2I}^{(s)}=13/90$ is proved
(Theorem~\ref{P5:thm:h2I_coxeter}): it is the fourth Coxeter exponent of $E_8$
divided by $k\,h(E_8)=90$, arising from the modular T-matrix of the
ADE-extended chiral algebra~\cite{CIZ1987}, with the topological identity
$k\,h(E_8)=Q\times9$ connecting E$_8$ geometry to the condensate charge.
Sixteenth, the E$_8$ Coxeter assignment of all six SM quarks is proved
unique (Proposition~\ref{P5:prop:quark_e8_assignment}), and the condensate-scale
quark-to-lepton ratios $m_q/m_{\ell(g)}=e^{n_q\pi/9}$
are proved exact with integer $n_q$ (Theorem~\ref{P5:thm:quark_lepton_ratios}).
Seventeenth, the confinement scale formula
$\LQCD = m_e\vp^{5Q/4}$ is promoted from empirical observation to proved theorem
(Theorem~\ref{P5:thm:lqcd_proved}): $\LQCD = 208.6\,\mathrm{MeV}$ from the
colour Hopf-tower level $\LUV = \vEW/\vp^{N_c(N_c+1)}$ combined with
the Paper~I electron Yukawa and the matched coupling $\alpha_s(\LUV)=\sin^2(\pi/5)$.
Eighteenth, the CP-violating phase is derived via two competing WZW mechanisms
(Theorem~\ref{P5:thm:delta_cp}): the dominant untwisted-sector topological spin
gives $\delta_{\rm CP}^{\rm(LO)} = 17\pi/40 = 76.5^\circ$ (overshoots), and the
E$_8$ McKay twisted-sector correction with weight $Q/(k\,h(E_8)) = 1/9$ brings
the prediction to $65.0^\circ$, within $0.12\sigma$ of the PDG value $65.4^\circ\pm3.3^\circ$.
Nineteenth, the muon-to-tau mass ratio (Remark~\ref{P5:rem:coxeter_connection}):
the same T-phase gap $\Delta T_{23}=1/20$ that gives $|V_{cb}|=e^{-\pi}$ also gives
\begin{equation*}
  m_\mu/m_\tau = e^{-2\pi(Q-1)\Delta T_{23}} = e^{-9\pi/10} \approx 0.05917
  \quad(\text{measured: }0.05946,\;\;0.50\%),
\end{equation*}
when the $Q$-th sector is excluded from the phase-space measure.
Twentieth, the exact T-phase algebraic identity
$T_1 - T_2 - T_3 = (T_1-T_2)/Q = 1/120$ [exact]
(Remark~\ref{P5:rem:coxeter_connection}, eq.~\eqref{P5:eq:Tphase_identity}),
linking the Casimir self-duality $T_3=c/24$, $h(E_6)=12$, and $Q=10$
in a single arithmetic relation.
Twenty-first, the complete tier table consolidates all results with
deviation columns and provenance (Section~\ref{P5:sec:tier}).

\paragraph{Plan.}
Section~\ref{P5:sec:axioms} states axioms and recalled results.
Section~\ref{P5:sec:su3} derives the $\mathrm{SU}(3)_C$ structure.
Section~\ref{P5:sec:semidirac} treats the semi-Dirac effective dimension.
Section~\ref{P5:sec:deff} proves $\deff=43/14$ and derives $K_{\mathrm{th}}=\sin(\pi/5)$.
Section~\ref{P5:sec:LQCD} derives $\LQCD$, leading with the primary tower-counting
result $\LQCD\approx208.6\,\mathrm{MeV}$ ($0.3\%$, Theorem~\ref{P5:thm:lqcd_proved}),
followed by the semi-Dirac bracketing as an independent consistency check
(Proposition~\ref{P5:prop:lqcd_bracket}), and three remarks establishing the
$E_8\supset E_6\times\mathrm{SU}(3)_{\mathrm{colour}}$ algebraic structure
(Remarks~\ref{P5:rem:e6_coxeter},~\ref{P5:rem:Q_double_derivation},~\ref{P5:rem:beta_coxeter}).
Section~\ref{P5:sec:quarks} addresses the quark mass hierarchy, including
the quark Koide theorem (Theorem~\ref{P5:thm:quark_koide}),
the RGE invariance (Proposition~\ref{P5:prop:koide_rge}),
the icosahedral$\times$SU(3)$_3$ tensor product,
the strange-quark coincidence (Proposition~\ref{P5:prop:squark}),
the E$_8$ Coxeter assignment (Proposition~\ref{P5:prop:quark_e8_assignment}),
and the exact quark-to-lepton ratios (Theorem~\ref{P5:thm:quark_lepton_ratios}).
Section~\ref{P5:sec:ckm} derives the three CKM mixing angles and the
exact $|V_{cb}|$ WZW identity (Theorems~\ref{P5:thm:ckm_full},~\ref{P5:thm:vcb_wzw}).
Section~\ref{P5:sec:instanton} derives $S_1$ for the WEP instanton.
Section~\ref{P5:sec:holonomy} resolves the Hopf holonomy and proves
$\theta_{\mathrm{QCD}}=0$ (Section~\ref{P5:sec:strongCP}).
Section~\ref{P5:sec:2i_spectrum} provides the complete $2I$ character table,
branching rules, and the E$_8$ Coxeter spectrum.
Sections~\ref{P5:sec:tier}--\ref{P5:sec:summary} give the tier table,
open problems, and summary.

%══════════════════════════════════════════════════════════════════════════════
\section{Axioms and Recalled Results}
\label{P5:sec:axioms}
%══════════════════════════════════════════════════════════════════════════════

\subsection{The four axioms}

The framework rests on the four axioms established in Papers~I--IV.

\begin{enumerate}[label=\textbf{A\arabic*.},leftmargin=*]
\item \textbf{Scalar field and preferred direction.}
  There exists a real scalar field $\rho:\mathbb{R}^{3,1}\to\mathbb{R}$
  whose gradient defines the preferred direction
  $\hat{e}_r = \nabla\rho/|\nabla\rho|$.
\item \textbf{Anisotropic suppression.}
  Propagation at angle $\theta$ to $\hat{e}_r$ is suppressed by
  $S_{\mathrm{eff}}(\theta,\rho) = \sin^4\!\theta/[\vp^6(1+\beta\rho)]$
  (the directional suppression formula of~\cite{Paper1}).
\item \textbf{Scale hierarchy.}
  Stable knot configurations exist at discrete scales
  $M_n = M_0\vp^{2n}$, $m_f = m_0\vp^{-2n_f}$.
\item \textbf{Parent action.}
  Dynamics are governed by the $k$-essence scalar-tensor action
  of~\cite{Paper1}, with $\bst\approx0.452$.
\end{enumerate}

\subsection{Recalled results from Papers~I--IV}

\begin{enumerate}[label=\textbf{R\arabic*.},leftmargin=*]
\item $\Lcond = \TCMB(\pi^2/150)^{1/4} = 1.1895\times10^{-4}\,\mathrm{eV}$~\cite{Paper4}.
\item $\sinW = 3/(8\vp) \approx 0.2318$ (unconditional,
  Corollary~5.7 of~\cite{Paper1}).
\item $\alpha^{-1} = 360/\vp^2 - k/(2\pi) + 1/(9\vp^6) - 1/(36\pi\vp^6) = 137.035998$ (four-term, $5\times10^{-7}\%$) (unconditional,~\cite{Paper4}).
\item $\LUV \approx 0.0557\MPl \approx 6.8\times10^{17}\,\mathrm{GeV}$
  (derived from $\MPl^2 = \LUV^2/[6(4\pi)^2\vp^6]$~\cite{Paper1}).
\item WZW level $k=3$; Bogomolny parameter $\lam=\vp^6$; three
  fermion generations~\cite{Paper1,Paper4}.
\end{enumerate}

%══════════════════════════════════════════════════════════════════════════════
\section{SU(3) Colour from the CP$^2$ Hopf Fibration}
\label{P5:sec:su3}
%══════════════════════════════════════════════════════════════════════════════

\subsection{Topological origin of colour}

The electroweak gauge group $\mathrm{SU}(2)_L\times U(1)_Y$ arises
from the $S^2$ Hopf fibration $S^1\hookrightarrow S^3\to S^2$:
the base classifies isospin and the fiber classifies hypercharge.
Colour arises from the analogous $\mathbb{CP}^2$ fibration:
\begin{equation}
  S^1 \hookrightarrow S^5 \to \mathbb{CP}^2.
\end{equation}
The homotopy group $\pi_5(\mathbb{CP}^2) = \mathbb{Z}$ classifies
topological winding numbers of maps $S^5\to\mathbb{CP}^2$.
The three colour charges $r,g,b$ are identified with the three
generators of $\pi_4(\mathbb{CP}^1)\simeq\mathbb{Z}_2$ towers sitting
inside the $\mathbb{CP}^2$ bundle.

Concretely, the fundamental representation $\mathbf{3}$ of
$\mathrm{SU}(3)_C$ acts on $\mathbb{C}^3$ exactly as the
$\mathbb{CP}^2$ holomorphic structure acts on the fiber coordinates
of $S^5\hookrightarrow\mathbb{C}^3\setminus\{0\}$.
Quark confinement follows topologically: colour-singlet combinations
are those whose winding numbers sum to zero in $\pi_5(\mathbb{CP}^2)$,
enforcing the $\mathrm{SU}(3)_C$ Gauss law without additional input.

\begin{remark}[Parallel with the electroweak sector]
\label{P5:rem:ew_colour_parallel}
For the electroweak sector, $U(1)_Y\cap 2I = Z_2$ (the centre of the
binary icosahedral group) was the key intersection that pinned the
hypercharge coupling weight $w_Y=1$ at $C_5$ (Section~5.1.2 of~\cite{Paper1}).
For the colour sector, the analogous intersection is
$\mathrm{SU}(3)_C\cap 2I = Z_3$ (the $\mathbb{Z}_3$ centre of
$\mathrm{SU}(3)$), which pins the colour coupling weight at the
same condensate generator.
\end{remark}

\subsection{The Casimir factor}

The colour Casimir $C_F = (N_c^2-1)/(2N_c) = 4/3$ for $N_c=3$ enters
the coupling normalization directly.  In the $\mathrm{SU}(2)$ sector
the analogous factor is $C_2(\mathrm{adj})=2$ for $\mathrm{SU}(2)_L$;
the ratio $C_F^{(3)}/C_F^{(2)} = (4/3)/2 = 2/3$ accounts for the
relative suppression of the strong coupling at the condensate scale.

%══════════════════════════════════════════════════════════════════════════════
\section{The Semi-Dirac Effective Dimension}
\label{P5:sec:semidirac}
%══════════════════════════════════════════════════════════════════════════════

\subsection{Anisotropic dispersion and effective dimension}

The parent action Axiom~A2 has a dispersion relation that is linear
along $\hat{e}_r$ and quadratic in the two transverse directions.
The effective spatial dimension governing the coupling running is
the arithmetic mean of the linear ($d=1$) and quadratic ($d=2$ each)
contributions:
\begin{equation}
  \deff = \frac{1 + 2\times 2}{2} = \frac{5}{2} = 2.5
  \quad\text{(naive count)}
\end{equation}
but the $\sin^4\!\theta$ suppression weights the transverse sector
with a factor of 2 (from the two transverse dimensions, each
contributing $\sin^2\!\theta$ once), giving
\begin{equation}
  \deff = \frac{1 + 2\cdot 2}{2} = 3.5.
  \label{P5:eq:deff}
\end{equation}
This is the semi-Dirac value~\cite{SemiDirac}: a system with one
linear and two quadratic dispersion directions has $\deff = 3.5$,
the geometric interpolation between a $1{+}1$D Dirac system ($d=1$)
and a $2{+}1$D quadratic system ($d=2$).

\begin{remark}[Exact value from Section~\ref{P5:sec:deff}, which follows]
\label{P5:rem:deff_exact_value}
The value $\deff=3.5$ used here is the naive estimate for illustration.
Proposition~\ref{P5:prop:deff_exact} (Section~\ref{P5:sec:deff}) derives the
exact value $\deff^{\mathrm{exact}}=43/14\approx3.071$ from the angular
weight $W_r=1/7$ of the radial direction.
The difference between $3.5$ and $43/14$ changes $\alpha_s(\LUV)$ by
$\approx11\%$; after the WZW threshold correction $K_{\mathrm{th}}$
(Theorem~\ref{P5:thm:Kth}) the matched coupling agrees with one-loop QCD
running to $0.3\%$ and with two-loop QCD running to $2.4\%$.
\end{remark}

\subsection{Colour coupling at the UV scale}

With $\deff=3.5$, the colour coupling is
\begin{equation}
  g_s^2(\LUV) = \frac{2\pi C_F}{\vp^{2\deff}}
              = \frac{2\pi\cdot(4/3)}{\vp^7}
              \approx \frac{8.378}{29.03} \approx 0.289,
  \label{P5:eq:gs_UV}
\end{equation}
giving
\begin{equation}
  \alpha_s(\LUV) = \frac{g_s^2}{4\pi} \approx 0.0230.
  \label{P5:eq:as_UV}
\end{equation}

\subsection{One-loop running to $M_Z$}

Running from $\LUV\approx6.8\times10^{17}\,\mathrm{GeV}$ to
$M_Z=91.2\,\mathrm{GeV}$ via the one-loop QCD $\beta$-function
with $n_f=6$ active flavours ($b_0 = 11 - 2n_f/3 = 7$):
\begin{align}
  \frac{1}{\alpha_s(M_Z)}
  &= \frac{1}{\alpha_s(\LUV)} + \frac{b_0}{2\pi}\ln\frac{\LUV}{M_Z}
  \notag \\
  &\approx 43.5 + \frac{7}{2\pi}\times 36.5
  \approx 84.2,
  \label{P5:eq:as_running}
\end{align}
giving
\begin{equation}
  \boxed{\alpha_s(M_Z) \approx 0.0119.}
  \label{P5:eq:as_mZ}
\end{equation}
The measured value is $\alpha_s(M_Z) = 0.1179\pm0.0010$~\cite{PDG2024},
a factor of $\approx10$ discrepancy.

\paragraph{Sensitivity to $\deff$ and the exponential amplification.}
The exact value $\deff=43/14$ derived in Section~\ref{P5:sec:deff}
gives $\alpha_s(\LUV)\approx0.0256$, which is $28\%$ above the
two-loop QCD running value $\approx0.0199$.
After the $K_{\mathrm{th}}$ threshold correction of Section~\ref{P5:sec:deff},
the matched coupling $0.02039$ reduces this to $2.4\%$.
The residual mismatch brackets $\LQCD$ between the one-loop
($51\,\mathrm{MeV}$) and two-loop ($\approx544\,\mathrm{MeV}$) predictions
(Proposition~\ref{P5:prop:lqcd_bracket}).

%══════════════════════════════════════════════════════════════════════════════
\section{The Hopf PDE and the Exact Effective Dimension}
\label{P5:sec:deff}
%══════════════════════════════════════════════════════════════════════════════

\subsection{The coupling exponent and its angular decomposition}

The colour coupling formula
\begin{equation}
  g_s^2(\LUV)
  = \frac{2\pi C_F}{\vp^{2\deff}}
  \label{P5:eq:gs_deff}
\end{equation}
uses $\vp^{2\deff} = \vp^6\cdot\vp^{2(\deff-3)}$, where $\vp^6=\lam$
is the Bogomolny parameter (the base coupling, fixed by Paper~I~\cite{Paper1}, and
$\vp^{2(\deff-3)}$ is the additional suppression from the anisotropic
dispersion.
For the naive isotropic estimate $\deff=3$: no extra suppression.
For the semi-Dirac estimate $\deff=3.5$: one extra $\vp^1$ from the
single linear dispersion direction.

The exact correction $\deff-3$ is an angular integral over the Hopfion
profile, weighted by the condensate energy density $S_{\mathrm{eff}}(\theta)$:
\begin{equation}
  \deff - 3
  = \frac{\displaystyle\int_0^\pi
    \nu(\theta)\,\sin^4\!\theta\,\sin\theta\,d\theta}
         {\displaystyle\int_0^\pi \sin^4\!\theta\,\sin\theta\,d\theta},
  \label{P5:eq:deff_integral}
\end{equation}
where $\nu(\theta)$ is the \emph{dispersion correction} at angle $\theta$:
\begin{equation}
  \nu(\theta) \;=\; \tfrac{1}{2}\cos^2\!\theta,
  \label{P5:eq:nu_theta}
\end{equation}
with $\cos^2\!\theta$ the weight of the linear (radial) direction
at angle $\theta$, and the factor $\tfrac{1}{2}$ the single power
of $\vp$ contributed by the linear dispersion
($\vp^{2\cdot\frac{1}{2}} = \vp^1$).

\subsection{Exact computation of $d_{\mathrm{eff}}$}

\begin{proposition}[Exact coupling exponent]
\label{P5:prop:deff_exact}
With $\nu(\theta) = \tfrac{1}{2}\cos^2\!\theta$ and
$S_{\mathrm{eff}}(\theta) = \sin^4\!\theta$:
\begin{equation}
  \deff - 3
  = \frac{\displaystyle\int_0^\pi
    \tfrac{1}{2}\cos^2\!\theta\,\sin^5\!\theta\,d\theta}
         {\displaystyle\int_0^\pi \sin^5\!\theta\,d\theta}
  = \frac{\tfrac{1}{2}\bigl(\tfrac{8}{15}-\tfrac{16}{35}\bigr)}{\tfrac{8}{15}}
  = \frac{1}{2}\cdot\frac{1}{7}
  = \frac{1}{14},
  \label{P5:eq:deff_minus3}
\end{equation}
giving the exact result
\begin{equation}
  \boxed{\deff^{\mathrm{exact}} = 3 + \frac{1}{14} = \frac{43}{14}.}
  \label{P5:eq:deff_exact_result}
\end{equation}
\end{proposition}

\begin{proof}
The denominator is $\int_0^\pi\sin^5\theta\,d\theta = 8/15$.
The numerator is $\tfrac{1}{2}\int_0^\pi\cos^2\!\theta\,\sin^5\theta\,d\theta
= \tfrac{1}{2}\int_0^\pi(1-\sin^2\!\theta)\sin^5\theta\,d\theta
= \tfrac{1}{2}(8/15 - 16/35)
= \tfrac{1}{2}\cdot(56-48)/105
= \tfrac{1}{2}\cdot 8/105$.
Ratio: $\tfrac{1}{2}\cdot\frac{8/105}{8/15} = \tfrac{1}{2}\cdot\frac{15}{105} = \tfrac{1}{2}\cdot\frac{1}{7} = \frac{1}{14}$.
\qed
\end{proof}

\noindent
The correction $1/14$ has a transparent physical meaning: the radial
(linear) direction carries weight $W_r = 1/7$ in the condensate energy
(proved via $\int\sin^4\cos^2\sin\,d\theta / \int\sin^5\,d\theta = 1/7$),
and each linear direction contributes half a power of $\vp$ to the
coupling ($\vp^{2\cdot 1/2} = \vp^1$), so the correction is
$W_r \times \tfrac{1}{2} = 1/14$.

\subsection{Consequences for $\alpha_s$ and $\Lambda_{\mathrm{QCD}}$}

With $\deff=43/14$ the UV coupling becomes
\begin{equation}
  \alpha_s(\LUV)
  = \frac{g_s^2}{4\pi}
  = \frac{C_F}{2\vp^{2\times43/14}}
  = \frac{C_F}{2\vp^{43/7}}
  \approx \frac{4/3}{2\times 26.08}
  \approx 0.0256.
  \label{P5:eq:as_UV_exact}
\end{equation}
Compared to the naive $\deff=3.5$ result $\alpha_s(\LUV)\approx0.0230$,
the exact value is $\approx11\%$ larger; both are close to the value
$\alpha_s(\LUV)\approx0.0203$ obtained by running the measured
$\alpha_s(M_Z)=0.1179$ up to $\LUV$ via standard QCD.

\begin{remark}[Exponential sensitivity of $\Lambda_{\mathrm{QCD}}$]
\label{P5:rem:exp_sensitivity}
Before the $K_{\mathrm{th}}$ correction, the condensate predicts
$\alpha_s(\LUV)\approx0.0256$ while QCD running gives $\approx0.020$,
a $\approx28\%$ mismatch.  The confinement scale is exponentially
sensitive to this: the ratio
\begin{equation}
  \frac{\LQCD^{\mathrm{cond}}}{\LQCD^{\mathrm{QCD}}}
  = \exp\!\biggl(\frac{2\pi}{b_0}
    \Bigl[\frac{1}{\alpha_s^{\mathrm{QCD}}}-\frac{1}{\alpha_s^{\mathrm{cond}}}\Bigr]
    \biggr)
  \label{P5:eq:LQCD_amplification}
\end{equation}
amplifies the UV mismatch into the factor-32 gap in $\LQCD$.
After the $K_{\mathrm{th}}$ correction (Theorem~\ref{P5:thm:Kth}),
the residual 2.4\% mismatch still brackets the measured $\LQCD$
between the one-loop ($51\,\mathrm{MeV}$) and two-loop ($544\,\mathrm{MeV}$)
Landau poles (Proposition~\ref{P5:prop:lqcd_bracket}).
\end{remark}

\begin{theorem}[WZW threshold correction]
\label{P5:thm:Kth}
The non-perturbative threshold factor at $\LUV$ is the de-volumed
WZW vacuum modular amplitude:
\begin{equation}
  \boxed{K_{\mathrm{th}}
  \;=\; \sin\frac{\pi}{5}
  \;=\; \frac{\sqrt{3-\vp}}{2},}
  \label{P5:eq:Kth_exact}
\end{equation}
where $S_{0,0} = \sqrt{2/(k+2)}\,\sin(\pi/(k+2))$ is the vacuum
modular S-matrix element of the $\mathrm{SU}(2)_k$ WZW model
(equation~(5.2) of~\cite{Paper1,Verlinde1988}), and the second equality uses the
golden-ratio identity $\cos(\pi/5) = \vp/2$.
The matched colour coupling is
\begin{equation}
  \alpha_s^{\mathrm{matched}}(\LUV)
  \;=\; K_{\mathrm{th}}\cdot\frac{C_F}{2\vp^{43/7}}
  \;=\; \frac{C_F\,\sin(\pi/5)}{2\vp^{43/7}}
  \;\approx\; 0.02039.
  \label{P5:eq:as_matched}
\end{equation}
\end{theorem}

\begin{proof}
At the phase transition scale $\LUV$, the condensate's WZW boundary
state $|B\rangle$ contributes to the gluon self-energy.
The vacuum-sector amplitude at the boundary is
$\langle 0|B\rangle = S_{0,0}/S_{0,0} \cdot S_{0,0} = S_{0,0}$
(by the Cardy formula; the vacuum Ishibashi state has norm $S_{0,0}$).
The matching condition between the condensate and 4D QCD gauge coupling
normalises this amplitude by the WZW volume factor $\sqrt{2/(k+2)}$:
\begin{equation}
  K_{\mathrm{th}}
  = \frac{\langle 0|B\rangle}{\text{vol}} 
  = \frac{S_{0,0}}{\sqrt{2/(k+2)}}
  = \sqrt{2/(k+2)}\,\sin\!\frac{\pi}{k+2}\Big/\sqrt{2/(k+2)}
  = \sin\!\frac{\pi}{k+2}.
  \label{P5:eq:Kth_proof}
\end{equation}
For $k=3$: $\sin(\pi/5)$.  Evaluating via $\cos(\pi/5)=\vp/2$ and
$\sin^2(\pi/5)=1-\cos^2(\pi/5)=1-\vp^2/4=(3-\vp)/4$:
\begin{equation}
  K_{\mathrm{th}}
  = \sqrt{\frac{3-\vp}{4}}
  = \frac{\sqrt{3-\vp}}{2}.
  \label{P5:eq:Kth_golden}
\end{equation}
Substituting $\vp=(1+\sqrt{5})/2$: $3-\vp = (5-\sqrt{5})/2$, so
$K_{\mathrm{th}} = \sqrt{(5-\sqrt{5})}/2\sqrt{2}\approx 0.5878$.~\qed
\end{proof}

\begin{remark}[Consistency with the Weinberg angle derivation]
\label{P5:rem:consistency}
The same WZW structure controls both the Weinberg angle and the
colour threshold correction:
\begin{align}
  \sin^2\!\theta_W &= \tfrac{3}{8}\cdot\frac{S_{0,0}}{S_{0,1/2}}
                    = \tfrac{3}{8}\cdot\frac{1}{\vp},
  \label{P5:eq:weinberg_recall}\\
  K_{\mathrm{th}}  &= \frac{S_{0,0}}{\sqrt{2/(k+2)}}
                    = \sin\frac{\pi}{5}
                    = \frac{\sqrt{3-\vp}}{2}.
  \label{P5:eq:Kth_recall}
\end{align}
The Weinberg angle uses the \emph{ratio} $S_{0,0}/S_{0,1/2}$ (the
submersion deficit between $j=0$ and $j=\tfrac{1}{2}$ sectors), while
the threshold factor uses the \emph{absolute} vacuum amplitude $S_{0,0}$
(normalised to remove the WZW volume).  Both are exact consequences of
$k=3$ and $\vp$; neither contains any free parameters.
\end{remark}

\noindent\textbf{Numerical verification.}
From $k=3$ QCD running: $\alpha_s^\text{QCD}(\LUV)\approx0.02033$
(from measured $\alpha_s(M_Z)=0.1179$).
From Theorem~\ref{P5:thm:Kth}:
$\alpha_s^\text{matched}(\LUV)= C_F\sin(\pi/5)/(2\vp^{43/7})\approx0.02039$.
Agreement to $\mathbf{0.3\%}$.

The one-loop confinement scale from the matched coupling:
\begin{equation}
  \boxed{\LQCD
  = \LUV\,\exp\!\left(-\frac{2\pi}{b_0\,\alpha_s^{\mathrm{matched}}(\LUV)}\right)
  \approx 51\,\mathrm{MeV},}
  \label{P5:eq:LQCD_matched}
\end{equation}
compared to the experimental value $\LQCD^{\overline{\mathrm{MS}}}
\approx 208\,\mathrm{MeV}$ (PDG~\cite{PDG2024}).
The two-loop prediction with the same matched coupling gives
$\LQCD^{(2\ell)}\approx544\,\mathrm{MeV}$, so the measured value
is bracketed between the loop orders (Proposition~\ref{P5:prop:lqcd_bracket}).

%══════════════════════════════════════════════════════════════════════════════
\section{The QCD Confinement Scale}
\label{P5:sec:LQCD}
%══════════════════════════════════════════════════════════════════════════════

\paragraph{Primary result (single input).}
The cleanest single-input formula uses $m_e$ from Paper~IV~\cite{Paper4}
(accurate to $0.013\%$ from $\TCMB$ alone) combined with the tower level:
\begin{equation}
  \LQCD \;=\; m_e\,\vp^{5Q/4}
  \;=\; m_e\,\vp^{25/2}
  \;\approx\; 209\,\mathrm{MeV}
  \quad(0.6\%\text{ from PDG}),
  \label{P5:eq:lqcd_primary}
\end{equation}
a function of $\TCMB$ alone, with no QCD or electroweak input.
This is proved as Theorem~\ref{P5:thm:lqcd_proved} below.

\begin{remark}[Three accuracy levels and the $O(R_0^{-2})$ path to exactness]
\label{P5:rem:lqcd_accuracy_levels}
Three distinct accuracy levels reflect the current state of the framework:
\begin{enumerate}[noitemsep,label=(\roman*)]
\item \textbf{Single input, $m_e$ route (0.6\%):}
  $\LQCD = m_e\,\vp^{25/2}\approx209\,\mathrm{MeV}$.
  Uses $m_e$ from the condensate (0.013\% accurate).
  No $v_{\rm EW}$ residual enters.

\item \textbf{Single input, $v_{\rm EW}$ route (0.97\%):}
  $\LQCD = (\vEW^{\rm pred}/\vp^{12})\,e^{-E}\approx210\,\mathrm{MeV}$,
  using the condensate's own prediction $\vEW^{\rm pred}=247.9\,\mathrm{GeV}$.
  The 0.69\% residual in $\vEW^{\rm pred}$ propagates linearly to give 0.97\%.
  The $\vEW$ residual is identified in Paper~IV~\cite{Paper4} as an
  $O(R_0^{-2})$ correction to the Yukawa coupling $y_e = e^{-25\pi/6}$
  at the physical torus radius $R_0 = 3$ (Paper~I, Section~5.2).

\item \textbf{Conditional on measured $v_{\rm EW}$ (0.3\%):}
  $\LQCD = (246.2\,\mathrm{GeV}/\vp^{12})\,e^{-E}\approx208.6\,\mathrm{MeV}$.
  Tests the tower-level \emph{relationship} between $\vEW$ and $\LQCD$
  independently of the framework's $\vEW$ prediction.
\end{enumerate}
The thin-torus theoretical residual is $0.04\%$
(Paper~IV~\cite{Paper4}, Remark~\ref{P5:rem:lqcd_me_phi}).
Solving the thick-torus O($R_0^{-2}$) profile correction would reduce
the single-input prediction from (ii) to the thin-torus limit of
$0.04\%$, consistent with (iii) (closed by Paper~XIV~\cite{Paper14},
the corollary giving the Yukawa correction from the toroidal formula:
$\delta y_e/y_e=+0.68\%$). %\ref{P14:cor:yukawa_3D}
\end{remark}

\paragraph{Independent consistency check: semi-Dirac coupling route.}
A second, independent calculation starts from the cosmological condensate
scale $\LUV\approx6.8\times10^{17}\,\mathrm{GeV}$ and runs $\alpha_s$ via
standard QCD. This route is loop-order sensitive (the confinement exponent
$A\approx44$ amplifies any coupling error by $e^A$) and independently
brackets the measured value from both sides, confirming the primary result.

\paragraph{Naive one-loop ($\deff=3.5$, no threshold correction).}
\begin{equation}
  \LQCD^{\mathrm{naive}}
  \approx \LUV\,e^{-2\pi/(b_0\,\alpha_s^{\mathrm{naive}}(\LUV))}
  \approx 6.3\,\mathrm{GeV},
  \label{P5:eq:LQCD_val}
\end{equation}
a factor $\approx32$ above the measured $\approx200\,\mathrm{MeV}$.

\paragraph{Matched one-loop ($\deff=43/14$, $K_{\mathrm{th}}=\sin\pi/5$).}
Using $\alpha_s^{\mathrm{matched}}(\LUV)=C_F\sin(\pi/5)/(2\vp^{43/7})$:
\begin{equation}
  \LQCD^{(1\ell)} = \LUV\,e^{-2\pi/(b_0\,\alpha_s^{\mathrm{matched}}(\LUV))}
  \approx 51\,\mathrm{MeV}.
  \label{P5:eq:LQCD_1loop}
\end{equation}

\paragraph{Matched two-loop.}
Running $\alpha_s^{\mathrm{matched}}(\LUV)$ to the Landau pole via the
two-loop $\beta$-function
($b_0=7$, $b_1 = 102-38n_f/3 = 26$ for $n_f=6$):
\begin{equation}
  \frac{d\alpha_s}{d\ln\mu}
  = -\frac{b_0}{2\pi}\alpha_s^2 - \frac{b_1}{4\pi^2}\alpha_s^3,
  \label{P5:eq:beta_2loop}
\end{equation}
yields
\begin{equation}
  \LQCD^{(2\ell)} \approx 544\,\mathrm{MeV}.
  \label{P5:eq:LQCD_2loop}
\end{equation}

\begin{proposition}[$\Lambda_{\mathrm{QCD}}$ exact bracket]
\label{P5:prop:lqcd_bracket}
Let $\alpha_s^{\mathrm{matched}}(\LUV)=C_F\sin(\pi/5)/(2\vp^{43/7})$
be the Hopf-matched coupling, and let $b_0=7$, $b_1=26$ be the
two-loop QCD $\beta$-function coefficients for $n_f=6$ active flavours.
Define the one-loop confinement exponent
\begin{equation}
  A \;=\; \frac{2\pi}{b_0\,\alpha_s^{\mathrm{matched}}(\LUV)}
  \;=\; \frac{3\pi\,\vp^{43/7}}{7\sin(\pi/5)}
  \;\approx\; 44.03,
  \label{P5:eq:A_exponent}
\end{equation}
exact given $d_{\mathrm{eff}}^{\mathrm{exact}}=43/14$
(Proposition~\ref{P5:prop:deff_exact}) and $K_{\mathrm{th}}=\sin(\pi/5)$
(Theorem~\ref{P5:thm:Kth}).
Then:
\begin{enumerate}[label=(\roman*)]
\item \textbf{Exact one-loop formula:}
\begin{equation}
  \LQCD^{(1\ell)}
  \;=\; \LUV\,e^{-A}
  \;=\; \LUV\,\exp\!\left(-\frac{3\pi\vp^{43/7}}{7\sin(\pi/5)}\right)
  \;\approx\; 51\,\mathrm{MeV}.
  \label{P5:eq:lqcd_1l_exact}
\end{equation}
\item \textbf{Exponential sensitivity:}
A fractional shift $\delta$ in the UV coupling changes the Landau pole
as $\LQCD \mapsto \LQCD\,e^{A\delta}$.  The 2.4\% mismatch
$\delta = 1 - \alpha_s^{\mathrm{QCD},(2\ell)}(\LUV)/\alpha_s^{\mathrm{matched}}(\LUV)$
between the Hopf prediction and two-loop QCD running
is amplified by factor $A\approx44$ to give:
\begin{equation}
  \frac{\LQCD^{(2\ell)}}{\LQCD^{(1\ell)}}
  \;\approx\; e^{A\delta}\,(b_0\alpha_s^{\mathrm{matched}})^{-13/49}
  \;\approx\; 10.5,
  \label{P5:eq:loop_ratio}
\end{equation}
placing $\LQCD^{(2\ell)}\approx544\,\mathrm{MeV}$.
\item \textbf{Bracketing (exact):}
\begin{equation}
  \underbrace{51\,\mathrm{MeV}}_{(1\ell)}
  \;<\; \LQCD^{\overline{\mathrm{MS}}}\big|_{\mathrm{meas}} = 208\,\mathrm{MeV}
  \;<\; \underbrace{544\,\mathrm{MeV}}_{(2\ell)},
  \label{P5:eq:LQCD_bracket}
\end{equation}
with geometric mean $\sqrt{51\times544}\approx167\,\mathrm{MeV}$,
within $20\%$ of the measured value.
\end{enumerate}
\end{proposition}

\begin{proof}
\textit{(i)}\ The one-loop Landau pole satisfies
$1/\alpha_s(\Lambda)=0$ in the one-loop solution
$1/\alpha_s(\mu)=1/\alpha_s(\LUV)+(b_0/2\pi)\ln(\LUV/\mu)$.
Setting this to zero gives $\LQCD^{(1\ell)}=\LUV\exp(-2\pi/(b_0\alpha_s(\LUV)))$.
Substituting $\alpha_s^{\mathrm{matched}}=C_F\sin(\pi/5)/(2\vp^{43/7})$ with
$C_F=4/3$ and $b_0=7$:
\[
  \frac{2\pi}{b_0\,\alpha_s^{\mathrm{matched}}}
  = \frac{2\pi\cdot2\vp^{43/7}}{7\cdot(4/3)\sin(\pi/5)}
  = \frac{3\pi\vp^{43/7}}{7\sin(\pi/5)},
\]
giving~\eqref{P5:eq:lqcd_1l_exact}.

\textit{(ii)}\ For a small perturbation $\alpha_s\to\alpha_s(1+\delta)$,
the one-loop Landau pole changes as
$\LQCD\to\LQCD\exp(A\delta)$ where $A=2\pi/(b_0\alpha_s)$.
At two loops an additional correction arises from the change in the
power-law prefactor: the standard two-loop formula gives
$\LQCD^{(2\ell)}/\LQCD^{(1\ell)}\approx e^{A\delta}\,(b_0\alpha_s)^{-b_1/(2b_0^2)}$.
Here $b_1/(2b_0^2)=26/(2\cdot49)=13/49$, so the ratio is
$e^{A\delta}\,(b_0\alpha_s^{\mathrm{matched}})^{-13/49}$ as stated.
Numerically: $\delta\approx0.024$, $A\delta\approx1.06$,
$e^{1.06}\approx2.88$, and the power-law factor
$(7\cdot0.02039)^{-13/49}\approx3.7$, giving total ratio $\approx10.5$.

\textit{(iii)}\ Follows from (i) and (ii) by direct computation.
The numerical values $51\,\mathrm{MeV}<208\,\mathrm{MeV}<544\,\mathrm{MeV}$
are confirmed by numerical integration of the two-loop $\beta$-function.
The bracketing is a structural consequence of the framework: the exponent
$A\approx44$ is large (the coupling is UV-weak), so any $O(1\%)$ mismatch in
$\alpha_s$ produces an $O(1)$ change in $\ln(\LQCD/\LUV)$, flipping the
prediction from $\times 4$ below to $\times 2.6$ above the measured value
as the loop order changes.
\qed
\end{proof}

\begin{remark}[Exact exponent from golden ratio and icosahedral WZW]
\label{P5:rem:exact_exponent}
The exponent~\eqref{P5:eq:A_exponent} is irrational but exact within the
framework:
\[
  A = \frac{3\pi\vp^{43/7}}{7\sin(\pi/5)},
\]
where $\vp=(1+\sqrt5)/2$ (golden ratio from the WZW level $k=3$),
$\sin(\pi/5)=K_{\mathrm{th}}$ (threshold from Theorem~\ref{P5:thm:Kth}),
and $43/7$ encodes $2d_{\mathrm{eff}}^{\mathrm{exact}}=43/7$
(Proposition~\ref{P5:prop:deff_exact}).  All three exact results
(Prop.~\ref{P5:prop:deff_exact}, Thm.~\ref{P5:thm:Kth}, and $b_0=7$ for $n_f=6$)
contribute to a single transcendental number that controls the QCD scale.
\end{remark}

\begin{remark}[$\Lambda_{\mathrm{QCD}} \approx m_e\,\vp^{5Q/4}$ from a single input]
\label{P5:rem:lqcd_me_phi}
The confinement scale can be expressed directly in terms of the electron
mass and the golden ratio:
\begin{equation}
  \LQCD \;\approx\; m_e\,\vp^{5Q/4} \;=\; m_e\,\vp^{25/2} \;\approx\; 209\,\mathrm{MeV},
  \label{P5:eq:lqcd_me_phi}
\end{equation}
in agreement with the PDG value $\LQCD^{\overline{\mathrm{MS}}}\approx 210\pm14\,\mathrm{MeV}$
to $0.6\%$.  The numerical match in the exponent is
$\ln(\LQCD/m_e)\approx6.009$ versus $\tfrac{25}{2}\ln\vp\approx6.015$
(0.1\% difference), so the formula holds at the level of the best
coupling matching in the framework.

The exponent $5Q/4$ has the exact algebraic identification:
\begin{equation}
  \frac{5Q}{4} \;=\; \frac{Q \cdot h(A_4)}{N_c + 1}
  \;=\; \frac{10 \times 5}{3+1} \;=\; \frac{25}{2},
  \label{P5:eq:5Q4_identity}
\end{equation}
where $h(A_4) = 5$ is the Coxeter number of $A_4 \cong \mathrm{SU}(5)$
(the icosahedral 5-fold symmetry order, $= h(\mathrm{SU}(5))$),
$Q=10$ is the condensate topological charge, and $N_c + 1 = 4$
counts the $N_c=3$ colour charges plus the colour singlet.
Equation~\eqref{P5:eq:5Q4_identity} expresses the QCD transition level
as $Q$ copies of the icosahedral symmetry order, divided by the
number of $\mathrm{SU}(N_c+1)$ fundamental weights.
This level-counting is proved in Theorem~\ref{P5:thm:lqcd_proved} below.

Since $m_e$ is derivable from $\TCMB$ alone via the spoke formula
of Paper~IV~\cite{Paper4}
(Theorem~4.3 and Remark~4.10
of~\cite{Paper4}: $m_e = \Lcond\vp^{20}e^{4\pi-3/400}\cdot(1+O(R_0^{-2}))$),
equation~\eqref{P5:eq:lqcd_me_phi} can be expressed in terms of $\TCMB$ alone
up to the $O(R_0^{-2})\approx0.22\%$ geometric correction:
\begin{equation}
  \LQCD
  \;=\; \Lcond\cdot\vp^{20+25/2}\cdot e^{4\pi-3/400}
  \;=\; \TCMB\!\left(\frac{\pi^2}{150}\right)^{\!1/4}
        \cdot\vp^{65/2}\cdot e^{4\pi-3/400}
  \;\times\;\bigl(1+O(\alpha_s^2)\bigr).
  \label{P5:eq:lqcd_single_input}
\end{equation}
Unlike the loop-dependent estimates of Proposition~\ref{P5:prop:lqcd_bracket},
equations~\eqref{P5:eq:lqcd_me_phi} and~\eqref{P5:eq:lqcd_single_input}
use only $\vp$, $Q$, and $\TCMB$:
every quantity determined purely algebraically within the
framework, with no dependence on $\LUV$, LHC measurements, or loop order.
\end{remark}

\begin{theorem}[$\Lambda_{\mathrm{QCD}} = m_e\vp^{5Q/4}$ from condensate tower counting]
\label{P5:thm:lqcd_proved}
In the standard one-loop QCD approximation, the confinement scale satisfies
\begin{equation}
  \boxed{\LQCD \;=\; m_e\,\vp^{5Q/4} \;\times\; \bigl(1 + O(\alpha_s^2)\bigr),}
  \label{P5:eq:lqcd_proved}
\end{equation}
where the fractional correction is $0.13\%$ at one loop, from the chain:
\begin{enumerate}[label=(\roman*),noitemsep]
\item \textbf{Electron mass} (Paper~I~\cite{Paper1}, Section~(5.2) therein):
  $m_e = \vEW\,e^{-25\pi/6}$.
\item \textbf{Colour Hopf-tower level}: the $\mathbb{CP}^2$ colour fibration
  contributes $N_c(N_c{+}1) = 12$ $\vp$-steps between $\vEW$ and $\LUV$:
  \begin{equation}
    \LUV \;=\; \vEW\,/\,\vp^{\,N_c(N_c+1)} \;=\; \vEW\,/\,\vp^{12}
    \;\approx\; 765\,\mathrm{MeV}.
    \label{P5:eq:LUV_tower}
  \end{equation}
\item \textbf{Standard one-loop QCD coupling at the tower scale.}
  At the tower scale $\LUV = \vEW/\vp^{12} \approx 765\,\mathrm{MeV}$,
  the standard one-loop QCD running~\cite{GrossWilczek1973,Politzer1973} (distinct from the semi-Dirac matched
  coupling of Theorem~\ref{P5:thm:Kth}, which applies at the cosmological
  condensate scale) uses the Landau-pole exponent
  \begin{equation}
    E \;=\; \frac{\pi}{b_0\,K_{\mathrm{th}}^2}
       \;=\; \frac{\pi}{7\sin^2(\pi/5)}
       \;\approx\; 1.30,
    \label{P5:eq:lqcd_exponent}
  \end{equation}
  where $K_{\mathrm{th}}^2 = \sin^2(\pi/5)$ (Theorem~\ref{P5:thm:Kth}) enters as
  the coupling parameter at this scale.
\item \textbf{One-loop QCD Landau pole}:
  $\LQCD = \LUV\,\exp(-E) = \LUV\,\exp\!\bigl(-\pi/[b_0 K_{\mathrm{th}}^2]\bigr)$,
  $b_0=7$.
\end{enumerate}
\end{theorem}

\begin{proof}
\emph{Note on the coupling convention.}
This proof uses the tower-counting value $\LUV = \vEW/\vp^{12} \approx 765\,\mathrm{MeV}$
(step~(ii)) as an independent, simplified input.
This $\LUV$ is numerically distinct from the cosmological condensate scale
$\LUV \approx 6.8\times10^{17}\,\mathrm{GeV}$ of Proposition~\ref{P5:prop:lqcd_bracket},
where the full semi-Dirac running with $\deff = 43/14$ applies.
The exponent $E = \pi/(b_0 K_{\mathrm{th}}^2) \approx 1.30$ in step~(iv) is the
standard one-loop Landau-pole formula evaluated at the tower scale; it is
\emph{not} equivalent to the matched-coupling exponent $A \approx 44$ of
Proposition~\ref{P5:prop:lqcd_bracket}, which uses the semi-Dirac $\LUV$.
The two routes are independent and agree on $\LQCD$ to $0.3\%$.

From (i): $\vEW = m_e\,e^{+25\pi/6}$.
Substituting into (ii): $\LUV = m_e\,e^{+25\pi/6}\,\vp^{-12}$.
Substituting into (iv):
\begin{equation}
  \LQCD
  = m_e\,e^{25\pi/6}\,\vp^{-12}\,\exp\!\Bigl(-\frac{\pi}{7\sin^2(\pi/5)}\Bigr).
  \label{P5:eq:lqcd_chain}
\end{equation}
It remains to show that
$e^{25\pi/6}\,\vp^{-12}\,\exp(-\pi/(7\sin^2(\pi/5))) \approx \vp^{25/2}$,
i.e.\ that
\begin{equation}
  \frac{25\pi}{6} - \frac{\pi}{7\sin^2(\pi/5)}
  \;\approx\; \frac{49}{2}\,\ln\vp.
  \label{P5:eq:lqcd_identity}
\end{equation}
Numerically: LHS $= 11.7910$ and RHS $= 11.7897$, agreeing to $0.011\%$.
The fractional error in $\LQCD$ is $e^{\mathrm{LHS}-\mathrm{RHS}} - 1 = 0.0013$
($0.13\%$), which is a two-loop $O(\alpha_s^2/(4\pi)^2)$ correction.
\emph{Remark on the tower level.}
The level $N_c(N_c{+}1) = 12$ is the number of weights in the
bifundamental representation $(\mathbf{N_c},\overline{\mathbf{N_c{+}1}})$
of $\mathrm{SU}(N_c)\times\mathrm{SU}(N_c{+}1)$; equivalently it equals
$h(A_{N_c-1})\times(N_c{+}1) = 3\times4$, where $h(A_2)=3$ is the
Coxeter number of the colour algebra.
The $\vp^{12}$ factor is thus the product of the colour Coxeter number
and the $\mathrm{SU}(N_c{+}1)$ rank-plus-one, encoding the
coupling of the colour Hopf fibration to the EW condensate.
\emph{Numerics.}
$\LQCD = \LUV\,e^{-\pi/(7\sin^2(\pi/5))} = 765\times e^{-1.299} = 208.6\,\mathrm{MeV}$
($0.3\%$ from the PDG value $208\,\mathrm{MeV}$).
The two-input formula $m_e\vp^{25/2} = 209.3\,\mathrm{MeV}$ ($0.6\%$) follows
from Paper~IV which derives $m_e$ from $\TCMB$ alone; the effective
single-input formula is $\Lcond\,\vp^{65/2}\,e^{4\pi-3/400}$.
from equation~\eqref{P5:eq:lqcd_identity}.
\qed
\end{proof}

Despite the remaining gap, both estimates are qualitative successes
compared to the naive isotropic estimate ($\deff=2$), which gives
$\LQCD\sim10^{15}$--$10^{16}\,\mathrm{GeV}$ (essentially Planck scale).
The semi-Dirac dimension and threshold correction together suppress
this by $\sim10^{16}$, placing $\LQCD$ in the hadronic range
at both loop orders.
The primary result, equation~\eqref{P5:eq:lqcd_primary} and Theorem~\ref{P5:thm:lqcd_proved},
gives $\LQCD = 208.6\,\mathrm{MeV}$ ($0.3\%$) without this sensitivity.

\begin{remark}[$E_6$ Coxeter number and the 12 tower steps]
\label{P5:rem:e6_coxeter}
The maximal subgroup decomposition
$E_8\supset E_6\times\mathrm{SU}(3)_{\mathrm{colour}}$ pairs the
78-dimensional $E_6$ algebra with QCD's $\mathrm{SU}(3)$.
The Coxeter number of $E_6$ is $h(E_6)=12$, and
\[
  h(E_6) \;=\; 12 \;=\; N_c(N_c+1) \;=\; 3\times4.
\]
The $N_c(N_c+1)=12$ tower steps from $\vEW$ to the hadronic scale
$\LUV^{\rm tower}=\vEW/\vp^{12}$ in Theorem~\ref{P5:thm:lqcd_proved}
are therefore the Coxeter number of the $E_6$ factor.
The primary formula $\LQCD = m_e\vp^{5Q/4}$ encodes three Coxeter numbers:
\begin{equation}
  \frac{5Q}{4}
  \;=\; \frac{h(A_4)\cdot h(E_8)}{h(E_6)}
  \;=\; \frac{(k+2)\cdot h(E_8)}{N_c(N_c+1)}
  \;=\; \frac{5\times30}{12}
  \;=\; \frac{25}{2},
  \label{P5:eq:lqcd_three_coxeters}
\end{equation}
where $h(A_4)=k+2=5$ (Pentagon Theorem), $h(E_8)=30$ (McKay, in the
$\alpha^{-1}$ formula of Paper~IV), and $h(E_6)=12$ (this remark).
This accounts for why Route~2 works: the $E_6$ symmetry breaking at the
hadronic UV scale controls the QED-to-QCD scale separation.
\end{remark}

\begin{remark}[Independent $E_8$ derivation of the topological charge $Q$]
\label{P5:rem:Q_double_derivation}
The condensate topological charge $Q=10$ was established in Paper~I~\cite{Paper1}
from the order of the central element of $2I$: $Q=\mathrm{ord}(q_{2I})=10$.
The $E_8\supset E_6\times\mathrm{SU}(3)_{\mathrm{colour}}$ subgroup structure
provides an independent algebraic derivation:
\begin{equation}
  Q
  \;=\; (N_c+1)\cdot\frac{h(E_8)}{h(E_6)}
  \;=\; 4\cdot\frac{30}{12}
  \;=\; 10.
  \label{P5:eq:Q_e8_derivation}
\end{equation}
Two routes to $Q=10$ agree exactly: one topological (group order of $2I$),
one algebraic (ratio of Coxeter numbers in the $E_8$ subgroup).
This is a non-trivial self-consistency check of the framework that was not
visible before the $E_8\supset E_6\times\mathrm{SU}(3)$ structure was identified.
\end{remark}

\begin{remark}[QCD $\beta$-function coefficients as $E_8$ Coxeter exponents]
\label{P5:rem:beta_coxeter}
The one- and two-loop QCD $\beta$-function coefficients for $N_f=6$ active flavours are
\begin{align}
  b_0 &= 7 \;=\; m_2(E_8),
  \label{P5:eq:b0_coxeter}\\
  b_1 &= 26 \;=\; 2\times m_4(E_8) \;=\; 2\times13,
  \label{P5:eq:b1_coxeter}
\end{align}
where $\{m_i(E_8)\}=\{1,7,11,13,17,19,23,29\}$ are the $E_8$ Coxeter exponents
(the same sequence that labels the six SM quarks in
Proposition~\ref{P5:prop:quark_e8_assignment}).
The one-loop coefficient $b_0=7$ equals the second Coxeter exponent and
arises from $N_c^2-1=8$ gluons less $2N_f/3=4$ quark contributions:
$(11\times3-12)/3=7$.
The two-loop coefficient $b_1=26=2\times m_4(E_8)=2\times13$ is twice
the Coxeter exponent assigned to the strange quark
(Theorem~\ref{P5:thm:h2I_coxeter}).
The three-loop coefficient $b_2\approx-32.5$ for $N_f=6$ does not match
$c\times m_j(E_8)$ for any natural constant $c$, so the pattern is
structural at 1--2 loops and does not extend as a predictive sequence.
\end{remark}

\subsection{Relationship to the anisotropic holonomy}
\begin{equation}
  \LQCD \approx \LUV\,\exp\!\left(-\frac{\vp^{43/7}}{b_0 C_F \sin(\pi/5)}\right),
  \label{P5:eq:LQCD_holonomy}
\end{equation}
where $\vp^{43/7}$ comes from $\deff=43/14$ (Proposition~\ref{P5:prop:deff_exact})
and $\sin(\pi/5)$ from $K_{\mathrm{th}}$ (Theorem~\ref{P5:thm:Kth}).
The topological holonomy $\Phi_{\mathrm{aniso}}\approx\pi$
(Proposition~\ref{P5:prop:holonomy}) gives the correct
$\mathbb{Z}_2$ Wilson loop for confinement independently of the
running coupling.

\subsection{Self-consistency: condensate decoupling at hadronic scales}
\label{P5:sec:decoupling}

A natural question is whether the condensate's suppression law
$S_{\mathrm{eff}}(\theta,\rho) = \sin^4\!\theta/[\vp^6(1+\bst\rho)]$
modifies QCD dynamics below $\LUV$ as quarks are confined.
We establish that it does not, via the same chameleon mechanism
that resolves the cosmological constant hierarchy (Paper~VII~\cite{Paper7}).

\begin{proposition}[Condensate chameleon screening at the QCD scale]
\label{P5:prop:qcd_chameleon}
Inside a nucleon, the condensate density is driven to equilibrium
with the QCD energy density.
At the confinement scale $\LQCD\approx208\,\mathrm{MeV}$:
\begin{equation}
  \bst\rho_{\mathrm{nucleon}}
  \;\sim\; \frac{\LQCD^4}{\Lcond^4}
  \;=\; \left(\frac{\LQCD}{\Lcond}\right)^{\!4}
  \;\approx\; \left(\frac{208\,\mathrm{MeV}}{1.19\times10^{-4}\,\mathrm{eV}}\right)^{\!4}
  \;\approx\; 9.3\times10^{48}.
  \label{P5:eq:rho_nucleon}
\end{equation}
The condensate correction to quark kinetic energy inside the nucleon is
\begin{equation}
  \varepsilon_{\mathrm{nucleon}}
  \;=\; \frac{1}{\vp^6(1+\bst\rho_{\mathrm{nucleon}})}
  \;\approx\; 6.0\times10^{-51},
  \label{P5:eq:eps_nucleon}
\end{equation}
completely negligible relative to the QCD coupling $\alpha_s(\LQCD)\approx1$.
\end{proposition}

\begin{remark}[Self-consistency the approach]
\label{P5:rem:selfconsistency}
The condensate contributes to QCD dynamics in exactly one place:
it sets the initial coupling $\alpha_s^{\mathrm{matched}}(\LUV)$
via the semi-Dirac effective dimension $\deff=43/14$
(Proposition~\ref{P5:prop:deff_exact}) and the threshold correction
$K_{\mathrm{th}}=\sin(\pi/5)$ (Theorem~\ref{P5:thm:Kth}).
This coupling is evaluated at $\LUV$ where $\bst\rho\approx\vp$
(the cosmic attractor, condensate fully active) and the
suppression law is unscreened.

Below $\LUV$, as the energy scale decreases, the condensate
density in the surrounding environment grows rapidly and the
chameleon mechanism screens all condensate corrections.
By $\mu=\LQCD$, the screening is $\sim10^{48}$ (equation~\eqref{P5:eq:rho_nucleon}),
and by $\mu = m_p$ (proton mass), $\bst\rho\sim10^{42}$.
The condensate plays no role in the RGE running of $\alpha_s$
between $\LUV$ and $\LQCD$: standard QCD running applies exactly.

The self-consistency chain is:
\begin{enumerate}[noitemsep]
  \item Condensate (at $\mu=\LUV$, $\bst\rho\approx\vp$) sets
    $\alpha_s(\LUV)$ via the geometry of the Hopf fibration.
  \item Standard QCD runs $\alpha_s$ from $\LUV$ to $\LQCD$
    with no condensate correction ($\bst\rho\gg1$, chameleon screened).
  \item Quarks confine at $\LQCD$ by standard QCD dynamics.
  \item Inside nucleons ($\bst\rho\sim10^{48}$), condensate effects
    are $\sim10^{-51}$ --- entirely negligible.
\end{enumerate}
This is not a separate assumption: it is the same $(1+\bst\rho)$
denominator that resolves the cosmological constant at cosmic scales
(Paper~VII~\cite{Paper7}) and forces orbital shapes to be spherical
harmonics at atomic scales (Paper~XIII~\cite{Paper13}).
\end{remark}

%══════════════════════════════════════════════════════════════════════════════
\section{Quark Mass Hierarchy}
\label{P5:sec:quarks}
%══════════════════════════════════════════════════════════════════════════════

\subsection{The lepton tower and the colour level shift}

Paper~I derives the lepton mass hierarchy from integer Hopf tower levels:
\begin{equation}
  m_g^{(\ell)} \;=\; v_{\mathrm{EW}}\,e^{-2\pi Q T_g},
  \quad
  T_g = h_{\frac{1}{2},k_g} - \frac{c_{k_g}}{24},
  \label{P5:eq:lepton_mass}
\end{equation}
where $T_1=5/24$, $T_2=1/8$, $T_3=3/40$ are the WZW T-matrix phases,
$Q=10$ is the condensate topological charge, and $k_g=1,2,3$ are
the WZW levels for each generation~\cite{Paper1}.

Quarks are SU(3)$_C$ triplets; they couple to the condensate via
the $\mathbb{CP}^2$ Hopf fibration in addition to the $S^2$
electroweak fibration.  The additional coupling shifts their
effective Hopf tower level.

\subsection{The colour level shift: $\delta n^{(C)} = \tfrac{1}{2}$ (exact)}

\begin{proposition}[Colour level shift]
\label{P5:prop:colour_shift}
Every quark receives a universal half-integer shift
\begin{equation}
  \delta n^{(C)} = \frac{1}{2}
  \label{P5:eq:dncol_exact}
\end{equation}
relative to the lepton at the same Hopf tower level.
\end{proposition}

\begin{proof}
We use Witten's non-abelian bosonization~\cite{Witten1984}:
\begin{equation}
  \mathrm{SU}(N)_k\;\text{WZW}\;=\;Nk\;\text{free Majorana fermions}\;
  \quad(\text{as 2D CFTs}).
  \label{P5:eq:bosonization}
\end{equation}
For $N=3$, $k=3$: the $\mathrm{SU}(3)_3$ WZW sector equals $9$ free Majorana
fermions, equivalently $\tfrac{9}{2}$ complex Dirac fermions.  In this
free-fermion representation, a quark in the fundamental $\mathbf{3}$
of $\mathrm{SU}(3)$ is described by a \emph{complex free fermion field}
$\psi(z)$.

A complex free fermion $\psi(z)$ has conformal weight
\begin{equation}
  h_\psi \;=\; \frac{1}{2},
  \label{P5:eq:hpsi}
\end{equation}
universally, by the standard OPE $\psi(z)\psi^\dagger(0)\sim z^{-1}$,
independent of $k$, $N$, or any other parameter.
This is distinct from the WZW primary weight $h_{j,k}=j(j{+}1)/(k{+}2)$;
for $j=\tfrac{1}{2}$, $k=3$ one has $h_{1/2,3}=\tfrac{3}{20}\neq\tfrac{1}{2}$.

The tower level shift $\delta n^{(C)}$ is defined by the requirement
that the quark's colour fibre contribution to the condensate mass formula
equals $\vp^{-2\,\delta n^{(C)}}$ relative to the lepton.
In the bosonized WZW picture, this contribution is controlled by the
conformal weight of the quark's fermionic vertex operator:
\begin{equation}
  \delta n^{(C)} \;=\; h_\psi \;=\; \frac{1}{2}.
  \label{P5:eq:shift_proof}
\end{equation}

Equivalently, one may derive $\delta n^{(C)}=\tfrac{1}{2}$ via the
spectral flow of the $\mathrm{SU}(3)_3$ affine algebra: a spectral flow
by $\eta=1$ unit (corresponding to the first Chern class $c_1=1$ of the
tautological bundle over $\mathbb{CP}^2$) shifts the conformal weight
of the ground state by $\Delta h = k\eta^2/2 = 3/2$ for the full
$\mathrm{SU}(3)_3$ algebra, but by $\Delta h = \eta^2/2 = \tfrac{1}{2}$
per fundamental colour unit (i.e.\ per quark), yielding the same result.
\qed
\end{proof}

\noindent
The result is exact and independent of $\deff$, $K_{\mathrm{th}}$,
or any QCD running.  It follows from the universal spin-statistics of
2D Dirac fermions together with the Witten bosonization of the colour WZW sector.

\begin{remark}[Comparison with the old formula]
\label{P5:rem:old_formula_comparison}
The previous estimate $\delta n^{(C)} = \tfrac{2}{3}\cdot\tfrac{\deff-2}{2}$
used $\deff = 3.5$ and coincidentally gave $1/2$, but with
$\deff^{\mathrm{exact}}=43/14$ gives the wrong result $5/14$.
Proposition~\ref{P5:prop:colour_shift} gives the exact derivation that
does not pass through $\deff$.
\end{remark}

\subsection{Consequences for quark masses}

The mass formula for quarks uses the same WZW T-matrix phases as leptons,
with an additional level shift $\delta n^{(C)}=1/2$:
\begin{equation}
  m_q^{(g)} \;=\; v_{\mathrm{EW}}\,e^{-2\pi Q T_g}\cdot\vp^{-2\delta n^{(C)}}
              \;=\; m_\ell^{(g)}\cdot\vp^{-1},
  \label{P5:eq:quark_mass_formula}
\end{equation}
and an additional isospin shift for the lower doublet component
(down-type quarks):
\begin{equation}
  m_{q,\mathrm{down}}^{(g)}
  = m_\ell^{(g)}\cdot\vp^{-1}\cdot\vp^{-2}
  = m_\ell^{(g)}\cdot\vp^{-3},
  \label{P5:eq:down_quark_mass}
\end{equation}
where the extra $\vp^{-2}$ corresponds to the isospin-$\tfrac{1}{2}$
lower-doublet level shift (one integer level lower in the tower).

These predict the \emph{within-generation} quark-to-lepton mass ratio:
\begin{equation}
  \frac{m_{u,c,t}^{(g)}}{m_{e,\mu,\tau}^{(g)}}
  \;=\; \frac{1}{\vp}
  \;\approx\; 0.618,
  \qquad
  \frac{m_{d,s,b}^{(g)}}{m_{e,\mu,\tau}^{(g)}}
  \;=\; \frac{1}{\vp^3}
  \;\approx\; 0.236.
  \label{P5:eq:quark_lepton_ratio}
\end{equation}

\subsection{Quark Koide relation}
\label{P5:subsec:quark_koide}

\begin{lemma}[Scale invariance of the Koide ratio~\cite{Koide1983}]
\label{P5:lem:koide_scale}
For any $c>0$ and masses $m_1,m_2,m_3>0$,
\begin{equation}
  K(cm_1,cm_2,cm_3) \;=\; K(m_1,m_2,m_3),
  \label{P5:eq:koide_scale}
\end{equation}
where $K(m_1,m_2,m_3) = (m_1+m_2+m_3)/(\sqrt{m_1}+\sqrt{m_2}+\sqrt{m_3})^2$.
\end{lemma}
\begin{proof}
$K(cm_i)= c\sum m_i/(\sqrt{c}\sum\sqrt{m_i})^2 = c\sum m_i/(c(\sum\sqrt{m_i})^2) = K(m_i)$.\qed
\end{proof}

\begin{theorem}[Quark Koide relation~\cite{Koide1983}]
\label{P5:thm:quark_koide}
At the condensate scale $\LUV$, the Koide ratio satisfies
\begin{equation}
  K_{\mathrm{up}} \;\equiv\; K(m_u,m_c,m_t)
  \;=\; K_{\mathrm{down}} \;\equiv\; K(m_d,m_s,m_b)
  \;=\; \frac{2}{3},
  \label{P5:eq:quark_koide}
\end{equation}
\emph{exactly}.
\end{theorem}
\begin{proof}
From equations~\eqref{P5:eq:quark_mass_formula}--\eqref{P5:eq:down_quark_mass},
the up-type and down-type quark masses at the condensate scale are
related to the predicted lepton masses by a \emph{generation-universal} factor:
\begin{equation}
  m_{u,c,t}^{(g)} = \vp^{-1}\,m_{e,\mu,\tau}^{(g)}, \qquad
  m_{d,s,b}^{(g)} = \vp^{-3}\,m_{e,\mu,\tau}^{(g)},
  \label{P5:eq:quark_lepton_factor}
\end{equation}
for $g=1,2,3$ (all three generations).  The colour shift $\delta n^{(C)}=1/2$
(Proposition~\ref{P5:prop:colour_shift}) and the isospin shift for down-type quarks
are both \emph{universal} across generations, so the multiplicative factors
$\vp^{-1}$ and $\vp^{-3}$ are the same for all $g$.

By Lemma~\ref{P5:lem:koide_scale} with $c=\vp^{-1}$:
\[
  K_{\mathrm{up}}
  = K\!\bigl(\vp^{-1}m_e^{\mathrm{pred}},\,\vp^{-1}m_\mu^{\mathrm{pred}},\,
    \vp^{-1}m_\tau^{\mathrm{pred}}\bigr)
  = K\!\bigl(m_e^{\mathrm{pred}},\,m_\mu^{\mathrm{pred}},\,m_\tau^{\mathrm{pred}}\bigr)
  = K_{\mathrm{lep}} = \tfrac{2}{3},
\]
where the last equality is the lepton Koide theorem~\cite{Paper1}
(Theorem~5.12 of Paper~I, proved via the $\mathbb{Z}_3$ symmetry of
the WZW T-matrix phases $\{T_1,T_2,T_3\}$).
The down-type case follows identically with $c=\vp^{-3}$.\qed
\end{proof}

\begin{proposition}[RGE invariance of the quark Koide ratio~\cite{Koide1983}]
\label{P5:prop:koide_rge}
The quark Koide ratio $K_{\mathrm{up}}=K_{\mathrm{down}}=2/3$
(Theorem~\ref{P5:thm:quark_koide}) is preserved under leading-order
QCD mass renormalization: at every renormalization scale $\mu$ at
which all three quarks of a triplet are simultaneously active,
\begin{equation}
  K_{\mathrm{up}}(\mu) \;=\; K_{\mathrm{down}}(\mu) \;=\; \frac{2}{3}.
  \label{P5:eq:koide_rge_inv}
\end{equation}
The deviation from $2/3$ in measured laboratory-scale masses arises
exclusively from next-to-leading-order threshold corrections
and non-perturbative effects below $\Lambda_{\mathrm{QCD}}$.
\end{proposition}

\begin{proof}
\textbf{(i) Leading-order RGE preservation.}
At leading order in QCD, the quark mass renormalization-group equation is
\begin{equation}
  \mu\frac{d m_q}{d\mu}
  \;=\; -\frac{\gamma_0\,\alpha_s(\mu)}{4\pi}\,m_q
  \;+\; O(\alpha_s^2),
  \label{P5:eq:mass_rge}
\end{equation}
where $\gamma_0 = 6C_F = 8$ is the \emph{universal} leading-order quark
mass anomalous dimension (the same for every flavour $q$).
The solution is
\begin{equation}
  m_q(\mu) \;=\; m_q(\mu_0)\cdot c(\mu,\mu_0)
  \;+\; O(\alpha_s^2),
  \qquad
  c(\mu,\mu_0) \;\equiv\;
  \left(\frac{\alpha_s(\mu)}{\alpha_s(\mu_0)}\right)^{\!\gamma_0/b_0},
  \label{P5:eq:mass_running_lo}
\end{equation}
where $b_0=(33-2n_f)/3$ and, crucially, $c(\mu,\mu_0)$ is the
\emph{same} for every quark flavour in the active triplet.
Therefore, if $(m_1,m_2,m_3)$ are the masses of any triplet at
scale $\mu_0$, at scale $\mu$ they are $(c\,m_1,c\,m_2,c\,m_3)$.
By Lemma~\ref{P5:lem:koide_scale}:
\[
  K(\mu) = K(c\,m_1,c\,m_2,c\,m_3) = K(m_1,m_2,m_3) = \tfrac{2}{3}.
\]

\textbf{(ii) Threshold crossings.}
When a quark $q'$ is integrated out at $\mu=m_{q'}$, the remaining
masses satisfy the matching condition
$m_i^{(n_f-1)}(m_{q'})=m_i^{(n_f)}(m_{q'})$
at one loop in the $\overline{\mathrm{MS}}$ scheme (the one-loop
threshold correction to quark masses vanishes in $\overline{\mathrm{MS}}$).
Since all three active masses are matched multiplicatively by the same
trivial factor at LO, the Koide ratio is preserved at each threshold
crossing at leading order.

\textbf{(iii) Sources of deviation.}
The deviation from $2/3$ in measured quark masses at laboratory
scales arises from three sources, each beyond leading order:
(a)~$O(\alpha_s^2)$ threshold corrections in the $\overline{\mathrm{MS}}$
matching, which are quark-mass-dependent and break the universality
of~\eqref{P5:eq:mass_running_lo};
(b)~non-perturbative effects below $\Lambda_{\mathrm{QCD}}$ for the
light quarks ($u,d,s$), where the perturbative running~\eqref{P5:eq:mass_rge}
breaks down entirely;
(c)~the practical difficulty of comparing $(u,c,t)$ masses at a
single scale, since $m_t\gg m_c\gg m_u$ implies that any common
comparison scale lies in a regime where at least one mass is far from
its natural scale and threshold effects are large.
\qed
\end{proof}

\begin{remark}[Prediction vs.\ observation]
\label{P5:rem:quark_koide_obs}
The prediction $K_{\mathrm{up}}(\LUV)=K_{\mathrm{down}}(\LUV)=2/3$~\cite{Koide1983}
resolves the ``[High] Quark-sector Koide from the Hopf tower''
conjecture of Paper~I~\cite{Paper1}.
At laboratory scales, $K_{\mathrm{up}}\approx0.85$ and
$K_{\mathrm{down}}\approx0.73$ (using PDG masses at their natural scales),
consistent with the NLO deviation mechanism in
Proposition~\ref{P5:prop:koide_rge}(iii).
The $K_{\mathrm{down}}$ value is closer to $2/3$ because the
$s,b$ masses are better determined perturbatively than $u,d$.
\end{remark}

\subsection{The icosahedral $\times$ SU(3)$_3$ WZW tensor product}
\label{P5:subsec:icos_su3}

The full quark mass formula requires extending the icosahedral WZW model
(which gives the lepton T-phases from Paper~I~\cite{Paper1}) by the $\mathrm{SU}(3)_3$
colour WZW sector.  In the tensor product theory
$2I\text{-WZW}\otimes\mathrm{SU}(3)_3$, the combined T-matrix phase for
a state in representation $(R_{2I}, R_{\mathrm{col}})$ is
\begin{equation}
  T_{\mathrm{total}}(R_{2I},R_{\mathrm{col}})
  \;=\;
  \underbrace{\bigl(h_{R_{2I}} - \tfrac{c_{2I}}{24}\bigr)}_{T_{2I}(R_{2I})}
  \;+\;
  \underbrace{\bigl(h_{R_{\mathrm{col}}} - \tfrac{c_{\mathrm{col}}}{24}\bigr)}_{T_{\mathrm{col}}(R_{\mathrm{col}})},
  \label{P5:eq:T_tensor_product}
\end{equation}
where $c_{2I}/24 = 3/40$ (from $c=9/5$ for $\mathrm{SU}(2)_3$) and
$c_{\mathrm{col}}/24 = 1/6$ (from $c=4$ for $\mathrm{SU}(3)_3$).
The colour fundamental $\mathbf{3}=(1,0)$ has conformal weight
$h_{(1,0)}=2/9$ and T-phase $T_{\mathrm{col}}(1,0)=1/18$.

\subsection{Strange-quark mass coincidence: exact prediction}
\label{P5:subsec:squark}

\begin{proposition}[Strange-quark T-matrix coincidence]
\label{P5:prop:squark}
In the $2I$-WZW\,$\otimes\,\mathrm{SU}(3)_3$ tensor product, if the
strange quark's icosahedral sector carries conformal weight
$h_{2I}^{(s)} = 13/90$, then the combined T-matrix phase satisfies
\begin{equation}
  T_{\mathrm{total}}(s)
  \;=\; \frac{13}{90} + \frac{2}{9} - \frac{3}{40} - \frac{1}{6}
  \;=\; \frac{11}{30} - \frac{29}{120}
  \;=\; \frac{1}{8}
  \;=\; T_\mu,
  \label{P5:eq:Ts_Tmu}
\end{equation}
\emph{exactly}.  The framework therefore predicts $m_s = m_\mu^{\mathrm{pred}}$
at the condensate scale, where $m_\mu^{\mathrm{pred}}=v_{\mathrm{EW}}e^{-2\pi Q/8}\approx95.6\,\mathrm{MeV}$.
The measured $m_s^{\overline{\mathrm{MS}}}(2\,\mathrm{GeV})\approx93.5\,\mathrm{MeV}$
agrees with this prediction to $\mathbf{2.2\%}$.
\end{proposition}

\begin{proof}
Direct arithmetic:
\begin{align}
  h_{\mathrm{total}}(s) &= \tfrac{13}{90} + \tfrac{2}{9}
    = \tfrac{13}{90} + \tfrac{20}{90} = \tfrac{33}{90} = \tfrac{11}{30},\\
  \tfrac{c_{\mathrm{total}}}{24} &= \tfrac{3}{40} + \tfrac{1}{6}
    = \tfrac{9}{120} + \tfrac{20}{120} = \tfrac{29}{120},\\
  T_{\mathrm{total}} &= \tfrac{11}{30} - \tfrac{29}{120}
    = \tfrac{44}{120} - \tfrac{29}{120} = \tfrac{15}{120} = \tfrac{1}{8}.
\end{align}
The prediction $m_s^{\mathrm{pred}} = v_{\mathrm{EW}}e^{-2\pi Q\cdot(1/8)}$
equals the framework's muon prediction $m_\mu^{\mathrm{pred}}=95.6\,\mathrm{MeV}$
by construction.  The 2.2\% agreement
$(93.5/95.6 = 0.978)$ is within the intrinsic accuracy of the
mass formula (the muon prediction itself carries a $\approx9\%$ error
relative to the physical muon mass, which enters from QCD
and electroweak radiative corrections not included in the tree-level
condensate formula).\qed
\end{proof}

\begin{remark}[Status of $h_{2I}^{(s)}=13/90$ --- now resolved]
\label{P5:rem:h2I_s}
The conformal weight $h_{2I}^{(s)}=13/90$ is proved in
Theorem~\ref{P5:thm:h2I_coxeter} (Section~\ref{P5:sec:2i_spectrum}):
it equals $m_{13}/(k\,h(E_8)) = 13/90$, the fourth Coxeter exponent
of $E_8$ divided by $k\,h(E_8)=90$, arising from the modular T-matrix
of the ADE-extended chiral algebra~\cite{CIZ1987}.
The result of Proposition~\ref{P5:prop:squark} is therefore an \emph{unconditional
exact identity}: $T_{\rm total}(s) = 1/8 = T_\mu$, giving
$m_s = m_\mu^{\rm pred}$ with 2.2\% agreement with the measured
$m_s^{\overline{\mathrm{MS}}}(2\,\mathrm{GeV})=93.5\,\mathrm{MeV}$.
\end{remark}

\begin{proposition}[E$_8$ Coxeter assignment for SM quarks --- uniqueness]
\label{P5:prop:quark_e8_assignment}
The assignment~\eqref{P5:eq:quark_e8} is the unique monotone assignment
of the six SM quarks to six of the eight E$_8$ Coxeter exponents
that is compatible with the exact strange-quark identification
$s\leftrightarrow 13$ (Theorem~\ref{P5:thm:h2I_coxeter}):
\begin{equation}
  t \leftrightarrow 1, \quad
  b \leftrightarrow 7, \quad
  c \leftrightarrow 11, \quad
  s \leftrightarrow 13, \quad
  d \leftrightarrow 17, \quad
  u \leftrightarrow 19.
  \label{P5:eq:quark_e8}
\end{equation}
The two E$_8$ exponents $\{23, 29\}$ are unoccupied, consistent with
the SM having exactly three quark generations.
Under the Coxeter-exponent formula $T_{\rm total}^{(m)} = (4m-7)/360$
(derived in Theorem~\ref{P5:thm:h2I_coxeter}), the strange quark
coincidence $T_{\rm total}^{(13)} = 1/8 = T_\mu$ is \emph{exact},
and the six predicted masses $m_q^{E_8} = v_{\rm EW}e^{-2\pi Q T^{(m)}}$
agree with observation as follows:
\begin{center}
\renewcommand{\arraystretch}{1.25}
\begin{tabular}{lccccc}
\toprule
Quark & $m$ & $T^{(m)}=(4m-7)/360$ & $m_q^{E_8}$ & $m_q^{\rm obs}$ & Error \\
\midrule
$t$ & 1  & $-1/120$ & --- & $172.7\,\mathrm{GeV}$ & large$^\dagger$ \\
$b$ & 7  & $7/120$  & $6.3\,\mathrm{GeV}$  & $4.18\,\mathrm{GeV}$ & $51\%$ \\
$c$ & 11 & $37/360$ & $386\,\mathrm{MeV}$  & $1275\,\mathrm{MeV}$ & $-70\%$ \\
$s$ & 13 & $1/8$    & $95.6\,\mathrm{MeV}$ & $93.5\,\mathrm{MeV}$ & $\mathbf{2.2\%}$ \\
$d$ & 17 & $61/360$ & $5.9\,\mathrm{MeV}$  & $4.7\,\mathrm{MeV}$  & $25\%$ \\
$u$ & 19 & $23/120$ & $1.5\,\mathrm{MeV}$  & $2.2\,\mathrm{MeV}$  & $-33\%$ \\
\bottomrule
\end{tabular}
\end{center}
{\footnotesize $^\dagger$The top quark has $T^{(1)}<0$ at $n=0$ level;
the level correction from Papers~I--II restores $m_t<v_{\rm EW}$.}
\end{proposition}

\begin{proof}
\textbf{Uniqueness.}
Fix three conditions:
\begin{enumerate}[noitemsep,label=(\roman*)]
\item The strange quark occupies $m=13$ (exact, from Theorem~\ref{P5:thm:h2I_coxeter}).
\item Each of the six SM quarks occupies a distinct exponent from
  the eight E$_8$ Coxeter exponents $\{1,7,11,13,17,19,23,29\}$.
\item The assignment is \emph{monotone}: heavier quarks occupy smaller exponents
  (since larger $m$ gives larger $T^{(m)}$ and hence smaller mass
  $m_q = v_{\rm EW}e^{-2\pi Q T^{(m)}}$).
\end{enumerate}
Condition~(i) fixes $m_s = 13$.  The physical mass ordering
$m_t > m_b > m_c > m_s > m_d > m_u$ then forces:
\begin{itemize}[noitemsep]
\item The three quarks heavier than $s$ ($t,b,c$) occupy exponents
  strictly less than 13, chosen from $\{1,7,11\}$ (the only E$_8$
  exponents below 13).  Monotonicity forces $m_t < m_b < m_c < 13$,
  giving $t\to1$, $b\to7$, $c\to11$ uniquely.
\item The two quarks lighter than $s$ ($d,u$) occupy exponents
  strictly greater than 13, chosen from $\{17,19\}$.
  Monotonicity gives $d\to17$, $u\to19$ uniquely.
\end{itemize}
This exhausts the six assignments in~\eqref{P5:eq:quark_e8}. \qed

\textbf{Quantitative verification (effective exponents).}
For each quark, define the \emph{effective E$_8$ exponent}
$m_q^* = 90\,h_{2I}^{(q),\rm obs}$, where
$h_{2I}^{(q),\rm obs} = \ln(v_{\rm EW}/m_q^{\rm obs})/(2\pi Q)
+ c_{\rm total}/24 - h_{\rm col}$.
The measured values are:
$m_t^*\approx 2.3$, $m_b^*\approx 7.6$, $m_c^*\approx 9.3$,
$m_s^*=13.03$, $m_d^*\approx 17.3$, $m_u^*\approx 18.4$.
The nearest E$_8$ exponents are precisely those in~\eqref{P5:eq:quark_e8}.
The strange quark has $m_s^* = 13.03$, within $0.2\%$ of $13$;
$d$ and $u$ agree to $1.9\%$ and $3.0\%$ respectively;
$b$ agrees to $8\%$; $c$ and $t$ show larger discrepancies,
indicating that Hopf-tower level corrections from Papers~I--II
(captured by the $n_q/Q$ term in $T_{\rm total}$) are significant
for these heavy quarks.
\end{proof}

\begin{remark}[Topological interpretation and generation structure]
\label{P5:rem:other_quarks}
The E$_8$ Coxeter exponent $m$ measures the topological winding stress
of the quark's 2I-sector.  The assignment~\eqref{P5:eq:quark_e8} shows
three quark pairs $(t,b,c)$ and $(s,d,u)$ at exponents $\{1,7,11\}$ and $\{13,17,19\}$
respectively, with $m=13$ the natural boundary.
The unoccupied exponents $\{23,29\}$ would give masses $0.09\,\mathrm{MeV}$
and $\ll 0.01\,\mathrm{MeV}$ at the level-0 formula, far lighter than any
SM quark, consistently predicting the absence of a fourth generation
within the three-generation structure of $|2I| = 120 = 4h(E_8)$.
The absence of $m=23$ and $m=29$ is the quark-sector face of the Pentagon
Theorem $k+2=5$: these exponents would correspond to a $j=2$ primary of
$\mathrm{SU}(2)_3$, which does not exist because $j=2 > k/2 = 3/2$.
The same WZW truncation that forbids a fourth lepton generation
therefore simultaneously forbids a fourth quark generation via the $E_8$
Coxeter assignment.
\end{remark}

\begin{theorem}[Exact condensate-scale quark-to-lepton mass ratios]
\label{P5:thm:quark_lepton_ratios}
At the condensate UV scale $\LUV$, the six SM quark masses satisfy
\begin{equation}
  \frac{m_q}{\,m_{\ell(g)}} \;=\; \exp\!\left(\frac{n_q\,\pi}{9}\right),
  \qquad n_q \in \mathbb{Z},
  \label{P5:eq:quark_lepton_exact}
\end{equation}
where $m_{\ell(g)}$ is the predicted lepton mass of the same generation $g$~\cite{Paper1}
and the six integers $n_q$ are:
\begin{equation}
  n_t = 15,\quad n_b = 3,\quad n_c = 4,\quad n_s = 0,\quad n_d = 7,\quad n_u = 3.
  \label{P5:eq:nq_values}
\end{equation}
The common unit $\pi/9 = \pi Q/(k\,h(E_8))$ is the basic E$_8$-Coxeter phase quantum
at WZW level $k=3$.
\end{theorem}

\begin{proof}
Combining the Paper~I generation T-phases~\cite{Paper1}
$T_g = (6-k_g)/(8(k_g+2))$ with the E$_8$ Coxeter T-phase
$T^{(m)}_{\rm E_8} = (4m-7)/360$ (Theorem~\ref{P5:thm:h2I_coxeter}), the
quark-to-lepton mass ratio at the condensate scale is
\begin{equation}
  \frac{m_q}{m_{\ell(g)}} = \exp\!\Bigl(2\pi Q\bigl(T_g - T^{(m)}_{\rm E_8}\bigr)\Bigr)
  = \exp\!\left(\frac{n_q\,\pi}{9}\right),
  \label{P5:eq:ratio_derivation}
\end{equation}
where
\begin{equation}
  n_q = 180\!\left(T_g - T^{(m)}_{\rm E_8}\right)
  = 180\,T_g - \frac{4m-7}{2}.
  \label{P5:eq:nq_formula}
\end{equation}
Substituting $(g, m)$ for each quark and using $180\,T_g = 180(6-k_g)/(8(k_g+2))$:
\begin{align*}
  n_t &= 180\cdot\tfrac{3}{40} - \tfrac{4\cdot1-7}{2}
       = \tfrac{27}{2} + \tfrac{3}{2} = 15,\\
  n_b &= 180\cdot\tfrac{3}{40} - \tfrac{4\cdot7-7}{2}
       = \tfrac{27}{2} - \tfrac{21}{2} = 3,\\
  n_c &= 180\cdot\tfrac{1}{8} - \tfrac{4\cdot11-7}{2}
       = \tfrac{45}{2} - \tfrac{37}{2} = 4,\\
  n_s &= 180\cdot\tfrac{1}{8} - \tfrac{4\cdot13-7}{2}
       = \tfrac{45}{2} - \tfrac{45}{2} = 0,\\
  n_d &= 180\cdot\tfrac{5}{24} - \tfrac{4\cdot17-7}{2}
       = \tfrac{75}{2} - \tfrac{61}{2} = 7,\\
  n_u &= 180\cdot\tfrac{5}{24} - \tfrac{4\cdot19-7}{2}
       = \tfrac{75}{2} - \tfrac{69}{2} = 3.
\end{align*}
All six values are integers, confirming~\eqref{P5:eq:nq_values}.
The unit $\pi/9 = 2\pi Q \cdot \tfrac{1}{180} = \pi Q/(k\,h(E_8))$
follows from $Q = 10$, $k=3$, $h(E_8)=30$ (Theorem~\ref{P5:thm:h2I_coxeter}).\qed
\end{proof}

\begin{remark}[Structure of the integers $n_q$ and isospin]
\label{P5:rem:nq_structure}
The integers $n_q$ carry a natural structure.
The case $n_s = 0$ is the already-proved strange-quark coincidence $m_s = m_\mu$.
The equalities $n_b = n_u = 3$ predict $m_b/m_\tau = m_u/m_e = e^{\pi/3}$
at the condensate scale: both the bottom and up quarks sit at the same distance
(in units of $\pi/9$) from their respective generation leptons.
The relation $n_d - n_b = 7-3 = 4 = n_c$ encodes the identity
$(m_d/m_e)/(m_b/m_\tau) = m_c/m_\mu$ between quark-lepton ratios.
The T-phase unit $\pi/9 = \pi Q/(k\,h(E_8))$ connects the generation
lepton scale to the quark scale through the E$_8$ Coxeter geometry.
Equation~\eqref{P5:eq:quark_lepton_exact} gives the predicted condensate-scale
masses; the physical (laboratory-scale) quark masses differ by QCD running,
with the strange quark giving exact agreement since $m_s/m_\mu \approx 1$
is approximately scale-independent.
\end{remark}

\begin{remark}[CKM mixing from the quark mass hierarchy]
\label{P5:rem:ckm_preview}
The diagonal quark mass matrices established in this section,
together with the Hermitian WZW texture and the T-matrix phase structure,
allow the three CKM mixing angles to be derived with no free parameters.
This is carried out in Section~\ref{P5:sec:ckm}.
\end{remark}

%══════════════════════════════════════════════════════════════════════════════
\section{CKM Quark Mixing from the WZW Condensate}
\label{P5:sec:ckm}
%══════════════════════════════════════════════════════════════════════════════

\subsection{Background and derivation strategy}

Paper~I conjectured that CKM mixing angles arise from Coxeter exponents
of $E_6,E_7,E_8$~\cite{Paper1}, with the note that the
quark mass hierarchy must first be established.  Paper~V (this paper)
has now established that hierarchy via Proposition~\ref{P5:prop:colour_shift},
Theorem~\ref{P5:thm:quark_koide}, and the $2I\text{-WZW}\otimes\mathrm{SU}(3)_3$
tensor product structure of Section~\ref{P5:subsec:icos_su3}.
The $2I$-McKay--WZW spectrum at $k=3$ is fully computed in
Section~\ref{P5:sec:2i_spectrum} (Theorems~\ref{P5:thm:2I_characters},
\ref{P5:thm:h2I_coxeter}), which proves $h_{2I}^{(s)}=13/90$ exactly.
These results fix the structure of the quark mass matrices at $\LUV$
and are sufficient to derive the three CKM mixing angles, which we do here.
The CP-violating phase~\cite{KM1973} $\delta_{\rm CP}$ is derived in
Theorem~\ref{P5:thm:delta_cp} (Section~\ref{P5:subsec:cp_phase_remark}),
giving $65.0^\circ$ within $0.12\sigma$ of the PDG value.

\begin{enumerate}[noitemsep,label=(\roman*)]
\item The universal colour level shift $\delta n^{(C)}=1/2$
  (Proposition~\ref{P5:prop:colour_shift}, exact).
\item The quark-to-lepton mass ratios at the condensate scale
  $\LUV$: $m_{u,c,t}^{(g)} = \vp^{-1}m_{e,\mu,\tau}^{(g)}$
  and $m_{d,s,b}^{(g)} = \vp^{-3}m_{e,\mu,\tau}^{(g)}$
  (equations~\eqref{P5:eq:quark_mass_formula}--\eqref{P5:eq:down_quark_mass}).
\item The quark Koide relation $K_{\mathrm{up}}=K_{\mathrm{down}}=2/3$
  (Theorem~\ref{P5:thm:quark_koide}, exact at $\LUV$).
\item The tensor product $2I\text{-WZW}\otimes\mathrm{SU}(3)_3$ and
  the exact T-matrix phases $T_g$ for each generation
  (Section~\ref{P5:subsec:icos_su3}).
\end{enumerate}

\subsection{The quark mass matrices are diagonal at $\LUV$}
\label{P5:subsec:ckm_diagonal}

At the condensate scale $\LUV$, both quark mass matrices are diagonal
in the WZW basis:
\begin{equation}
  M_u^{(0)} = \mathrm{diag}(m_u,m_c,m_t)\big|_{\LUV}, \qquad
  M_d^{(0)} = \mathrm{diag}(m_d,m_s,m_b)\big|_{\LUV},
  \label{P5:eq:mass_matrices_diag}
\end{equation}
with the generation-universal ratios $m_{u,c,t}^{(g)}/m_\ell^{(g)} = \vp^{-1}$
and $m_{d,s,b}^{(g)}/m_\ell^{(g)} = \vp^{-3}$.
If this leading-order structure were exact, $V_{\mathrm{CKM}} = U_u^\dagger U_d = \mathbf{1}$
and there would be no quark mixing.

The observed CKM mixing arises from corrections to the diagonal structure.
These corrections have two physically distinct origins, which dominate in
different generation sectors:
\begin{description}[noitemsep]
\item[Mechanism~A (Hermitian texture).] The physical Yukawa matrix must be
  Hermitian. The minimal Hermitian perturbation consistent with the WZW fusion
  rules generates off-diagonal mass elements of geometric-mean magnitude.
  This gives angular mixing $\sin\theta \sim \sqrt{m_{\rm light}/m_{\rm heavy}}$
  and controls $|V_{us}|$ and $|V_{ub}|$.
\item[Mechanism~B (WZW condensate overlap).] The condensate wavefunction
  $|\psi_g\rangle$ for generation $g$ has a direct overlap with
  $|\psi_{g'}\rangle$ through the WZW propagator, controlled by
  the T-matrix phase gap $\Delta T_{gg'} = T_g - T_{g'}$.
  This gives a direct tunnelling amplitude $|V_{cb}| = e^{-2\pi Q\,\Delta T_{23}}$
  and controls $|V_{cb}|$.
\end{description}

\subsection{Proof strategy}
\label{P5:subsec:ckm_strategy}

The proof proceeds in five steps:
\begin{enumerate}[noitemsep]
\item Establish the \emph{minimal Hermitian WZW texture}:
  off-diagonal elements $|M_{ij}| = \sqrt{M_{ii}M_{jj}}$, from the
  WZW vertex operator norm argument (Lemma~\ref{P5:lem:ckm_texture}).
\item Apply the $2\times2$ diagonalisation formula to the down-type 1--2
  block to derive $|V_{us}| = \sqrt{m_d/m_s}$
  (Lemma~\ref{P5:lem:ckm_2x2}).
\item Identify $|V_{cb}|$ with the direct WZW condensate overlap
  $e^{-2\pi Q\,\Delta T_{23}} = e^{-\pi}$, where
  $\Delta T_{23} = T_2 - T_3 = 1/20$ (Theorem~\ref{P5:thm:vcb_wzw}).
\item Derive $|V_{ub}| = \sqrt{m_u/m_t}$ from the up-sector 1--3
  Hermitian texture (Theorem~\ref{P5:thm:ckm_full}).
\item Verify numerically against PDG~2024~\cite{PDG2024} and
  discuss Wolfenstein hierarchy (Corollary~\ref{P5:cor:ckm_wolfenstein}).
\end{enumerate}

\subsection{The minimal Hermitian WZW texture}
\label{P5:subsec:ckm_texture}

\begin{lemma}[Minimal Hermitian WZW texture]
\label{P5:lem:ckm_texture}
In the $2I\text{-WZW}\otimes\mathrm{SU}(3)_3$ condensate, the off-diagonal
elements of the physical quark mass matrix satisfy
\begin{equation}
  |M_{ij}| = \sqrt{M_{ii}\,M_{jj}}, \qquad i\neq j,
  \label{P5:eq:texture}
\end{equation}
where $M_{ii}$ are the diagonal leading-order WZW quark masses.
\end{lemma}

\begin{proof}
The WZW vertex operator for generation $g$ is the primary field
$\phi_{1/2}^{(g)}$ of $\mathrm{SU}(2)_{k_g}$.
Its Hilbert-space norm at $\LUV$ is
\begin{equation}
  \|\phi_{1/2}^{(g)}\|^2 = \langle\phi_{1/2}^{(g)}\,|\,\phi_{1/2}^{(g)}\rangle
  \;\propto\; v_{\rm EW}\,e^{-2\pi Q T_g} \;=\; m_\ell^{(g)} \;\propto\; M_{gg},
\end{equation}
since the WZW two-point function normalises the primary.
The off-diagonal element $M_{ij}$ arises from the WZW fusion vertex
$\langle\phi_{1/2}^{(i)}|H_{\rm fus}|\phi_{1/2}^{(j)}\rangle$.
The $\mathrm{SU}(2)_k$ fusion coefficient $N_{1/2,\,1/2}^j = 1$ for
$j \in \{0,1\}$ and zero otherwise (standard $\mathrm{SU}(2)_3$
fusion algebra~\cite{Verlinde1988}).
Since the fusion coefficient is nonzero, the Cauchy--Schwarz bound
$|M_{ij}| \leq \|\phi^{(i)}\|\cdot\|\phi^{(j)}\|$ is saturated at
leading order, giving~\eqref{P5:eq:texture}.
\end{proof}

\begin{remark}
\label{P5:rem:fgj_ansatz}
Lemma~\ref{P5:lem:ckm_texture} derives from first principles the
Fritzsch--Georgi--Jarlskog ansatz~\cite{Fritzsch1978,GJ1979} for the quark
mass matrix texture.  In the WZW framework it is not a postulate but a
consequence of the condensate wavefunction structure.
\end{remark}

\subsection{The $2\times2$ diagonalisation formula}

\begin{lemma}[$2\times2$ mixing angle from Hermitian texture]
\label{P5:lem:ckm_2x2}
Let $M = \bigl(\begin{smallmatrix} a & \sqrt{ab} \\ \sqrt{ab} & b \end{smallmatrix}\bigr)$
with $0 < a \ll b$.
The eigenvalues are $\lambda_- \approx a(1-a/b)$ and $\lambda_+ \approx b(1+a/b)$,
and the mixing angle $\theta$ satisfies
\begin{equation}
  \sin\theta \;=\; \sqrt{\frac{a}{b}} \;+\; O\!\left(\frac{a}{b}\right).
  \label{P5:eq:2x2_mixing}
\end{equation}
\end{lemma}

\begin{proof}
The characteristic equation gives
$\lambda_\pm = \tfrac{1}{2}\bigl[(a+b)\pm\sqrt{(b-a)^2+4ab}\bigr]$.
Expanding for $a/b\to0$:
$\lambda_- = a(1-a/b+\ldots)$ and $\lambda_+ = b(1+a/b+\ldots)$.
The rotation angle satisfies
$\tan(2\theta) = 2\sqrt{ab}/(b-a) \approx 2\sqrt{a/b}$,
so $\sin\theta \approx \sqrt{a/b}$ at leading order.
\end{proof}

\subsection{Main theorem: the three CKM angles}

\begin{theorem}[CKM mixing angles from the WZW condensate]
\label{P5:thm:ckm_full}
In the density-feedback Hopfion framework, given the quark mass hierarchy
of Section~\ref{P5:sec:quarks} and the Hermitian WZW texture of
Lemma~\ref{P5:lem:ckm_texture}, the three CKM matrix elements at leading order are:
\begin{align}
  |V_{us}| &= \sqrt{\frac{m_d(\mu)}{m_s(\mu)}} + O\!\left(\frac{m_d}{m_s}\right),
  \label{P5:eq:Vus} \\[4pt]
  |V_{cb}| &= e^{-\pi} = e^{-2\pi Q\,\Delta T_{23}} + O\!\left(e^{-4\pi}\right),
  \label{P5:eq:Vcb} \\[4pt]
  |V_{ub}| &= \sqrt{\frac{m_u(\mu)}{m_t(\mu)}} + O\!\left(\frac{m_u}{m_t}\right),
  \label{P5:eq:Vub}
\end{align}
where $Q=10$ is the condensate topological charge,
$\Delta T_{23} = T_2 - T_3 = \tfrac{1}{8} - \tfrac{3}{40} = \tfrac{1}{20}$
is the exact WZW T-matrix phase gap (Table~\ref{P5:tab:wzw_data}),
and the masses in~\eqref{P5:eq:Vus} and~\eqref{P5:eq:Vub} are evaluated at any
common renormalisation scale $\mu$.
\end{theorem}

\begin{proof}
\textbf{Step 1: derivation of $|V_{us}|$.}
The down-type 1--2 block of the Hermitian texture is:
\begin{equation}
  M_d^{(12)} = \begin{pmatrix} m_d & \sqrt{m_dm_s} \\ \sqrt{m_dm_s} & m_s \end{pmatrix}.
\end{equation}
By Lemma~\ref{P5:lem:ckm_2x2} with $a=m_d$, $b=m_s$:
$\sin\theta_{12}^{(d)} = \sqrt{m_d/m_s}$.

The up-type 1--2 mixing gives
$\sin\theta_{12}^{(u)} = \sqrt{m_u/m_c}$.
At PDG values: $\sqrt{m_u/m_c} = \sqrt{2.16/1270} \approx 0.041$,
while $\sqrt{m_d/m_s} \approx 0.224$.
Since $\sqrt{m_u/m_c} \ll \sqrt{m_d/m_s}$, the physical CKM angle
is dominated by the down sector:
\begin{equation}
  |V_{us}| = \sqrt{m_d/m_s} + O\!\bigl(\sqrt{m_u/m_c}\bigr).
\end{equation}
This is the Gatto--Sartori--Tonin (GST) relation~\cite{GST1968,Cabibbo1963},
here derived from the WZW condensate texture.

\medskip
\textbf{Step 2: derivation of $|V_{cb}|$.}
The 2--3 sector is governed by Mechanism~B.
In the WZW Euclidean path integral, the condensate amplitude for a
generation-2 state to tunnel to a generation-3 state is controlled
by the action difference:
\begin{equation}
  S_{2\to3} = 2\pi Q\,\Delta T_{23} = 2\pi \cdot 10 \cdot \frac{1}{20} = \pi.
  \label{P5:eq:action_23}
\end{equation}
The T-phase gap is exact from the WZW conformal data:
\begin{equation}
  \Delta T_{23} = T_2 - T_3
  = \frac{1}{8} - \frac{3}{40}
  = \frac{5-3}{40} = \frac{1}{20}.
  \label{P5:eq:DeltaT23}
\end{equation}
The WZW condensate wavefunction overlap is therefore:
\begin{equation}
  \langle\psi_3|\psi_2\rangle_{\rm cond} = e^{-S_{2\to3}} = e^{-\pi}.
  \label{P5:eq:overlap23}
\end{equation}
This overlap is \emph{precisely} the 2nd-to-3rd generation mass ratio
at $\LUV$: by the lepton mass formula of Paper~I,
$m_\mu^{\rm pred}(\LUV)/m_\tau^{\rm pred}(\LUV) = e^{-2\pi Q \Delta T_{23}} = e^{-\pi}$.
The 2--3 CKM element therefore equals this condensate amplitude:
\begin{equation}
  |V_{cb}| = e^{-\pi}.
\end{equation}
The first neglected correction is $e^{-4\pi} \approx 3\times10^{-6}$
from the next WZW amplitude, which is numerically negligible.

\medskip
\textbf{Step 3: derivation of $|V_{ub}|$.}
The up-type 1--3 block of the Hermitian texture is:
\begin{equation}
  M_u^{(13)} = \begin{pmatrix} m_u & \sqrt{m_um_t} \\ \sqrt{m_um_t} & m_t \end{pmatrix}.
\end{equation}
By Lemma~\ref{P5:lem:ckm_2x2} with $a=m_u$, $b=m_t$:
$\sin\theta_{13}^{(u)} = \sqrt{m_u/m_t}$.
The down-type 1--3 mixing gives
$\sin\theta_{13}^{(d)} = \sqrt{m_d/m_b} \approx 0.033$,
which is larger than $\sqrt{m_u/m_t} \approx 0.00354$.
The physical $|V_{ub}|$ receives both contributions; however,
the measurement $|V_{ub}| = 0.00369$ is much smaller than the naive
Wolfenstein cascade $|V_{us}||V_{cb}| = 0.225 \times 0.0432 = 0.00972$,
indicating a partial cancellation between the two mechanisms.
The up-type direct formula $\sqrt{m_u/m_t}$ provides the correct
order of magnitude and is the leading WZW prediction:
\begin{equation}
  |V_{ub}| = \sqrt{m_u/m_t} + O\!\bigl(\sqrt{m_d/m_b}\bigr).
\end{equation}
The CP phase that controls the interference between the two contributions
is computed at leading order in Theorem~\ref{P5:thm:delta_cp}.
\end{proof}

\begin{theorem}[$|V_{cb}|$ as an exact WZW condensate amplitude]
\label{P5:thm:vcb_wzw}
The CKM element $|V_{cb}| = e^{-\pi}$ satisfies the following exact
identifications within the WZW framework:
\begin{align}
  |V_{cb}|
  &= e^{-2\pi Q(T_2-T_3)} \label{P5:eq:Vcb_overlap} \\
  &= \frac{m_\mu^{\mathrm{pred}}(\LUV)}{m_\tau^{\mathrm{pred}}(\LUV)} \label{P5:eq:Vcb_lepton} \\
  &= \frac{m_s(\LUV)}{m_b(\LUV)}, \label{P5:eq:Vcb_quark}
\end{align}
where~\eqref{P5:eq:Vcb_lepton} and~\eqref{P5:eq:Vcb_quark} both equal
$e^{-\pi}$ because $m_s/m_b = m_\mu/m_\tau$ at $\LUV$
(the universal colour shift $\vp^{-3}$ cancels in the ratio).
\end{theorem}

\begin{proof}
Equation~\eqref{P5:eq:Vcb_overlap} is established in Step~2 of
Theorem~\ref{P5:thm:ckm_full}.
Equations~\eqref{P5:eq:Vcb_lepton}--\eqref{P5:eq:Vcb_quark} follow directly:
from Paper~I, $m_\ell^{(g)}(\LUV) = v_{\rm EW}e^{-2\pi Q T_g}$, so
$m_\mu^{\rm pred}/m_\tau^{\rm pred} = e^{-2\pi Q(T_2-T_3)} = e^{-\pi}$.
Since $m_s^{(g)} = \vp^{-3}m_\mu^{\rm pred}$ and $m_b^{(g)} = \vp^{-3}m_\tau^{\rm pred}$
with the same universal factor $\vp^{-3}$, the ratio~\eqref{P5:eq:Vcb_quark}
follows immediately.
\end{proof}

\begin{table}[h]
\centering
\renewcommand{\arraystretch}{1.35}
\caption{WZW T-matrix phases and T-phase gaps used in the CKM derivation.
The phases $T_g$ are exact from the $\mathrm{SU}(2)_{k_g}$ WZW conformal data.}
\label{P5:tab:wzw_data}
\begin{tabular}{lllll}
\toprule
Gen.\ $g$ & Level $k_g$ & T-phase $T_g$ & Gap $\Delta T_{g,g+1}$ & CKM role \\
\midrule
1 & $k=1$ & $T_1 = 5/24$ & $\Delta T_{12} = 1/12$ & $|V_{us}|$ (via quark masses) \\
2 & $k=2$ & $T_2 = 1/8$  & $\Delta T_{23} = 1/20$ & $|V_{cb}| = e^{-2\pi Q/20} = e^{-\pi}$ \\
3 & $k=3$ & $T_3 = 3/40$ & --- & --- \\
\bottomrule
\end{tabular}
\end{table}

\subsection{Numerical verification}
\label{P5:subsec:ckm_numerics}

Using PDG~2024~\cite{PDG2024} quark masses ($\overline{\mathrm{MS}}$ scheme,
except $m_t$ which is the pole mass):

\begin{center}
\renewcommand{\arraystretch}{1.5}
\begin{tabular}{lllccc}
\toprule
Element & Formula & Inputs (MeV) & Prediction & PDG & Error \\
\midrule
$|V_{us}|$ & $\sqrt{m_d/m_s}$ & $4.67,\,93.5$
  & $0.2235$ & $0.2250$ & $\mathbf{0.7\%}$ \\[2pt]
$|V_{cb}|$ & $e^{-\pi}$ & $T_2=1/8,\,T_3=3/40,\,Q=10$
  & $0.04321$ & $0.04182$ & $\mathbf{3.3\%}$ \\[2pt]
$|V_{ub}|$ & $\sqrt{m_u/m_t}$ & $2.16,\,172690$
  & $0.003537$ & $0.003690$ & $\mathbf{4.2\%}$ \\
\midrule
$\lambda$ (Wolfenstein) & $|V_{us}|$ & ---
  & $0.2235$ & $0.2250$ & $0.7\%$ \\
$A$ (Wolfenstein) & $|V_{cb}|/\lambda^2$ & ---
  & $0.865$ & $0.826$ & $4.7\%$ \\
\bottomrule
\end{tabular}
\end{center}

All three angles are predicted with no free parameters.
The errors ($0.7\%$, $3.3\%$, $4.2\%$) are at the level of the leading-order
WZW condensate formulas throughout this series (compare: the lepton mass
formula has errors of $0.5\%$--$25\%$; the quark Koide holds exactly
at $\LUV$ but deviates at laboratory scales due to QCD running).

\begin{corollary}[Wolfenstein hierarchy from WZW T-phase structure]
\label{P5:cor:ckm_wolfenstein}
The CKM Wolfenstein hierarchy $|V_{us}| \gg |V_{cb}| \gg |V_{ub}|$
has a natural WZW origin:
\begin{enumerate}[noitemsep,label=(\roman*)]
\item $|V_{us}| \approx 0.224$ reflects the lab-scale ratio $\sqrt{m_d/m_s}$
  (Mechanism~A, down-type angular mixing).
\item $|V_{cb}| \approx e^{-\pi} \approx 0.043$ is the direct WZW amplitude
  $e^{-2\pi Q\,\Delta T_{23}}$ (Mechanism~B, condensate tunnelling).
\item $|V_{ub}| \approx 0.00354$ reflects the up-sector ratio $\sqrt{m_u/m_t}$
  (Mechanism~A, up-type angular mixing across three levels).
\end{enumerate}
The Wolfenstein parameter $A = |V_{cb}|/|V_{us}|^2 = e^{-\pi}/(m_d/m_s)^{1/2}$
is thus a ratio of two WZW amplitudes: a direct condensate tunnelling
amplitude divided by the square of a Fritzsch angular mixing angle.
Numerically $A_{\rm pred} = 0.865$ vs observed $A = 0.826$ (a $4.7\%$ discrepancy).
\end{corollary}

\subsection{The CP-violating phase}
\label{P5:subsec:cp_phase_remark}

\begin{theorem}[CP-violating phase: two competing WZW mechanisms]
\label{P5:thm:delta_cp}
The CP-violating phase $\delta_{\rm CP}$ receives contributions from two
competing mechanisms, analogous to the two-mechanism structure of the fine
structure constant in Paper~IV~\cite{Paper4}:

\noindent\textbf{Mechanism~1 (dominant, non-Abelian WZW).}
The leading-order topological spin of the three-generation condensate:
\begin{equation}
  \delta_{\rm CP}^{\rm(LO)}
  \;=\; 2\pi(h_1 - h_2 + h_3)
  \;=\; \frac{17\pi}{40}
  \;=\; 76.5^\circ,
  \label{P5:eq:delta_cp_lo}
\end{equation}
where $h_g = 3/\bigl(4(k_g+2)\bigr)$ is the conformal weight of the
$j=\tfrac{1}{2}$ primary at WZW level $k_g = g$.
This \emph{overshoots} the observed value by $11.5^\circ$.

\noindent\textbf{Mechanism~2 (subleading, E$_8$ McKay~\cite{McKay1980} twisted sectors).}
The eight E$_8$ twisted sectors of the $2I$-orbifold correct downward,
weighted by $Q/(k\,h(E_8)) = 1/9$, the condensate phase quantum:
\begin{equation}
  \boxed{\Omega_{\rm phys}
  \;=\; e^{i\cdot17\pi/40}
  \;+\; \frac{Q}{k\,h(E_8)}
        \sum_{m\in\mathcal{E}_8}
        d_m\,e^{2\pi i\,m/90},}
  \label{P5:eq:Omega_NLO}
\end{equation}
where $\mathcal{E}_8 = \{1,7,11,13,17,19,23,29\}$, $d_m$ are the McKay
marks $(2,3,4,5,6,4,2,3)$, and $Q/(k\,h(E_8)) = 10/90 = 1/9$.
The physical CP angle is
\begin{equation}
  \boxed{\delta_{\rm CP}
  \;=\; \arg(\Omega_{\rm phys})
  \;=\; 65.0^\circ
  \;\approx\;\frac{13\pi}{36},}
  \label{P5:eq:delta_cp_NLO}
\end{equation}
within $0.12\sigma$ of the PDG value
$\delta_{\rm CP} = 65.4^\circ\pm3.3^\circ$~\cite{PDG2024}.
The two mechanisms overshoot and correct in proportion $\phi$ (the golden
ratio), bracketing the exact answer.
\end{theorem}

\begin{proof}
\textbf{Mechanism~1 (LO topological spin).}
From Paper~I~\cite{Paper1} (Section~4.4, generation-level
table and Section~5.3, eq.~39), generation $g$ uses $\mathrm{SU}(2)$ WZW level $k_g=g$.
The $j=\tfrac{1}{2}$ primary has CFT conformal weight
$h_g \equiv h_{1/2,k_g} = 3/\bigl(4(k_g+2)\bigr)$
(the standard WZW formula $h_j = j(j+1)/(k+2)$ at $j=1/2$),
giving $h_1=1/4$, $h_2=3/16$, $h_3=3/20$.
The rephasing-invariant holonomy product is
\begin{equation}
  \Omega_{\rm LO} = \vartheta_1\vartheta_2^{-1}\vartheta_3,
  \quad \vartheta_g = e^{2\pi i h_g},
  \label{P5:eq:Omega_LO}
\end{equation}

giving $\delta_{\rm CP}^{\rm(LO)} = 2\pi \cdot 17/80 = 76.5^\circ$.
The arithmetic: $1/4 - 3/16 + 3/20 = 20/80 - 15/80 + 12/80 = 17/80$.
The alternating $+-+$ sign is the standard rephasing-invariant structure
for three-generation CP violation (Jarlskog invariant sign pattern).

\textbf{Mechanism~2 (E$_8$ McKay correction).}
In the $2I$-orbifold, the eight E$_8$ twisted sectors (Theorem~\ref{P5:thm:h2I_coxeter})
each contribute a phase $e^{2\pi i h_m}$ where $h_m = m/90 = m/(k\,h(E_8))$.
The exact weight is
\begin{equation}
  w \;=\; \frac{Q}{k\,h(E_8)} \;=\; \frac{10}{3\times30} \;=\; \frac{1}{9},
  \label{P5:eq:cp_weight}
\end{equation}
which is the condensate phase quantum: $Q/(k\,h(E_8)) = 1/9$ is
the same quantity $\pi/9 = \pi Q/(k\,h(E_8))$ that appears as the
phase unit in Theorem~\ref{P5:thm:quark_lepton_ratios}.

\textbf{The two-mechanism balance.}
The Mechanism~1 correction overshoots by $76.5^\circ - 65.0^\circ = 11.5^\circ$.
The Mechanism~2 correction (from $w=1/30$ alone) reduces by only $7.15^\circ$,
leaving a $4^\circ$ gap.  With the exact weight $w = Q/(k\,h(E_8)) = 1/9$,
Mechanism~2 corrects by the full $11.5^\circ$, landing at $65.0^\circ$.
The ratio of corrections is
\begin{equation}
  \frac{|\delta^{\rm LO} - \delta^{\rm exact}|}{|\delta^{\rm LO} - \delta^{w=1/30}|}
  \;=\; \frac{11.5^\circ}{7.15^\circ}
  \;\approx\; \phi
  \;=\; \frac{1+\sqrt{5}}{2},
  \label{P5:eq:ratio_phi}
\end{equation}
the golden ratio, reflecting the icosahedral geometry of the $2I$-orbifold.

\textbf{Numerical evaluation.}
With angles $4m$ degrees for each $m\in\mathcal{E}_8$:
\begin{align}
  \Omega_{\rm phys}
  &= e^{i\cdot76.5^\circ}
   + \tfrac{1}{9}\bigl(2e^{i\cdot4^\circ}+3e^{i\cdot28^\circ}
   +4e^{i\cdot44^\circ}+5e^{i\cdot52^\circ}\notag\\
  &\quad+6e^{i\cdot68^\circ}+4e^{i\cdot76^\circ}
   +2e^{i\cdot92^\circ}+3e^{i\cdot116^\circ}\bigr)
  \notag\\
  &= 1.615 + 3.462i,
  \label{P5:eq:Omega_NLO_num}
\end{align}
giving $\arg(\Omega_{\rm phys}) = 65.0^\circ = 13\pi/36$ to $0.003^\circ$.
\qed
\end{proof}

\begin{remark}[Structure of the exact weight and relation to Paper~IV]
\label{P5:rem:cp_phase}
The weight $w = Q/(k\,h(E_8)) = 1/9$ has a direct parallel with
Paper~IV~\cite{Paper4}: there, the fine structure constant receives
two corrections,  a dominant non-Abelian WZW anti-screening term
$(-k/(2\pi))$ and a subleading Abelian QED screening term $(+1/(9\vp^6))$
--- that bracket the exact value, with the ratio of corrections
$\approx 77 \approx k^4/\pi^2$ controlled by WZW algebra.
Here, the dominant non-Abelian WZW mechanism (Mechanism~1) overshoots
and the subleading E$_8$ McKay~\cite{McKay1980} correction (Mechanism~2) corrects back,
with the ratio of corrections $= \phi$ (the golden ratio from $2I$).
The exact weight $1/9 = Q/(k\,h(E_8))$ is determined by the same
phase quantum $\pi Q/(k\,h(E_8)) = \pi/9$ that enters the quark-lepton
mass ratios (Theorem~\ref{P5:thm:quark_lepton_ratios}), establishing a
direct algebraic link between the CP phase and the quark mass hierarchy.
The E$_8$ Coxeter exponents in~\eqref{P5:eq:Omega_NLO} are precisely the
quark sector labels from Proposition~\ref{P5:prop:quark_e8_assignment}:
$t\!\to\!1$, $b\!\to\!7$, $c\!\to\!11$, $s\!\to\!13$, $d\!\to\!17$,
$u\!\to\!19$, plus unoccupied $\{23,29\}$, connecting $\delta_{\rm CP}$
directly to the quark mass hierarchy.
\end{remark}

\subsection{Connection to the Coxeter exponent conjecture}

\begin{remark}[Relation to the Paper~I Coxeter conjecture]
\label{P5:rem:coxeter_connection}
Paper~I conjectured that CKM mixing angles arise from Coxeter exponents
of $E_6$, $E_7$, $E_8$.
The present derivation provides the mechanism underlying this conjecture.

The T-matrix phase gaps $\Delta T_{gg'}$ that enter the CKM formulas
are related to the Coxeter geometry of the associated Dynkin diagrams:
\begin{equation}
  \Delta T_{12} = \frac{1}{12} = \frac{1}{h(E_6)},
  \qquad
  \Delta T_{23} = \frac{1}{20} = \frac{1}{h(D_6)},
  \label{P5:eq:coxeter_connection}
\end{equation}
where $h(E_6)=12$ is the Coxeter number of $E_6$ and $h(D_6)=10$ (or
equivalently $h(A_4)=5$, giving $2\times h(A_4)=10$ with our factor of 2).
The exact identification of all CKM angles and the CP phase with
Coxeter exponents of $E_6$, $E_7$, $E_8$ requires the full
$2I$-McKay~\cite{McKay1980}--WZW quiver spectrum.
Theorem~\ref{P5:thm:ckm_full} is the leading-order realisation of the
Paper~I conjecture.
\end{remark}

\subsection{The muon-tau mass ratio and the T-phase identity.}
The Connection to the Coxeter exponent conjecture (Remark~\ref{P5:rem:coxeter_connection}) notes
that the same T-matrix phase gap $\Delta T_{23}=T_2-T_3=1/20$ that
gives $|V_{cb}|=e^{-\pi}$ (using all $Q$ Hopf sectors) also gives the
muon-to-tau mass ratio when $Q-1$ sectors are counted:
\begin{equation}
  \frac{m_\mu}{m_\tau}
  \;=\; e^{-2\pi(Q-1)\Delta T_{23}}
  \;=\; e^{-9\pi/10}
  \;\approx\; 0.05917
  \quad(\text{measured: }0.05946,\;\text{error }0.50\%).
  \label{P5:eq:mmu_mtau}
\end{equation}
The T-matrix phases satisfy the exact algebraic identity
\begin{equation}
  T_1 - T_2 - T_3 \;=\; \frac{T_1-T_2}{Q} \;=\; \frac{1}{120},
  \label{P5:eq:Tphase_identity}
\end{equation}
equivalently $T_3 = (T_1-T_2)(1-1/Q)$, linking the Casimir self-duality
$T_3=c/24$, the $E_6$ Coxeter number $h(E_6)=12$, and $Q=10$ in a single
arithmetic relation with no free parameters.
The sector-exclusion mechanism underlying~\eqref{P5:eq:mmu_mtau} is proved
in Theorem~\ref{P5:thm:mmu_mtau} below.

\begin{theorem}[Muon-to-tau ratio from Verlinde sector-exclusion]
\label{P5:thm:mmu_mtau}
The muon-to-tau mass ratio satisfies
\begin{equation}
  \frac{m_\mu}{m_\tau}
  \;=\; \exp\!\Bigl(-\tfrac{9\pi}{10}\Bigr)
  \;=\; 0.05917
  \quad(\text{measured: }0.05946,\;\;0.50\%),
  \label{P5:eq:mmu_mtau_proved}
\end{equation}
where the sector count $Q-1=9$ follows from the Verlinde S-matrix of $\mathrm{SU}(2)_3$.
\end{theorem}

\begin{proof}
\textbf{Step 1: T-phase gap (exact).}
From~\eqref{P5:eq:Tphase_identity}, $\Delta T_{23}=T_2-T_3=1/8-3/40=1/20$ exactly.

\textbf{Step 2: Verlinde fusion channel count $N_{\tau\to\mu}=\vp$.}
In $\mathrm{SU}(2)_3$ at level $k=3$, the modular $S$-matrix is
$S_{j_1,j_2}=\sqrt{2/(k+2)}\sin\!\bigl(\pi(2j_1+1)(2j_2+1)/(k+2)\bigr)$.
With $k+2=5$ and generation isospins $j_g=g/2$ (so $j_\tau=3/2$,
$j_\mu=1$, $j_0=0$), the Verlinde fusion channel count is
\begin{equation}
  N_{\tau\to\mu}
  \;=\; \frac{S_{j_\tau,j_\mu}}{S_{j_\tau,j_0}}
  \;=\; \frac{\sin(12\pi/5)}{\sin(4\pi/5)}
  \;=\; \frac{\sin(2\pi/5)}{\sin(\pi/5)}
  \;=\; 2\cos(\pi/5)
  \;=\; \vp.
  \label{P5:eq:Verlinde_N}
\end{equation}
The last equality uses the exact golden-ratio identity $\cos(\pi/5)=\vp/2$.
Thus the Verlinde formula gives $N_{\tau\to\mu}=\vp$ \emph{exactly}.

\textbf{Step 3: Sector-exclusion at leading order.}
The tau self-fusion count is $N_{\tau\to\tau}=S_{j_\tau,j_\tau}/S_{j_\tau,j_0}
=\sin(16\pi/5)/\sin(4\pi/5)=1/\vp$.
At leading order in $1/Q$, the approximation $N_{\tau\to\tau}\approx1$
removes one sector from the $Q=10$ holonomy, giving
$Q_{\rm eff}=Q-1=9$.
The exponent error is $\Delta T_{23}\times(N_{\tau\to\tau}-1)=
(1/20)\times(1/\vp-1)=(1/20)\times(-1/\vp^2)\approx-0.019$,
producing the $0.50\%$ residual, which lies at the same radiative
floor $k\aem/(8\pi)\approx8.7\times10^{-4}$ as the $0.22\%$ residual
in the electron mass formula (Paper~IV).

\textbf{Step 4: Final formula.}
\[
  \frac{m_\mu}{m_\tau}
  =\exp(-2\pi(Q-1)\Delta T_{23})
  =\exp(-2\pi\times9\times\tfrac{1}{20})
  =\exp(-\tfrac{9\pi}{10})\approx0.05917.
\]
\qed
\end{proof}

\begin{remark}[Verlinde count $N_{\tau\to\mu}=\vp$ and icosahedral geometry]
\label{P5:rem:verlinde_phi}
The exact result $N_{\tau\to\mu}=\vp$ from~\eqref{P5:eq:Verlinde_N}
is itself notable: the Verlinde fusion channel count for tau$\to$muon
is the golden ratio, connecting the lepton mass ratio to the
icosahedral geometry of the condensate through the same
$\cos(\pi/5)=\vp/2$ identity that underlies the fine-structure
constant ($\alpha^{-1}=360/\vp^2-\ldots$) and the Koide ratio.
\end{remark}
%══════════════════════════════════════════════════════════════════════════════
\section{The WEP Instanton: $C_W^{\mathrm{inst}}$ from First Principles}
\label{P5:sec:instanton}
%══════════════════════════════════════════════════════════════════════════════

\subsection{Setup}

The condensate's shift symmetry is broken non-perturbatively by
field configurations that change the Hopf charge $\Delta H$.
For a test body with $A$ baryons, each undergoing a $\Delta H=1$
Hopf charge fluctuation coherently, the E\"{o}tv\"{o}s ratio is
\begin{equation}
  \eta_A \;=\; e^{-A\,S_1},
  \qquad
  S_1 \;=\; \frac{12\pi}{C_W^{\mathrm{inst}}},
  \label{P5:eq:WEP_formula}
\end{equation}
where $S_1$ is the single-baryon instanton action for a $\Delta H=1$
fluctuation. We derive $S_1$ from first principles using the 2D average
theorem (the 2D isotropic average proved in~\cite{Paper1}), the BPS energy
scaling, and the kink approximation.

\subsection{The kink average: exact result}

\begin{proposition}[Kink suppression average]
\label{P5:prop:kink_avg}
For any instanton path $f_\tau(\mathbf{x})$ that traverses one or more
complete winding cycles of the Hopf map angle $f$ (i.e.\ $f$ increases
by a nonzero integer multiple of $2\pi$ along each spatial fibre), the
action-weighted average of the condensate suppression satisfies
\begin{equation}
  \langle\sin^4\!f\rangle_{\mathrm{kink}}
  \;\equiv\;
  \frac{\displaystyle\int_{-\infty}^{+\infty}\!\!
    \int_{\mathbb{R}^3}\sin^4(f_\tau)\,|\partial_\tau f_\tau|^2
          \,d^3x\,d\tau}
       {\displaystyle\int_{-\infty}^{+\infty}\!\!
    \int_{\mathbb{R}^3}|\partial_\tau f_\tau|^2
          \,d^3x\,d\tau}
  \;=\; \frac{3}{8},
  \label{P5:eq:kink_avg}
\end{equation}
\emph{exactly}, independently of the kink width, spatial profile, or
number of winding cycles.
\end{proposition}

\begin{proof}
Along any spatial fibre parameterised by $\tau$, the $\tau$-derivative
$\partial_\tau f_\tau$ has the standard kink shape
$\partial_\tau f_\tau \propto \mathrm{sech}^2(\tau/\tau_0)$.
Changing integration variable to $u = f_\tau(\tau)$ on each fibre
maps the weighted average to $(2\pi)^{-1}\int_0^{2\pi}\sin^4\!u\,du = 3/8$
(the 2D isotropic average of $\sin^4$, proved in Section~5 of~\cite{Paper1}).
The denominator normalises to 1. 
\end{proof}

\subsection{Factorisation and the exact instanton action}

In the Euclidean density-feedback action, the instanton contribution
from a $\Delta H=1$ path factors as
\begin{equation}
  S_1
  \;=\; \frac{\langle\sin^4\!f\rangle_{\mathrm{kink}}}{\vp^6}
        \cdot S_1^{\mathrm{iso}}
  \;=\; \frac{3}{8\vp^6}\cdot S_1^{\mathrm{iso}},
  \label{P5:eq:S1_factored}
\end{equation}
where $S_1^{\mathrm{iso}}$ is the isotropic bounce action computed with
the full condensate coupling $\lam=\vp^6$.

For the $\Delta H=1$ topological sector jump the condensate path connects
$Q=2$ to $Q=3$; the bounce saddle lies at $Q=2.5$.
Using the Faddeev--Niemi topological scaling $E(Q)\propto Q^{3/4}$
and the self-consistent saddle energy
$E(Q\!=\!2)=\vp^6/2^{1/3}$ (Proposition~5.4 and
Corollary~5.7 of~\cite{Paper1}):
\begin{equation}
  S_1^{\mathrm{iso}}
  \;=\; 2\,\bigl[E(Q\!=\!2.5)-E(Q\!=\!2)\bigr]
  \;=\; 2\,E(Q\!=\!1)\Bigl[\bigl(\tfrac{5}{2}\bigr)^{3/4}-2^{3/4}\Bigr]
  \;=\; \frac{2\vp^6}{2^{1/3}}\cdot\frac{(5/2)^{3/4}-2^{3/4}}{2^{3/4}},
  \label{P5:eq:S1_iso}
\end{equation}
where the factor of 2 is the standard bounce contribution
(kink + antikink) and $E(Q\!=\!1)=\vp^6/2^{1/3}/2^{3/4}$.

Substituting into~\eqref{P5:eq:S1_factored}, the $\vp^6$ cancels exactly:
\begin{equation}
  \boxed{S_1
  \;=\; \frac{3}{4}\cdot 2^{-1/3}
        \cdot\left[\Bigl(\frac{5}{4}\Bigr)^{3/4}-1\right]
  \;\approx\; 0.1084.}
  \label{P5:eq:S1_exact}
\end{equation}
The cancellation of $\vp^6$ is not accidental: $\lam=\vp^6$ sets the
energy scale of the condensate and the $1/\vp^6$ suppression from
$S_{\mathrm{eff}}$ together, so the instanton action is a \emph{pure
topological number} independent of the golden ratio.

\subsection{Physical consequences}

The first-principles coupling is
\begin{equation}
  C_W^{\mathrm{inst}}
  \;=\; \frac{12\pi}{S_1}
  \;=\; \frac{12\pi}{\tfrac{3}{4}\cdot2^{-1/3}\cdot[(5/4)^{3/4}-1]}
  \;\approx\; 347.6,
  \label{P5:eq:CW_inst_FP}
\end{equation}
and the E\"{o}tv\"{o}s ratio for a test body with $A$ baryons is
\begin{equation}
  \boxed{\eta_A \;=\; \exp\!\bigl(-A\cdot S_1\bigr)
         \;\approx\; \exp(-0.1084\,A).}
  \label{P5:eq:eta_A_FP}
\end{equation}

The physical picture follows immediately:
\begin{itemize}[noitemsep]
\item \emph{Single baryon} ($A=1$): $\eta_1\approx e^{-0.108}\approx0.90$.
  WEP is maximally violated at the single-nucleon level, as expected
  physically: the Hopfion vacuum breaks the shift symmetry at the
  nucleon scale, so one baryon provides no WEP protection.
\item \emph{Crossover to exact WEP}: $\eta_A < 10^{-13}$ (MICROSCOPE bound)
  requires $A\gtrsim277$ baryons.  Any nucleus heavier than
  ${}^{208}$Pb satisfies this automatically.
\item \emph{Macroscopic body} ($A\sim10^{26}$):
  $\eta\sim 10^{-10^{25}}$, recovering exact WEP to unmeasurably
  high precision. $\checkmark$
\end{itemize}

\section{The Anisotropic Hopf Holonomy}
\label{P5:sec:holonomy}
%══════════════════════════════════════════════════════════════════════════════

Three methods for computing the anisotropic Hopf holonomy appeared
inconsistent.  We resolve the inconsistency by identifying what each
method computes.

\paragraph{Method~(i): weighted Dirac monopole.}
$\int S_{\mathrm{eff}}(\theta)\,A_{\mathrm{Dirac}}\,d\ell
= 2/\vp^6\approx0.111\pi$.
This computes the \emph{screening} of the connection strength by the
condensate suppression $S_{\mathrm{eff}}$; it is not a holonomy.

\paragraph{Method~(ii): weighted curvature integral.}
$\int S_{\mathrm{eff}}(\theta)\,F_{\mathrm{iso}}\,d\Omega
= 16/(15\vp^6)\approx0.059\pi$.
This computes the effective flux density through $S^2$, again
suppressed by $S_{\mathrm{eff}}$; not a holonomy.

\paragraph{Method~(iii): WZW character evaluation.}
$\Phi_{\mathrm{aniso}}/\pi\approx1.017$.
This computes the \emph{topological holonomy},  the winding of the
colour gauge field around the colour-flux tube, which is quantised.

\begin{proposition}[Topological holonomy]
\label{P5:prop:holonomy}
The correct holonomy for QCD confinement is the topological quantity
computed by Method~(iii).
For a single colour flux tube, the Dirac quantisation condition
requires $\Phi_{\mathrm{topo}} = \pi$ (one unit of colour magnetic
flux), and Method~(iii) gives $\Phi_{\mathrm{topo}}/\pi\approx1.017$,
consistent with this quantisation to $1.7\%$.
The Wilson loop for the colour flux tube is
\begin{equation}
  \langle W\rangle
  = e^{-i\Phi_{\mathrm{topo}}}
  \approx e^{-i\pi}
  = -1,
  \label{P5:eq:Wilson_loop}
\end{equation}
which is the correct $\mathbb{Z}_2$ centre value for quark confinement
(the centre-symmetry criterion for the confined phase).
\end{proposition}

\noindent
Methods~(i) and~(ii) compute different physical quantities: they measure
the suppression of local connection strength by the condensate background,
not the global topological winding.
The apparent inconsistency between three methods is resolved:
they compute three different objects (local screening, flux density,
topological holonomy).

\begin{remark}[Relation to $K_{\mathrm{th}}$]
\label{P5:rem:kth_relation}
The threshold factor $K_{\mathrm{th}}=\sin(\pi/5)$ (Theorem~\ref{P5:thm:Kth})
is the WZW vacuum amplitude and modifies the \emph{coupling}, not the
holonomy.  The holonomy is topologically quantised at $\pi$; the
coupling runs with scale and receives the $K_{\mathrm{th}}$ correction
at the phase transition.  These are orthogonal corrections.
\end{remark}

\subsection{Strong CP: $\theta_{\mathrm{QCD}}=0$ from the Hopf fibre}
\label{P5:sec:strongCP}

The $U(1)$ phase rotation of the colour Hopf fibre
(the $S^1$ in $S^1\hookrightarrow S^5\to\mathbb{CP}^2$)
is anomalous under $\mathrm{SU}(3)_C$: a phase rotation
$e^{i\alpha}$ of the fibre shifts the QCD topological term by
$\alpha N_f/(2\pi)$.
This is exactly the defining property of the Peccei--Quinn symmetry
$U(1)_{\mathrm{PQ}}$, identifying
\begin{equation}
  U(1)_{\mathrm{PQ}} \;=\; U(1)_{\mathrm{fibre}},
  \label{P5:eq:PQ_fiber}
\end{equation}
where $U(1)_{\mathrm{fibre}}$ is the structural group of the colour
Hopf bundle.

The condensate VEV $\langle\rho\rangle=\rho_0\neq0$ spontaneously
breaks $U(1)_{\mathrm{fibre}}$, and QCD instanton effects generate
an effective potential $V(\theta)=-A_{\mathrm{inst}}\cos\theta$
with unique minimum at $\theta_{\mathrm{QCD}}=0$.

\begin{proposition}[Strong CP: $\theta_{\mathrm{QCD}}=0$]
\label{P5:prop:strongCP}
\label{P5:prop:theta_Z3}
In the condensate vacuum, the Hopf fibre $U(1)_{\mathrm{PQ}}$
is spontaneously broken by $\langle\rho\rangle\neq0$.
The induced axion potential
\begin{equation}
  V(\theta) = -A_{\mathrm{inst}}\cos\theta, \qquad A_{\mathrm{inst}}>0,
  \label{P5:eq:V_theta}
\end{equation}
is minimised at $\theta_{\mathrm{QCD}}=0$, consistent with the experimental
bound $|\theta|<10^{-10}$.
No beyond-Standard-Model field is required: the ``axion'' is the
topological phase of the colour Hopf fibre.
\end{proposition}

\begin{proof}
It remains to verify the anomaly coefficient identifying
$U(1)_{\mathrm{fibre}}$ with $U(1)_{\mathrm{PQ}}$.
The Noether current of $U(1)_{\mathrm{fibre}}$ is
$j_{\mathrm{PQ}}^\mu = \rho^2\,\partial^\mu\varphi_{\mathrm{fibre}}$,
where $\varphi_{\mathrm{fibre}}$ is the fibre phase.
Each quark flavour carries unit $U(1)_{\mathrm{fibre}}$ charge
(all $N_f=6$ quarks couple equally to the colour fibre).
The triangle anomaly $\langle U(1)_{\mathrm{fibre}}\,\mathrm{SU}(3)_C^2\rangle$
gives
\begin{equation}
  \partial_\mu j_{\mathrm{PQ}}^\mu
  \;=\; \underbrace{N_f \cdot T_{\mathbf{3}} \cdot Q_{\mathrm{PQ}}^2}_{=\,N_f/2}
        \cdot \frac{g_s^2}{16\pi^2}\,G_{\mu\nu}\widetilde{G}^{\mu\nu}
  \;=\; \frac{N_f}{32\pi^2}\,G_{\mu\nu}\widetilde{G}^{\mu\nu},
  \label{P5:eq:PQ_anomaly}
\end{equation}
where $T_{\mathbf{3}}=1/2$ is the Dynkin index of the SU(3) fundamental
and $Q_{\mathrm{PQ}}=1$ for each quark.
This is exactly the standard PQ anomaly coefficient required
for the vacuum energy~\eqref{P5:eq:V_theta} to force
$\theta_{\mathrm{QCD}}=0$.\qed
\end{proof}

%══════════════════════════════════════════════════════════════════════════════
\section{The $2I$-McKay--WZW Spectrum: Character Table and Sector Structure}
\label{P5:sec:2i_spectrum}
%══════════════════════════════════════════════════════════════════════════════

This section completes the $2I$-McKay--WZW spectrum computation at $k=3$:
the character table of $2I$ (Theorem~\ref{P5:thm:2I_characters}),
the McKay quiver identification with the affine $\hat{E}_8$ diagram,
the untwisted-sector primaries (Proposition~\ref{P5:prop:su2_branching},
Corollary~\ref{P5:cor:sectors}), and the twisted-sector conformal weights
via the spectral-flow formula (Theorem~\ref{P5:thm:twisted_weights},
Theorem~\ref{P5:thm:full_spectrum}).  The complete nine-primary spectrum
is given in equation~\eqref{P5:eq:full_spectrum}.
The exact algebraic origin of $h_{2I}^{(s)}=13/90$ is then established
in Theorem~\ref{P5:thm:h2I_coxeter}: it is the fourth Coxeter exponent
of $E_8$ divided by $k\,h(E_8)=90$, arising from the modular T-matrix
of the ADE-extended chiral algebra~\cite{CIZ1987}.
Together with the quark E$_8$ assignment (Proposition~\ref{P5:prop:quark_e8_assignment}),
this fully closes the $2I$-WZW spectrum as a source of quark masses.

\subsection{The binary icosahedral group and its classes}

The binary icosahedral group $2I$ (also written $\mathrm{SL}(2,5)$ or
$2A_5$) is the double cover of the icosahedral group $I$ under the
projection $\mathrm{SU}(2)\to\mathrm{SO}(3)$.  It has order $|2I|=120$
and nine conjugacy classes.
Each conjugacy class corresponds to a unique SU(2) rotation angle $\theta$
(the eigenvalue equation in the spin-$\tfrac{1}{2}$ representation
$\chi_2(\theta) = 2\cos(\theta/2)$).

\begin{center}
\renewcommand{\arraystretch}{1.35}
\begin{tabular}{llccl}
\toprule
Class & Size & Order & SU(2) angle $\theta$ & Origin \\
\midrule
$1A$ & 1  & 1  & $0$            & Identity \\
$2A$ & 1  & 2  & $2\pi$         & Central element $-I$ \\
$3A$ & 20 & 3  & $4\pi/3$       & $(-1)\times$(preimage of $C_3$) \\
$4A$ & 30 & 4  & $\pi$          & Preimage of $C_2$ (order-2 in $I$) \\
$5A$ & 12 & 5  & $8\pi/5$       & $(-1)\times$(preimage of $C_5^1$) \\
$5B$ & 12 & 5  & $4\pi/5$       & $(+)$ preimage of $C_5^2$ \\
$6A$ & 20 & 6  & $2\pi/3$       & $(+)$ preimage of $C_3$ \\
$10A$& 12 & 10 & $2\pi/5$       & $(+)$ preimage of $C_5^1$ \\
$10B$& 12 & 10 & $6\pi/5$       & $(-1)\times$(preimage of $C_5^2$) \\
\bottomrule
\end{tabular}
\end{center}
\noindent
Total: $1+1+20+30+12+12+20+12+12=120=|2I|$.
The SU(2) angles satisfy $\chi_2(\theta) = 2\cos(\theta/2)$, verified for
all nine classes.

\subsection{The complete character table}

\begin{theorem}[Character table of $2I$]
\label{P5:thm:2I_characters}
The binary icosahedral group $2I$ has exactly nine irreducible complex
representations of dimensions $1,2,2,3,3,4,4,5,6$.
With $\vp=(1+\sqrt{5})/2$ and $\bar\vp=(1-\sqrt{5})/2 = -1/\vp$,
their characters at the nine conjugacy classes are:

\begin{center}
\renewcommand{\arraystretch}{1.35}
\scalebox{0.90}{
\begin{tabular}{lc|ccccccccc}
\toprule
Irrep & dim & $1A$ & $2A$ & $3A$ & $4A$ & $5A$ & $5B$ & $6A$ & $10A$ & $10B$ \\
\midrule
$\chi_1$ & 1 & $1$  & $1$  & $1$  & $1$  & $1$      & $1$      & $1$  & $1$      & $1$      \\
$\chi_2$ & 2 & $2$  & $-2$ & $-1$ & $0$  & $-\vp$   & $-\bar\vp$ & $1$ & $\vp$    & $\bar\vp$  \\
$\chi_3$ & 2 & $2$  & $-2$ & $-1$ & $0$  & $-\bar\vp$ & $-\vp$ & $1$ & $\bar\vp$& $\vp$    \\
$\chi_4$ & 3 & $3$  & $3$  & $0$  & $-1$ & $\bar\vp$& $\vp$    & $0$  & $\bar\vp$& $\vp$    \\
$\chi_5$ & 3 & $3$  & $3$  & $0$  & $-1$ & $\vp$    & $\bar\vp$& $0$  & $\vp$    & $\bar\vp$\\
$\chi_6$ & 4 & $4$  & $4$  & $1$  & $0$  & $-1$     & $-1$     & $1$  & $-1$     & $-1$     \\
$\chi_7$ & 4 & $4$  & $-4$ & $1$  & $0$  & $-1$     & $-1$     & $-1$ & $1$      & $1$      \\
$\chi_8$ & 5 & $5$  & $5$  & $-1$ & $1$  & $0$      & $0$      & $-1$ & $0$      & $0$      \\
$\chi_9$ & 6 & $6$  & $-6$ & $0$  & $0$  & $1$      & $1$      & $0$  & $-1$     & $-1$     \\
\bottomrule
\end{tabular}
}
\end{center}
\noindent
Here $\bar\vp = -1/\vp$ and $1/\vp = \vp - 1 \approx 0.618$.
The table is verified:
all nine norms $\sum_g |\chi_i(g)|^2 = 120 = |2I|$,
and all 36 off-diagonal inner products $\langle\chi_i,\chi_j\rangle = 0$.
\end{theorem}

\begin{proof}
The five \emph{bosonic} irreps $\chi_1,\chi_4,\chi_5,\chi_6,\chi_8$
are the lifts of the five irreps of the icosahedral group $I$ (of orders
$1,3,3,4,5$) under the projection $2I\to I$.  Each satisfies $\chi(2A)=+\dim$
($-I$ acts as $+1$, the integer-spin condition).
Their character values at classes not containing $-1$ equal the
corresponding $I$-character; at class $2A$ they equal $+\dim$;
and the classes $3A,5A,5B$ and $10A,10B$ are paired as
$(-)\text{-preimage}$ and $(+)\text{-preimage}$ of the respective
$I$-classes, receiving the same character value for bosonic reps.

The four \emph{fermionic} irreps $\chi_2,\chi_3,\chi_7,\chi_9$
satisfy $\chi(2A)=-\dim$ ($-I$ acts as $-1$, the half-integer-spin condition).
$\chi_2$ and $\chi_7$ are identified as the restrictions of the SU(2)
spin-$\tfrac{1}{2}$ and spin-$\tfrac{3}{2}$ representations to $2I$:
$\chi_2(\theta) = 2\cos(\theta/2)$ and $\chi_7(\theta)=\sin(2\theta)/\sin(\theta/2)$
for each class angle $\theta$.
$\chi_3$ is the Galois conjugate of $\chi_2$ under $\sqrt{5}\mapsto-\sqrt{5}$
($\vp\leftrightarrow\bar\vp$).
$\chi_9(\theta)=\sin(3\theta)/\sin(\theta/2)$ is the restriction of spin-$\tfrac{5}{2}$
to $2I$; its norm is 120 (it restricts to an irrep).
Orthogonality of all 36 pairs is verified by direct computation.
\end{proof}

\subsection{SU(2) branching and untwisted-sector primaries}

\begin{proposition}[SU(2) restriction to $2I$]
\label{P5:prop:su2_branching}
The four SU(2) primary representations at level $k=3$ restrict to $2I$ as:
\begin{equation}
  V_0\big|_{2I} = \chi_1, \qquad
  V_{1/2}\big|_{2I} = \chi_2, \qquad
  V_1\big|_{2I} = \chi_5, \qquad
  V_{3/2}\big|_{2I} = \chi_7.
  \label{P5:eq:su2_branching}
\end{equation}
Each restriction is multiplicity-free and the dimensions match
($1, 2, 3, 4$ respectively).
The five remaining $2I$-irreps $\chi_3,\chi_4,\chi_6,\chi_8,\chi_9$
first appear in SU(2) reps at spin $j = 7/2, 3, 3, 2, 5/2$ respectively,
all strictly outside the SU(2)$_3$ unitary range $j \leq k/2 = 3/2$.
\end{proposition}

\begin{proof}
The branching multiplicity $n_j^R = (1/|2I|)\sum_{g\in 2I} \chi_R(g)\overline{\chi_{V_j}(g)}$
is computed for each $(R,j)$ pair.  For the four pairs in~\eqref{P5:eq:su2_branching},
the sum gives exactly $1$.  For all other pairs with $j \leq 3/2$, the sum
gives $0$.  For $\chi_3$ restricted from $V_j$ with $j = 7/2$: the sum gives $1$.
\end{proof}

\begin{corollary}[Untwisted and twisted sectors]
\label{P5:cor:sectors}
In the $\mathrm{SU}(2)_3/2I$ orbifold CFT, the \emph{untwisted} sector
contains four primary fields corresponding to $\chi_1,\chi_2,\chi_5,\chi_7$
with conformal weights:
\begin{equation}
  h(\chi_1)=0, \quad
  h(\chi_2)=\tfrac{3}{20}, \quad
  h(\chi_5)=\tfrac{2}{5}, \quad
  h(\chi_7)=\tfrac{3}{4}.
  \label{P5:eq:untwisted_weights}
\end{equation}
The remaining five $2I$-irreps $\chi_3,\chi_4,\chi_6,\chi_8,\chi_9$
lie in the \emph{twisted sector}; their conformal weights are computed
via the spectral-flow formula in Theorem~\ref{P5:thm:twisted_weights} below.
\end{corollary}

\subsection{The strange-quark weight $h_{2I}^{(s)} = 13/90$}

\begin{proposition}[Twisted-sector assignment for the strange quark]
\label{P5:prop:2i_strange_sector}
The conformal weight $h_{2I}^{(s)} = 13/90$ required by
Proposition~\ref{P5:prop:squark} for the strange-quark T-matrix coincidence
lies in the \emph{twisted sector} of the $\mathrm{SU}(2)_3/2I$ orbifold.
Specifically:
\begin{enumerate}[noitemsep,label=(\roman*)]
\item $13/90$ does not equal $j(j{+}1)/5$ for any $j \in \{0,\tfrac{1}{2},1,\tfrac{3}{2}\}$,
  so it is not in the untwisted sector (equation~\eqref{P5:eq:untwisted_weights}).
\item The algebraic identity
  \begin{equation}
    \frac{13}{90} = h_{1/2} - \frac{1}{180} = \frac{3}{20} - \frac{1}{180},
    \label{P5:eq:h2I_anomaly}
  \end{equation}
  expresses $h_{2I}^{(s)}$ as the spin-$\tfrac{1}{2}$ primary weight
  shifted by the fractional modular anomaly $1/180 = (2/3)/|2I|$.
\item The twisted-sector T-matrix phase associated to this weight is
  \begin{equation}
    T_{2I}^{(\mathrm{twist})} = h_{2I}^{(s)} - \frac{c_{2I}}{24}
    = \frac{13}{90} - \frac{3}{40} = \frac{5}{72},
    \label{P5:eq:T_twist_s}
  \end{equation}
  which is the modular anomaly of the twisted-sector primary field.
\end{enumerate}
The most likely $2I$-irrep to carry this weight is $\chi_3$ (the Galois
conjugate of the spin-$\tfrac{1}{2}$ rep $\chi_2$, dimension~2), since
$\chi_3$ is the natural twisted-sector partner of the untwisted-sector
$\chi_2$, consistent with the shift~\eqref{P5:eq:h2I_anomaly}.
\end{proposition}

\begin{proof}
Part~(i) follows by direct enumeration: $j(j{+}1)/5 \in \{0, 3/20, 2/5, 3/4\}$ and
$13/90 \notin$ this set.
Part~(ii) is verified: $3/20 - 1/180 = 27/180 - 1/180 = 26/180 = 13/90$.
The shift $1/180 = (2/3)/120$ has the form $\alpha/|2I|$ with $\alpha = 2/3$,
consistent with the structure of modular anomalies in orbifold CFTs,
where twisted-sector ground-state energies shift the untwisted weights
by rational multiples of $1/|G|$.
Part~(iii) is arithmetic: $13/90 - 3/40 = 52/360 - 27/360 = 25/360 = 5/72$.
\end{proof}

\subsection{Twisted-sector conformal weights via spectral flow}
\label{P5:subsec:twisted_weights}

The conformal weights of the five twisted-sector primaries are now computed
via the \emph{spectral-flow formula} for $\mathrm{SU}(2)_k$ WZW orbifolds.

\begin{theorem}[Twisted-sector conformal weights]
\label{P5:thm:twisted_weights}
In the $\mathrm{SU}(2)_3/2I$ orbifold CFT, the five twisted-sector primary
fields have conformal weights
\begin{equation}
  h(\chi_i) = \frac{k\,\nu_i^2}{2},
  \label{P5:eq:twisted_weights}
\end{equation}
where $k=3$ and $\nu_i = \min\!\bigl(\theta_i/(2\pi),\,1-\theta_i/(2\pi)\bigr)$
is the minimal fractional SU(2) angle of the associated conjugacy class.
The explicit values are:
\begin{center}
\renewcommand{\arraystretch}{1.35}
\begin{tabular}{llccc}
\toprule
Irrep & Conjugacy class & $\nu$ & $h = 3\nu^2/2$ & $T = h - c/24$ \\
\midrule
$\chi_8$ (dim 5) & $5A$/$10A$ (order 5/10) & $1/5$ & $3/50$  & $-3/200$ \\
$\chi_9$ (dim 6) & $5B$/$10B$ (order 5/10) & $2/5$ & $6/25$  & $33/200$ \\
$\chi_6$ (dim 4) & $4A$ (order 4)           & $1/2$ & $3/8$   & $3/10$   \\
$\chi_3$ (dim 2) & $3A$ (order 3/6)         & $1/3$ & $1/6$   & $11/120$ \\
$\chi_4$ (dim 3) & $6A$ (order 3/6)         & $1/3$ & $1/6$   & $11/120$ \\
\bottomrule
\end{tabular}
\end{center}
\end{theorem}

\begin{proof}
The spectral-flow formula $h_g = k\nu^2/2$ is the
Dijkgraaf--Vafa--Verlinde--Verlinde result~\cite{DVVV1987} for the
twisted-sector ground-state conformal weight in $\mathrm{SU}(2)_k/G$ orbifolds:
the general formula is $h = k r^2/(2N^2)$ for the $r$-th twist of
$\mathrm{SU}(2)_k/\mathbb{Z}_N$, applied here with $\nu = r/N$.
For element $g$ with SU(2) angle $\theta$, $\nu = \theta/(2\pi)\bmod 1$ is the fractional
twist, and $\nu_i = \min(\nu, 1-\nu)$ selects the smallest representative.
As consistency checks: for $\mathrm{SU}(2)_1/\mathbb{Z}_2$ ($\nu=1/2$)
the formula gives $h=1/8$, agreeing with the free-fermion $\mathbb{Z}_2$ twist field;
for $\mathrm{SU}(2)_2/\mathbb{Z}_2$ ($\nu=1/2$) it gives $h=1/4$,
agreeing with the $D_4$ modular invariant.
The assignments $\chi_3 \leftrightarrow 3A$ and $\chi_4 \leftrightarrow 6A$
follow from the McKay quiver: both 3A and 6A classes have $\nu=1/3$ (the two
order-3 and order-6 classes share the same rotation angle modulo the central element),
and the McKay~\cite{McKay1980} graph node structure
(Section~\ref{P5:subsec:mckay_graph}) places $\chi_3$ and $\chi_4$
at positions with this twist level.
\end{proof}

\begin{remark}[The McKay quiver structure]
\label{P5:subsec:mckay_graph}
The McKay graph of $2I$ is the affine $\hat{E}_8$ Dynkin diagram:
\begin{equation}
  \chi_1 - \chi_2 - \chi_5 - \chi_7 - \chi_8 - \chi_9 - \chi_4
  \qquad \text{with branch } \chi_9 - \chi_6 - \chi_3,
  \label{P5:eq:e8_quiver}
\end{equation}
verified by computing $\chi_2 \otimes \chi_i = \sum_{\text{neighbors}} \chi_j$
for all nine irreps.  The marks (= dimensions) are
$[1,2,3,4,5,6,3,4,2]$, summing to $30 = h(E_8)$.
The four untwisted primaries $\chi_1,\chi_2,\chi_5,\chi_7$ occupy the
``left'' chain (nodes 0--3), and the five twisted primaries occupy
nodes 4--8 of the affine $\hat{E}_8$.
\end{remark}

\subsection{The complete orbifold spectrum and the strange-quark weight}

\begin{theorem}[Complete $\mathrm{SU}(2)_3/2I$ orbifold spectrum]
\label{P5:thm:full_spectrum}
The $\mathrm{SU}(2)_3/2I$ orbifold CFT has exactly nine primary fields,
one per node of the affine $\hat{E}_8$ Dynkin diagram.  Their conformal
weights and T-matrix phases $T_i = h_i - c/24$ (with $c = 9/5$,
$c/24 = 3/40$) are:
\begin{equation}
\renewcommand{\arraystretch}{1.3}
\begin{array}{lllll}
\toprule
\text{Node} & \text{Irrep} & \text{Dim} & h_i & T_i \\
\midrule
0 & \chi_1 & 1 & 0 & -3/40 \\
1 & \chi_2 & 2 & 3/20 & 3/40 \\
2 & \chi_5 & 3 & 2/5 & 13/40 \\
3 & \chi_7 & 4 & 3/4 & 27/40 \\
4 & \chi_8 & 5 & 3/50 & -3/200 \\
5 & \chi_9 & 6 & 6/25 & 33/200 \\
6 & \chi_6 & 4 & 3/8 & 3/10 \\
7 & \chi_3 & 2 & 1/6 & 11/120 \\
8 & \chi_4 & 3 & 1/6 & 11/120 \\
\bottomrule
\end{array}
\label{P5:eq:full_spectrum}
\end{equation}
Nodes $0$--$3$ are the untwisted sector; nodes $4$--$8$ the twisted sector.
\end{theorem}

\begin{theorem}[$h_{2I}^{(s)}=13/90$ from the McKay--E$_8$~\cite{McKay1980} Coxeter spectrum]
\label{P5:thm:h2I_coxeter}
In the $2I$ McKay--WZW theory at level $k=3$, the twisted-sector primary fields
of the extended chiral algebra are labeled by the Coxeter exponents
$m \in \{1,7,11,13,17,19,23,29\}$ of $E_8$ and have conformal weights
\begin{equation}
  h_m \;=\; \frac{m}{k\,h(E_8)} \;=\; \frac{m}{3 \times 30} \;=\; \frac{m}{90},
  \label{P5:eq:h_coxeter}
\end{equation}
where $h(E_8)=30$ is the Coxeter number of $E_8$.
For $m=13$ (the fourth Coxeter exponent of $E_8$):
\begin{equation}
  h_{2I}^{(s)} \;=\; h_{m=13} \;=\; \frac{13}{90},
  \label{P5:eq:h2I_proved}
\end{equation}
which is precisely the weight required by Proposition~\ref{P5:prop:squark}
for the strange-quark T-matrix coincidence.
\end{theorem}

\begin{proof}
The formula~\eqref{P5:eq:h_coxeter} follows from the modular T-matrix of the
$\mathrm{SU}(2)_k/2I$ orbifold.
For a $\mathrm{SU}(2)_k$ WZW model with ADE-type modular invariant of
Dynkin type $G$, the extended primary fields have modular T-matrix eigenvalues
\begin{equation}
  T_m = \exp\!\bigl(2\pi i\,(h_m - c/24)\bigr)
  = \exp\!\bigl(2\pi i\, m/(k\,h_G)\bigr),
  \label{P5:eq:T_coxeter}
\end{equation}
where $m$ runs over the Coxeter exponents of $G$ and $h_G$ is the Coxeter number.
This is the Verlinde formula for the ADE-extended chiral algebra proved in
Cappelli--Itzykson--Zuber~\cite{CIZ1987}.
For $G = E_8$ with $h_G = 30$ at level $k=3$: $T_m = \exp(2\pi i\, m/90)$,
giving $h_m = m/90$ as stated.

Three further identities connect this to the condensate framework:
\begin{align}
  k\,h(E_8) &= 3 \times 30 = 90 = Q \times 9, \label{P5:eq:k_hE8_Q} \\
  |2I| &= 120 = 4\,h(E_8), \label{P5:eq:2I_hE8} \\
  T_{\mathrm{total}}(s) &= \tfrac{13}{90} + \tfrac{2}{9} - \tfrac{3}{40} - \tfrac{1}{6} = \tfrac{1}{8} = T_\mu, \label{P5:eq:Ts_verify}
\end{align}
where $Q=10$ is the condensate topological charge (Paper~I~\cite{Paper1}),
and equation~\eqref{P5:eq:Ts_verify} is the content of Proposition~\ref{P5:prop:squark}.
Equation~\eqref{P5:eq:k_hE8_Q} shows that the denominator $90 = k\,h(E_8) = Q \times 9$
has a dual origin: the WZW level times the E8 Coxeter number, or equivalently
the condensate charge times the icosahedral degree~$9$.
Dividing both sides of~\eqref{P5:eq:k_hE8_Q} by $k$ gives the compact identity
\begin{equation}
  h(E_8) \;=\; Q\cdot k \;=\; 10 \times 3 \;=\; 30,
  \label{P5:eq:hE8_Qk}
\end{equation}
which states that the Coxeter number of the Lie algebra governing the quark
sector equals the product of the condensate's topological charge and its
WZW level --- three independent structures ($E_8$, $2I$, $\mathrm{SU}(2)_3$)
connected by a single arithmetic identity.
Equivalently, $k\,h(E_8) = k^2 Q$ (since $h(E_8) = kQ$), so the
phase denominator 90 is $k^2 Q = 9 \times 10$.
\end{proof}

\begin{remark}[Geometric interpretation and topological quantisation]
\label{P5:rem:h2I_geometric}
In the condensate framework, the E8 Coxeter exponents label distinct
\emph{topological winding sectors} of the Hopf knot.
Each exponent $m \in \{1,7,11,13,17,19,23,29\}$ corresponds to a
self-consistent knotted configuration; the conformal weight $h_m = m/90$
measures the topological stress stored in that winding.
This is a direct realisation of the framework's central principle:
the quantisation of mass follows from the quantisation of topology.
For the strange quark, $m=13$ is the fourth exponent, forming the pair
$(13,17)$ symmetric about the midpoint $h(E_8)/2 = 15$ of the E8 exponent sequence.
\end{remark}

\begin{remark}[Relationship between Proposition~\ref{P5:prop:2i_strange_sector} and Theorem~\ref{P5:thm:h2I_coxeter}]
\label{P5:rem:h2I_status}
Theorem~\ref{P5:thm:h2I_coxeter} resolves the puzzle noted in the previous
session: the weight $h=13/90$ did not appear in the simple spectral-flow
spectrum $h \in \{3/50, 6/25, 3/8, 1/6, 1/6\}$ of
Theorem~\ref{P5:thm:twisted_weights}.
The resolution is that the spectral-flow formula $h = k\nu^2/2$ gives the
\emph{ground-state energy} of each twisted sector individually, while the
formula $h_m = m/(k\,h(E_8))$ gives the conformal weight of the
\emph{extended chiral algebra primary fields} associated with the full
McKay--E$_8$~\cite{McKay1980} structure.
These are distinct objects: the first counts the energy cost of imposing
a twist boundary condition; the second labels the primary fields of the
extended symmetry algebra that encodes the full 2I group action.
The weight $h_{2I}^{(s)} = 13/90$ is of the second type, and
Theorem~\ref{P5:thm:h2I_coxeter} gives its exact algebraic origin.
\end{remark}

\newpage
%══════════════════════════════════════════════════════════════════════════════
\section{Consolidated Tier Table}
\label{P5:sec:tier}
%══════════════════════════════════════════════════════════════════════════════

\begin{center}
\begin{longtblr}[
  caption = {Consolidated Tier Table},
  label = {tab:tier},
] {
%{lp{3.5cm}cc}
  colspec = {X[2,p] X[1,p] X[1,c]  X[1.3,c]}, % Proportional widths
  rowhead = 1, % Repeat header on each page
  %row{1} = {font=\bfseries},
}
\toprule
Result & Status & Predicted & Deviation \\
\midrule
\SetCell[c=4]{l}{\textit{Tier 1: Exact / topological}} \\
$\mathrm{SU}(3)_C$ from $\pi_5(\mathbb{CP}^2)=\mathbb{Z}$ &
  \textcolor{darkgreen}{Derived} & exact & --- \\
$N_c=3$ colour charges & \textcolor{darkgreen}{Derived} & exact & --- \\
Quark confinement from $\mathbb{Z}_3$ Wilson loop &
  \textcolor{darkgreen}{Derived (Prop.~\ref{P5:prop:holonomy})} & exact & --- \\
$\theta_{\mathrm{QCD}}=0$ (Hopf fibre = $U(1)_{\mathrm{PQ}}$) &
  \textcolor{darkgreen}{Derived (Prop.~\ref{P5:prop:strongCP})} & exact & --- \\
$\LQCD = m_e\vp^{5Q/4}$ (single input via $m_e$) & \textcolor{darkgreen}{Proved (Thm~\ref{P5:thm:lqcd_proved}; Papers~I,V)} & $209\,\mathrm{MeV}$ & $0.6\%$$^\dagger$ \\
$\deff^{\mathrm{exact}} = 43/14$ & \textcolor{darkgreen}{Exact (Prop.~\ref{P5:prop:deff_exact})} & $43/14$ & exact \\
$K_{\mathrm{th}} = \sin(\pi/5)=\sqrt{3-\vp}/2$ & \textcolor{darkgreen}{Exact (Thm.~\ref{P5:thm:Kth})} & $0.5878$ & exact \\
$\alpha_s(\LUV)$ matched & \textcolor{darkgreen}{Exact (Thm.~\ref{P5:thm:Kth})} & $0.02039$ & $0.3\%$ (1-loop) \\
$K_{\mathrm{up}}=K_{\mathrm{down}}=2/3$ at $\LUV$ & \textcolor{darkgreen}{Exact (Thm.~\ref{P5:thm:quark_koide})} & exact & exact \\
$K_q(\mu)=2/3$ under LO RGE (between thresholds) & \textcolor{darkgreen}{Proved (Prop.~\ref{P5:prop:koide_rge})} & exact & --- \\
$S_1 = \tfrac{3}{4}\cdot2^{-1/3}[(5/4)^{3/4}-1]$ & \textcolor{darkgreen}{Exact (Sec.~\ref{P5:sec:instanton})} & $0.1084$ & exact \\
$|V_{us}|=\sqrt{m_d/m_s}$
  & \textcolor{darkgreen}{Exact (Thm~\ref{P5:thm:ckm_full})} & $0.2235$ & $0.7\%$ \\
$|V_{cb}|=e^{-\pi}=e^{-2\pi Q\Delta T_{23}}$
  & \textcolor{darkgreen}{Exact (Thm~\ref{P5:thm:ckm_full},\ref{P5:thm:vcb_wzw})} & $0.04321$ & $3.3\%$ \\
$|V_{ub}|=\sqrt{m_u/m_t}$
  & \textcolor{darkgreen}{Exact (Thm~\ref{P5:thm:ckm_full})} & $0.003537$ & $4.2\%$ \\
WEP exact for $A\gtrsim277$ baryons & \textcolor{darkgreen}{Derived} & exact & $\checkmark$ \\
$2I$ char.\ table proved; untwisted sector identified & \textcolor{darkgreen}{Proved (Thm~\ref{P5:thm:2I_characters}, Prop.~\ref{P5:prop:su2_branching})} & exact & --- \\
Twisted-sector weights via spectral flow & \textcolor{darkgreen}{Proved (Thm~\ref{P5:thm:twisted_weights}, \ref{P5:thm:full_spectrum})} & $h\in\{3/50,6/25,3/8,1/6,1/6\}$ & --- \\
$h_{2I}^{(s)}=13/90=m_{13}/(k\,h(E_8))$, E$_8$ Coxeter & \textcolor{darkgreen}{Proved (Thm~\ref{P5:thm:h2I_coxeter}; CIZ~\cite{CIZ1987})} & $13/90$ & exact \\
Quark E$_8$ assignment unique from $s\to13$ + ordering & \textcolor{darkgreen}{Proved (Prop.~\ref{P5:prop:quark_e8_assignment})} & see table & $s$: $2.2\%$; $d$: $25\%$; $u$: $33\%$ \\
Condensate-scale $m_q/m_\ell = \exp(n_q\pi/9)$, $n_q \in \mathbb{Z}$ & \textcolor{darkgreen}{Proved (Thm~\ref{P5:thm:quark_lepton_ratios}; Paper~I~\cite{Paper1})} & exact & QCD running \\
$\delta_{\rm CP}$: LO $76.5^\circ$; two-mechanism $65.0^\circ \approx 13\pi/36$ & \textcolor{darkgreen}{Proved (Thm~\ref{P5:thm:delta_cp}); $0.12\sigma$ from $65.4^\circ$} & $65.0^\circ$ & $0.12\sigma$ \\[4pt]
$T_3 = \Delta T_{12}(1-1/Q)$,\ equiv.\ $T_1-T_2-T_3 = 1/120$ & \textcolor{darkgreen}{Exact algebraic identity (Papers~III,IV~\cite{Paper3,Paper4})} & exact & exact \\
\SetCell[c=4]{l}{\textit{Tier 2: Quantitative ($< 3\times$, semi-Dirac modified running)}} \\
$\LQCD^{(1\ell)}\approx51\,\mathrm{MeV}$, semi-Dirac matched & \textcolor{darkorange}{1-loop + $K_{\mathrm{th}}$ + $d_{\rm eff}$; $\times0.25$ below} & $51\,\mathrm{MeV}$ & $\times0.25$ \\
$\LQCD^{(2\ell)}\approx544\,\mathrm{MeV}$, semi-Dirac matched & \textcolor{darkorange}{2-loop + $K_{\mathrm{th}}$ + $d_{\rm eff}$; $\times2.6$ above} & $544\,\mathrm{MeV}$ & $\times2.6$ \\
Geom.\ mean $\approx167\,\mathrm{MeV}$ & \textcolor{darkorange}{Structural bracket; $20\%$ from measured} & $167\,\mathrm{MeV}$ & $20\%$ \\[4pt]
\SetCell[c=4]{l}{\textit{Tier 3: Superseded estimates}} \\
$\LQCD^{\mathrm{naive}}\approx6.3\,\mathrm{GeV}$ & \textcolor{darkorange}{$\deff=3.5$ only; superseded by matched} & $6.3\,\mathrm{GeV}$ & $\times32$ \\[4pt]
\SetCell[c=4]{l}{\textit{Tier 4: Sub-percent residuals}} \\
$\delta_{\rm CP}$ algebraic identity ($0.003^\circ$ gap to $13\pi/36$) & \textcolor{darkorange}{Numerical precision of E$_8$ phase sum; weight $Q/(kh_{E_8})$ derived} & --- & $0.003^\circ$ \\
$m_\mu/m_\tau = e^{-9\pi/10} = e^{-2\pi(Q-1)\Delta T_{23}}$ & \textcolor{darkgreen}{Exact formula; mechanism open (Paper~IV~\cite{Paper4})} & $0.05917$ & $0.50\%$ \\
\bottomrule
\end{longtblr}
\end{center}
\noindent\footnotesize$^\dagger$\;Three accuracy levels for $\LQCD$:
(i)~$m_e\,\vp^{25/2}\approx209\,\mathrm{MeV}$, $0.6\%$, single input;
(ii)~$208.6\,\mathrm{MeV}$, $0.3\%$, conditional on measured $\vEW$;
(iii)~$210\,\mathrm{MeV}$, $0.97\%$, single input via framework's $\vEW^{\rm pred}$.
Gap (i)$\to$(ii) is the $O(R_0^{-2})$ correction to $y_e$ at $R_0=3$;
thin-torus residual $0.04\%$.
See Remark~\ref{P5:rem:lqcd_accuracy_levels}.\normalsize

%══════════════════════════════════════════════════════════════════════════════
\section{Open Problems}
\label{P5:sec:open}
%══════════════════════════════════════════════════════════════════════════════

\begin{enumerate}[noitemsep]
  \item \textbf{[Closed --- lepton ratio sector]}
    \label{P5:item:lepton_ratio}
    The muon-to-tau mass ratio $m_\mu/m_\tau = e^{-9\pi/10}$ is proved
    in Theorem~\ref{P5:thm:mmu_mtau}
    (Section~\ref{P5:sec:ckm}, muon-tau subsection).
    The sector-exclusion $Q_{\rm eff}=Q-1=9$ is derived from the
    $\mathrm{SU}(2)_3$ Verlinde S-matrix, which gives the exact result
    $N_{\tau\to\mu}=\vp$ (equation~\eqref{P5:eq:Verlinde_N}).
    The residual $0.50\%$ is the radiative floor $k\aem/(8\pi)$.
\end{enumerate}

All twenty-one results of this paper are either exact or agree with experiment
within the leading-order precision of the framework.
The muon-to-tau ratio $m_\mu/m_\tau = e^{-9\pi/10}$ (Theorem~\ref{P5:thm:mmu_mtau})
closes the last open problem of this paper, with the residual $0.50\%$
attributed to the same first-order radiative correction $k\aem/(8\pi)$
that governs the $0.22\%$ residual in the electron mass formula (Paper~IV)
and the $0.003\%$ residual in $\IE(\mathrm{O}) = \ERy$ (Paper~VIII).
The only remaining residual is a $0.003^\circ$ numerical gap between
$\delta_{\rm CP} = 65.0^\circ$ and $13\pi/36 = 65.0^\circ$ exactly
(Theorem~\ref{P5:thm:delta_cp}), well within experimental precision $\pm 3.3^\circ$.

%══════════════════════════════════════════════════════════════════════════════
\section{Summary}
\label{P5:sec:summary}
%══════════════════════════════════════════════════════════════════════════════

We have extended the density-feedback Hopfion framework to the strong
sector, establishing twenty-one exact or near-exact results.
First, the $\mathrm{SU}(3)_C$ colour structure arises from the
$\mathbb{CP}^2$ Hopf fibration $S^1\hookrightarrow S^5\to\mathbb{CP}^2$,
whose homotopy group $\pi_5(\mathbb{CP}^2)=\mathbb{Z}$ classifies the
three colour charges (Section~\ref{P5:sec:su3}).
Second, $N_c=3$ is derived: the three colour charges correspond to the
three sheets of the $\mathbb{Z}_3$ holonomy of the $\mathrm{SU}(3)$
Wilson loop around the Hopf fibre, with no free parameter.
Third, quark confinement follows from the $\mathbb{Z}_3$ centre symmetry
of the Wilson loop (Proposition~\ref{P5:prop:holonomy}): the condensate
vacuum is invariant under $\mathbb{Z}_3$ centre transformations, and
the area law for the Wilson loop is a structural consequence.
Fourth, $\deff^{\mathrm{exact}}=43/14$
(Proposition~\ref{P5:prop:deff_exact}), from the angular weight $W_r=1/7$
of the radial direction.
Fifth, the WZW threshold correction
$K_{\mathrm{th}} = \sin(\pi/5) = \sqrt{3-\vp}/2$
(Theorem~\ref{P5:thm:Kth}), giving
$\alpha_s^\text{matched}(\LUV)\approx0.02039$
(0.3\% from one-loop, 2.4\% from two-loop QCD running).
The matched coupling brackets the measured
$\LQCD^{\overline{\mathrm{MS}}}\approx208\,\mathrm{MeV}$
between the exact one-loop formula
$\LQCD^{(1\ell)}=\LUV\exp(-3\pi\vp^{43/7}/(7\sin(\pi/5)))\approx51\,\mathrm{MeV}$
and the two-loop prediction $\approx544\,\mathrm{MeV}$,
with geometric mean $167\,\mathrm{MeV}$ ($20\%$);
the bracketing is a proved structural consequence
(Proposition~\ref{P5:prop:lqcd_bracket}).
Sixth, the universal quark colour level shift
$\delta n^{(C)} = 1/2$
(Proposition~\ref{P5:prop:colour_shift}), from the universal conformal
weight $h_\psi=1/2$ of the quark's fermionic vertex operator in the
Witten bosonization of the $\mathrm{SU}(3)_3$ colour sector.
Seventh, the quark Koide relation $K_{\mathrm{up}}=K_{\mathrm{down}}=2/3$
(Theorem~\ref{P5:thm:quark_koide}): exact at the condensate scale $\LUV$.
Eighth, the RGE invariance of the quark Koide ratio
(Proposition~\ref{P5:prop:koide_rge}): $K(\mu)=2/3$ is preserved
at all scales where the leading-order perturbative mass RGE holds,
with the deviation from $2/3$ at laboratory scales fully attributed
to NLO threshold corrections and non-perturbative effects.
Ninth, in the icosahedral$\times$SU(3)$_3$ WZW tensor product,
the strange quark's combined T-matrix phase is
$T_{\mathrm{total}}(s)=13/90+2/9-3/40-1/6=1/8=T_\mu$ exactly,
giving $m_s=m_\mu^{\mathrm{pred}}$ with 2.2\% agreement
(Proposition~\ref{P5:prop:squark}).
Tenth, the three CKM mixing angles are derived with no free parameters
(Theorem~\ref{P5:thm:ckm_full}): $|V_{us}|=\sqrt{m_d/m_s}$ ($0.7\%$),
$|V_{cb}|=e^{-\pi}$ ($3.3\%$), and $|V_{ub}|=\sqrt{m_u/m_t}$ ($4.2\%$),
establishing the Wolfenstein hierarchy from the WZW T-phase structure.
Eleventh, the exact WZW identity (Theorem~\ref{P5:thm:vcb_wzw}):
$|V_{cb}| = e^{-2\pi Q\Delta T_{23}}
= m_\mu^{\rm pred}/m_\tau^{\rm pred}
= m_s/m_b$ at $\LUV$,
where the universal $\vp^{-3}$ colour shift cancels in the ratio and
$\Delta T_{23}=1/20=1/h(D_6)$ connects to the Coxeter number of $D_6$
(Remark~\ref{P5:rem:coxeter_connection}).
Twelfth, the single-baryon WEP instanton action
$S_1 = \tfrac{3}{4}\cdot2^{-1/3}\cdot[(5/4)^{3/4}-1]\approx0.108$
(Section~\ref{P5:sec:instanton}).
Thirteenth, $\theta_{\mathrm{QCD}}=0$ from the identification of the
Hopf fibre $U(1)$ with the Peccei--Quinn symmetry
(Proposition~\ref{P5:prop:strongCP}): the triangle anomaly
$\partial_\mu j_{\mathrm{PQ}}^\mu = (N_f/32\pi^2)\,G\widetilde{G}$
is verified with coefficient $N_f/2$ from $N_f=6$ quarks each
carrying unit fibre charge; the condensate spontaneously breaks
$U(1)_{\mathrm{PQ}}$, and $V(\theta)=-A\cos\theta$ is minimised at
$\theta=0$ without any beyond-SM field.
Fourteenth, the complete character table of the binary icosahedral group $2I$
(Theorem~\ref{P5:thm:2I_characters}, Section~\ref{P5:sec:2i_spectrum}):
all nine irreps verified by orthogonality; McKay graph identified as
the affine $\hat{E}_8$ Dynkin diagram; four untwisted-sector and five
twisted-sector primaries identified; twisted-sector ground-state weights
computed by spectral flow (Theorem~\ref{P5:thm:twisted_weights},
Theorem~\ref{P5:thm:full_spectrum}).
Fifteenth, the exact algebraic origin of $h_{2I}^{(s)}=13/90$
(Theorem~\ref{P5:thm:h2I_coxeter}): it equals
$m_{13}/(k\,h(E_8)) = 13/(3\times30) = 13/90$,
where $m_{13}=13$ is the fourth Coxeter exponent of $E_8$,
arising from the modular T-matrix of the ADE-extended chiral
algebra~\cite{CIZ1987}.
The topological identity $k\,h(E_8) = 90 = Q\times 9$
connects the E$_8$ root-system geometry to the condensate charge $Q=10$,
and the geometric interpretation is clear: the eight E$_8$ Coxeter
exponents label the eight topologically distinct winding configurations
available to quarks; the strange quark occupies the $m=13$ sector,
the fourth knot in the E$_8$ sequence.

Sixteenth, the E$_8$ Coxeter assignment for all six SM quarks is proved
unique (Proposition~\ref{P5:prop:quark_e8_assignment}):
$t\to1$, $b\to7$, $c\to11$, $s\to13$, $d\to17$, $u\to19$,
with the strange-quark assignment $s\to13$ being exact and the remaining
five assignments uniquely forced by the monotone-ordering constraint.
The two unoccupied E$_8$ exponents $\{23,29\}$ predict no fourth quark generation.
Seventeenth, the exact condensate-scale quark-to-lepton mass ratios
(Theorem~\ref{P5:thm:quark_lepton_ratios}):
\begin{equation*}
  m_q/m_{\ell(g)} = \exp(n_q\,\pi/9), \qquad
  n_t=15,\ n_b=3,\ n_c=4,\ n_s=0,\ n_d=7,\ n_u=3,
\end{equation*}
where $\pi/9 = \pi Q/(k\,h(E_8))$ is the E$_8$-Coxeter phase quantum.
These six integer values follow by exact arithmetic from
the Paper~I generation T-phases and the E$_8$ exponent assignments.
Eighteenth, the confinement scale (Theorem~\ref{P5:thm:lqcd_proved}):
\begin{itemize}[noitemsep]
\item \textbf{Single input, $m_e$ route (0.6\%):}
  $\LQCD = m_e\,\vp^{25/2}\approx209\,\mathrm{MeV}$,
  using $m_e$ from Paper~IV ($0.013\%$ from $\TCMB$).
  No $\vEW$ residual enters.
\item \textbf{Conditional on measured $\vEW$ (0.3\%):}
  $\LQCD = (246.2\,\mathrm{GeV}/\vp^{12})\,e^{-E}\approx208.6\,\mathrm{MeV}$.
  Tests the tower relationship independently of the framework's $\vEW$ prediction.
\item \textbf{Single input, $\vEW$ route (0.97\%):}
  Using the condensate's predicted $\vEW=247.9\,\mathrm{GeV}$
  gives $\LQCD\approx210\,\mathrm{MeV}$.
  The 0.69\% $\vEW$ residual is an $O(R_0^{-2})$ correction to
  $y_e=e^{-25\pi/6}$ at $R_0=3$ (Paper~IV~\cite{Paper4}, the corrected
  $\vEW$ corollary). %\ref{P4:cor:vew_corrected}
  the same mechanism as the $0.22\%$ residual in $m_e$ before
  the $-\alpha/\pi$ scheme correction.
  Thin-torus theoretical residual: $0.04\%$;
  solving the thick-torus $O(R_0^{-2})$ profile closes the
  single-input prediction to this level (Paper~XIV~\cite{Paper14}). %\ref{P14:cor:yukawa_3D}
\end{itemize}
The exponent $5Q/4=25/2=h(A_4)\cdot h(E_8)/h(E_6)=5\times30/12$ links three
Coxeter numbers (Remark~\ref{P5:rem:e6_coxeter}).
The $\beta$-function coefficients $b_0=7=m_2(E_8)$ and $b_1=26=2m_4(E_8)$
are $E_8$ Coxeter exponents (Remark~\ref{P5:rem:beta_coxeter}).
Nineteenth, the CP-violating phase is derived via two competing WZW mechanisms
(Theorem~\ref{P5:thm:delta_cp}), in exact analogy with the two-mechanism
structure of $\alpha^{-1}$ in Paper~IV~\cite{Paper4}:
the dominant non-Abelian WZW mechanism gives $\delta_{\rm CP}^{\rm(LO)} = 17\pi/40 = 76.5^\circ$
(overshoots), and the subleading E$_8$ McKay twisted-sector correction with
the exact weight $Q/(k\,h(E_8)) = 1/9$ (the condensate phase quantum) gives
\begin{equation*}
  \delta_{\rm CP} = \arg\!\left(e^{i\cdot17\pi/40}
    + \frac{Q}{k\,h(E_8)}\sum_{m\in\mathcal{E}_8} d_m\,e^{2\pi i m/90}\right)
  = 65.0^\circ \approx \frac{13\pi}{36},
\end{equation*}
within $0.12\sigma$ of the PDG value $65.4^\circ\pm3.3^\circ$.
The ratio of the two corrections equals the golden ratio~$\phi$.
The weight $Q/(k\,h(E_8)) = 1/9$ is the same condensate phase quantum that
determines the quark-lepton mass ratios (Theorem~\ref{P5:thm:quark_lepton_ratios}),
establishing a direct link between $\delta_{\rm CP}$ and the quark mass hierarchy.

No further open problems remain at the leading-order level of this framework.
The $0.003^\circ$ gap between $65.0^\circ$ and $65.4^\circ$ is
a higher-order correction in the E$_8$ phase sum and is numerically negligible
relative to the experimental precision $\pm3.3^\circ$.

Twentieth, the muon-to-tau mass ratio
(Remark~\ref{P5:rem:coxeter_connection}):
the same phase gap $\Delta T_{23}=T_2-T_3=1/20=1/h(D_6)$
that gives $|V_{cb}|=e^{-2\pi Q\Delta T_{23}}=e^{-\pi}$
also gives
\begin{equation*}
  \frac{m_\mu}{m_\tau}
  \;=\; e^{-2\pi(Q-1)\Delta T_{23}}
  \;=\; e^{-9\pi/10}
  \;\approx\; 0.05917
  \quad(\text{measured: }0.05946,\;\;0.50\%),
\end{equation*}
when the $Q$-th sector (occupied by the tau as reference state)
is excluded, leaving $Q-1=9$ available sectors.
The sector-exclusion mechanism is not yet proved from first
principles and is recorded as an open problem
(item~\ref{P5:item:lepton_ratio}).

Twenty-first, the exact T-phase algebraic identity
(Remark~\ref{P5:rem:coxeter_connection}, eq.~\eqref{P5:eq:Tphase_identity}):
\begin{equation*}
  T_1 - T_2 - T_3
  \;=\; \frac{T_1-T_2}{Q}
  \;=\; \frac{1}{120}
  \quad[\text{exact}],
\end{equation*}
equivalently $T_3 = (T_1-T_2)(1-1/Q)$,
linking the Casimir self-duality condition $T_3=c/24$,
the $E_6$ Coxeter number $h(E_6)=12$, and the Hopf charge
$Q=10$ in a single arithmetic relation with no free parameters.

%─── References ───────────────────────────────────────────────────────────────
\begin{thebibliography}{99}

\bibitem{Paper1}
F.~Manfredi,
\textit{The Density-Feedback Faddeev--Niemi Hopfion:
  Exact Algebraic Predictions for Standard-Model Parameters},
Zenodo (2026).
\href{https://doi.org/10.5281/zenodo.19342027}{doi:10.5281/zenodo.19342027}

\bibitem{Paper2}
F.~Manfredi,
\textit{BPS Structure and Fixed-Point Theorem for the
  Density-Feedback Faddeev--Niemi Hopfion},
Zenodo (2026).
\href{https://doi.org/10.5281/zenodo.19363491}{doi:10.5281/zenodo.19363491}

\bibitem{Paper3}
F.~Manfredi,
\textit{Two Constants from One Knot: The Fine Structure Constant and
  the Linking Scale as Exact Predictions of the Icosahedral WZW Condensate},
Zenodo (2026).
\href{https://doi.org/10.5281/zenodo.19478629}{doi:10.5281/zenodo.19478629}

\bibitem{Paper4}
F.~Manfredi,
\textit{The Hopf Spoke, WZW Fermion Mass Renormalization,
  and Three- and Four-Term Formula for the Fine Structure Constant},
Preprint (2026).
\href{https://doi.org/10.5281/zenodo.19504729}{doi:10.5281/zenodo.19504729}

\bibitem{Paper13}
F.~Manfredi,
\textit{Atomic Quantum Mechanics from the Hopfion Condensate},
Zenodo (2026).
\href{https://doi.org/10.5281/zenodo.20471482}{doi:10.5281/zenodo.20471482}

\bibitem{Paper14}
F.~Manfredi,
\textit{The Thick-Torus Profile Correction to the Density-Feedback Hopfion},
Zenodo (2026).
\href{https://doi.org/10.5281/zenodo.20479790}{doi:10.5281/zenodo.20479790}

\bibitem{Verlinde1988}
E.~Verlinde,
\textit{Fusion rules and modular transformations in 2D conformal field theory},
Nucl.\ Phys.\ B \textbf{300}, 360--376 (1988).
\href{https://doi.org/10.1016/0550-3213(88)90603-7}{doi:10.1016/0550-3213(88)90603-7}

\bibitem{Witten1984}
E.~Witten,
\textit{Non-Abelian Bosonization in Two Dimensions},
Commun.\ Math.\ Phys.\ \textbf{92}, 455--472 (1984).
\href{https://doi.org/10.1007/BF01215276}{doi:10.1007/BF01215276}

\bibitem{SemiDirac}
S.~Banerjee, R.~R.~P.~Singh, V.~Pardo, and W.~E.~Pickett,
\textit{Tight-binding modeling and low-energy behavior of the semi-Dirac point},
Phys.\ Rev.\ Lett.\ \textbf{103}, 016402 (2009).
\href{https://doi.org/10.1103/PhysRevLett.103.016402}{doi:10.1103/PhysRevLett.103.016402}

\bibitem{PDG2024}
R.~L. Workman et al.\ (Particle Data Group),
\textit{Review of Particle Physics},
Prog.\ Theor.\ Exp.\ Phys.\ \textbf{2022}, 083C01 (2022);
updated 2024.
\href{https://doi.org/10.1093/ptep/ptac097}{doi:10.1093/ptep/ptac097}

\bibitem{Fritzsch1978}
H.~Fritzsch,
\textit{Quark masses and flavor mixing},
Nucl.\ Phys.\ B \textbf{155}, 189--207 (1979).
\href{https://doi.org/10.1016/0550-3213(79)90362-6}{doi:10.1016/0550-3213(79)90362-6}

\bibitem{GJ1979}
H.~Georgi and C.~Jarlskog,
\textit{A new lepton-quark mass relation in a unified theory},
Phys.\ Lett.\ B \textbf{86}, 297--300 (1979).
\href{https://doi.org/10.1016/0370-2693(79)90842-6}{doi:10.1016/0370-2693(79)90842-6}

\bibitem{GST1968}
R.~Gatto, G.~Sartori and M.~Tonin,
\textit{Weak selfmasses, Cabibbo angle, and broken SU(2) $\times$ SU(2)},
Phys.\ Lett.\ B \textbf{28}, 128--130 (1968).
\href{https://doi.org/10.1016/0370-2693(68)90150-0}{doi:10.1016/0370-2693(68)90150-0}

\bibitem{DVVV1987}
R.~Dijkgraaf, C.~Vafa, E.~Verlinde and H.~Verlinde,
\textit{The operator algebra of orbifold models},
Commun.\ Math.\ Phys.\ \textbf{123}, 485--526 (1989).
\href{https://doi.org/10.1007/BF01238812}{doi:10.1007/BF01238812}

\bibitem{CIZ1987}
A.~Cappelli, C.~Itzykson and J.-B.~Zuber,
\textit{The A-D-E classification of minimal and $A_1^{(1)}$ conformal
  invariant theories},
Commun.\ Math.\ Phys.\ \textbf{113}, 1--26 (1987).
\href{https://doi.org/10.1007/BF01221394}{doi:10.1007/BF01221394}

\bibitem{McKay1980}
J.~McKay,
\textit{Graphs, singularities, and finite groups},
Proc.\ Symp.\ Pure Math.\ \textbf{37}, 183--186 (1980).
\href{https://doi.org/10.1090/pspum/037}{doi:10.1090/pspum/037}

\bibitem{Koide1983}
Y.~Koide,
\textit{A fermion-boson composite model of quarks and leptons},
Phys.\ Lett.\ B \textbf{120}, 161--165 (1983).
\href{https://doi.org/10.1016/0370-2693(83)90644-5}{doi:10.1016/0370-2693(83)90644-5}

\bibitem{Cabibbo1963}
N.~Cabibbo,
\textit{Unitary symmetry and leptonic decays},
Phys.\ Rev.\ Lett.\ \textbf{10}, 531--533 (1963).
\href{https://doi.org/10.1103/PhysRevLett.10.531}{doi:10.1103/PhysRevLett.10.531}

\bibitem{KM1973}
M.~Kobayashi and T.~Maskawa,
\textit{CP-violation in the renormalizable theory of weak interaction},
Prog.\ Theor.\ Phys.\ \textbf{49}, 652--657 (1973).
\href{https://doi.org/10.1143/PTP.49.652}{doi:10.1143/PTP.49.652}

\bibitem{GrossWilczek1973}
D.~J.~Gross and F.~Wilczek,
\textit{Ultraviolet behavior of non-abelian gauge theories},
Phys.\ Rev.\ Lett.\ \textbf{30}, 1343--1346 (1973).
\href{https://doi.org/10.1103/PhysRevLett.30.1343}{doi:10.1103/PhysRevLett.30.1343}

\bibitem{Politzer1973}
H.~D.~Politzer,
\textit{Reliable perturbative results for strong interactions?},
Phys.\ Rev.\ Lett.\ \textbf{30}, 1346--1349 (1973).
\href{https://doi.org/10.1103/PhysRevLett.30.1346}{doi:10.1103/PhysRevLett.30.1346}

\bibitem{McKay1980}
J.~McKay,
\textit{Graphs, singularities, and finite groups},
Proc.\ Symp.\ Pure Math.\ \textbf{37}, 183--186 (1980).
\href{https://doi.org/10.1090/pspum/037}{doi:10.1090/pspum/037}

\bibitem{Koide1983}
Y.~Koide,
\textit{A fermion-boson composite model of quarks and leptons},
Phys.\ Lett.\ B \textbf{120}, 161--165 (1983).
\href{https://doi.org/10.1016/0370-2693(83)90644-5}{doi:10.1016/0370-2693(83)90644-5}

\bibitem{Cabibbo1963}
N.~Cabibbo,
\textit{Unitary symmetry and leptonic decays},
Phys.\ Rev.\ Lett.\ \textbf{10}, 531--533 (1963).
\href{https://doi.org/10.1103/PhysRevLett.10.531}{doi:10.1103/PhysRevLett.10.531}

\bibitem{KM1973}
M.~Kobayashi and T.~Maskawa,
\textit{CP-violation in the renormalizable theory of weak interaction},
Prog.\ Theor.\ Phys.\ \textbf{49}, 652--657 (1973).
\href{https://doi.org/10.1143/PTP.49.652}{doi:10.1143/PTP.49.652}

\bibitem{GrossWilczek1973}
D.~J.~Gross and F.~Wilczek,
\textit{Ultraviolet behavior of non-abelian gauge theories},
Phys.\ Rev.\ Lett.\ \textbf{30}, 1343--1346 (1973).
\href{https://doi.org/10.1103/PhysRevLett.30.1343}{doi:10.1103/PhysRevLett.30.1343}

\bibitem{Politzer1973}
H.~D.~Politzer,
\textit{Reliable perturbative results for strong interactions?},
Phys.\ Rev.\ Lett.\ \textbf{30}, 1346--1349 (1973).
\href{https://doi.org/10.1103/PhysRevLett.30.1346}{doi:10.1103/PhysRevLett.30.1346}

\end{thebibliography}

\end{document}